% Oscillateurs mécaniques — Physique Tle C — Cours signé OnBuch+
% Programme officiel MINESEC. Compiler avec : tectonic P09-oscillateurs-mecaniques.tex
\newcommand{\DOCMATIERE}{Physique}
\newcommand{\DOCNIVEAU}{Tle C}
\newcommand{\DOCMODULE}{Module 2 : Mouvements et interactions — évolutions temporelles des systèmes mécaniques}
\newcommand{\DOCLECON}{P09}
\newcommand{\DOCTITRE}{Oscillateurs mécaniques}
% OnBuch+ — préambule « cours signés » ENRICHI (compilé avec tectonic / XeLaTeX)
% Système visuel « Pop » d'OnBuch+ : fond crème (jamais blanc), bordures encre
% épaisses, ombres « dures » décalées, titres Archivo Black en capitales.
% Version MATHÉMATIQUES : graphes (pgfplots/tikz), boîtes pédagogiques
% (définition, propriété, méthode, attention, le savais-tu, exercice résolu,
% exercice, TP, bilan), page de garde et signature OnBuch+ ancrée en bas.
%
% Macros attendues dans main.tex AVANT \input{preamble} :
%   \DOCMATIERE  \DOCNIVEAU  \DOCMODULE  \DOCLECON (numéro)  \DOCTITRE
\documentclass[11pt]{article}

\usepackage{fontspec}
\usepackage[french]{babel}
\usepackage[a4paper,margin=2cm,top=3.4cm,bottom=2.8cm,headheight=1.4cm,footskip=1.3cm]{geometry}
\usepackage{amsmath,amssymb,amsfonts}
\usepackage{mathtools} % \xmapsto, \coloneqq... souvent utilisés par les modèles
\usepackage{cancel} % simplifications barrées (bilans de forces)
% \abs{z} (module, valeur absolue) et \norm{u} (norme), utilisés par les modèles
\providecommand{\abs}{}\let\abs\relax\DeclarePairedDelimiter{\abs}{\lvert}{\rvert}
\providecommand{\norm}{}\let\norm\relax\DeclarePairedDelimiter{\norm}{\lVert}{\rVert}
\usepackage[table]{xcolor} % [table] : fournit aussi \rowcolors (alternance de lignes)
\usepackage{pagecolor}
\usepackage{tikz}
\usetikzlibrary{arrows.meta,positioning,calc,shapes.geometric,shapes.misc,shapes.arrows,decorations.pathmorphing,decorations.markings,patterns,shadows,mindmap,trees,fit,backgrounds,matrix,intersections}
\usepackage{pgfplots}
\usepackage[version=4]{mhchem} % équations nucléaires \ce{^{235}_{92}U}
\usepackage[european,siunitx]{circuitikz} % schémas électriques
\pgfplotsset{compat=1.18}
\usepgfplotslibrary{fillbetween} % aires entre courbes (calcul intégral)
\usepackage{enumitem}
\usepackage{fancyhdr}
% [most] charge toutes les bibliothèques tcolorbox (listings, minted, raster,
% documentation...) et gonfle inutilement la mémoire TeX sur un document long
% (~300 boîtes) : on ne charge que ce qui est réellement utilisé.
\usepackage[skins,breakable]{tcolorbox}
\usepackage{graphicx}
\usepackage{colortbl}
\usepackage{booktabs}
\usepackage{tabularx}
\usepackage{diagbox} % case d'en-tête diagonale des tableaux à double entrée
% Colonnes extensibles centrée / gauche / droite (C, L, R), souvent utilisées par les modèles
\newcolumntype{C}{>{\centering\arraybackslash}X}
\newcolumntype{L}{>{\raggedright\arraybackslash}X}
\newcolumntype{R}{>{\raggedleft\arraybackslash}X}
\usepackage{array}
\usepackage{multicol}
\usepackage{siunitx}
% parse-numbers=false : accepte aussi \SI{2,5 \times 10^{-3}}{...}
\sisetup{locale=FR,output-decimal-marker={,},group-minimum-digits=4,parse-numbers=false}
\DeclareSIUnit\division{div} % réglages d'oscilloscope (ms/div, V/div)
\usepackage{qrcode}
% \degree : commande fréquemment utilisée par les modèles pour un angle ou une
% température en dehors de \SI{}{} (non fournie par les paquets chargés).
\providecommand{\degree}{^{\circ}}
% \qed (fin de démonstration), utilisé par les modèles sans amsthm
\providecommand{\qed}{\hfill$\square$}
% \barre : double barre des valeurs interdites dans un tableau de variations (tkz-tab)
\providecommand{\barre}{\ensuremath{\|}}
% environnement proof (amsthm non chargé) utilisé par les modèles
\ifdefined\proof\else\newenvironment{proof}[1][Démonstration]{\par\noindent\textit{#1.}\ }{\qed\par}\fi
\usepackage{lastpage}
\usepackage{hyperref}
\hypersetup{hidelinks}

% ============================================================
% Palette « Pop » OnBuch+ — mêmes hex que la classe P (plus_ui.dart)
% ============================================================
\definecolor{poporange}{HTML}{F59321}
\definecolor{popdark}{HTML}{E07A0C}
\definecolor{popcream}{HTML}{FFF6E8}
\definecolor{popink}{HTML}{1C1714}
\definecolor{poppurple}{HTML}{7A5AE0}
\definecolor{popgreen}{HTML}{2E9E6A}
\definecolor{poppink}{HTML}{E0457B}
\definecolor{popblue}{HTML}{2F7DE1}
\definecolor{popgold}{HTML}{C98A12}
\definecolor{popmuted}{HTML}{8A7F76}
\definecolor{popline}{HTML}{EADFD2}
\definecolor{popgray}{HTML}{8A7F76}
\colorlet{popgrey}{popgray}
% Alias tolérants (le modèle nomme parfois les couleurs différemment)
\colorlet{popwhite}{white}
\colorlet{popblack}{popink}
\colorlet{popred}{poppink}
\colorlet{popyellow}{popgold}
% Teintes claires pour fonds de tableaux / zones
\colorlet{popblueL}{popblue!10!white}
\colorlet{poporangeL}{poporange!14!white}
\colorlet{popgreenL}{popgreen!12!white}
\colorlet{poppurpleL}{poppurple!12!white}
\colorlet{poppinkL}{poppink!12!white}
\colorlet{popgoldL}{popgold!14!white}

\pagecolor{popcream}

% ============================================================
% Typographie
% ============================================================
\defaultfontfeatures{Ligatures=TeX}
\setmainfont{PlusJakartaSans-Medium.ttf}[Path=../../../fonts/,BoldFont=PlusJakartaSans-Bold.ttf]
% Maths en sans-serif (Fira Math) avec chiffres et lettres droites de Plus
% Jakarta, pour que les formules se fondent dans le texte.
\usepackage[math-style=ISO,bold-style=ISO]{unicode-math}
\setmathfont{FiraMath-Regular.otf}[Path=../../../fonts/]
\setmathfont{PlusJakartaSans-Medium.ttf}[Path=../../../fonts/,range={up/num,up/{Latin,latin},\mathup}]
\usepackage{icomma} % virgule décimale sans espace en mode math
% Caractères mathématiques Unicode absents des polices de texte (Plus Jakarta,
% Archivo Black) : rendus par leur équivalent mathématique, en texte comme en
% formule — sinon ils s'affichent en carré vide (ex. « Arithmétique dans ℤ »).
\usepackage{newunicodechar}
\newunicodechar{ℝ}{\ensuremath{\mathbb{R}}}
\newunicodechar{ℤ}{\ensuremath{\mathbb{Z}}}
\newunicodechar{ℂ}{\ensuremath{\mathbb{C}}}
\newunicodechar{ℕ}{\ensuremath{\mathbb{N}}}
\newunicodechar{ℚ}{\ensuremath{\mathbb{Q}}}
\newunicodechar{𝔻}{\ensuremath{\mathbb{D}}}
\newunicodechar{α}{\ensuremath{\alpha}}
\newunicodechar{β}{\ensuremath{\beta}}
\newunicodechar{γ}{\ensuremath{\gamma}}
\newunicodechar{δ}{\ensuremath{\delta}}
\newunicodechar{ε}{\ensuremath{\varepsilon}}
\newunicodechar{θ}{\ensuremath{\theta}}
\newunicodechar{λ}{\ensuremath{\lambda}}
\newunicodechar{μ}{\ensuremath{\mu}}
\newunicodechar{σ}{\ensuremath{\sigma}}
\newunicodechar{τ}{\ensuremath{\tau}}
\newunicodechar{φ}{\ensuremath{\varphi}}
\newunicodechar{ψ}{\ensuremath{\psi}}
\newunicodechar{ω}{\ensuremath{\omega}}
\newunicodechar{Σ}{\ensuremath{\Sigma}}
% majuscules produites par \MakeUppercase dans les titres (α-aminés -> Α)
\newunicodechar{Α}{\ensuremath{\alpha}}
\newunicodechar{Β}{\ensuremath{\beta}}
\newunicodechar{Γ}{\ensuremath{\Gamma}}
\newunicodechar{Δ}{\ensuremath{\Delta}}
\newunicodechar{Ε}{\ensuremath{\varepsilon}}
\newunicodechar{Θ}{\ensuremath{\Theta}}
\newunicodechar{Λ}{\ensuremath{\Lambda}}
\newunicodechar{Μ}{\ensuremath{\mu}}
\newunicodechar{Τ}{\ensuremath{\tau}}
\newunicodechar{Φ}{\ensuremath{\Phi}}
\newunicodechar{Ψ}{\ensuremath{\Psi}}
\newunicodechar{Ω}{\ensuremath{\Omega}}
\newunicodechar{↦}{\ensuremath{\mapsto}}
\newunicodechar{∈}{\ensuremath{\in}}
\newunicodechar{∉}{\ensuremath{\notin}}
\newunicodechar{∧}{\ensuremath{\wedge}}
\newunicodechar{⇔}{\ensuremath{\Leftrightarrow}}
\newunicodechar{⟺}{\ensuremath{\Longleftrightarrow}}
\newunicodechar{⇒}{\ensuremath{\Rightarrow}}
\newunicodechar{⟹}{\ensuremath{\Longrightarrow}}
\newunicodechar{∘}{\ensuremath{\circ}}
\newunicodechar{∪}{\ensuremath{\cup}}
\newunicodechar{∩}{\ensuremath{\cap}}
\newunicodechar{≡}{\ensuremath{\equiv}}
\newunicodechar{⊂}{\ensuremath{\subset}}
\newunicodechar{⊥}{\ensuremath{\perp}}
\newunicodechar{∃}{\ensuremath{\exists}}
\newunicodechar{∀}{\ensuremath{\forall}}
\newunicodechar{∛}{\ensuremath{\sqrt[3]{\,}}}
\newunicodechar{∜}{\ensuremath{\sqrt[4]{\,}}}
\newunicodechar{⃗}{\ensuremath{{}^{\to}}}
\newunicodechar{ⁿ}{\ifmmode^{n}\else\textsuperscript{\ensuremath{n}}\fi}
\newunicodechar{ˣ}{\ifmmode^{x}\else\textsuperscript{\ensuremath{x}}\fi}
\newunicodechar{ᵇ}{\ifmmode^{b}\else\textsuperscript{\ensuremath{b}}\fi}
\newunicodechar{ʳ}{\ifmmode^{r}\else\textsuperscript{\ensuremath{r}}\fi}
\newunicodechar{ᵘ}{\ifmmode^{u}\else\textsuperscript{\ensuremath{u}}\fi}
\newunicodechar{ʸ}{\ifmmode^{y}\else\textsuperscript{\ensuremath{y}}\fi}
\newunicodechar{ᵃ}{\ifmmode^{a}\else\textsuperscript{\ensuremath{a}}\fi}
\newunicodechar{ᵉ}{\ifmmode^{e}\else\textsuperscript{\ensuremath{e}}\fi}
\newunicodechar{ᵏ}{\ifmmode^{k}\else\textsuperscript{\ensuremath{k}}\fi}
\newunicodechar{ᵖ}{\ifmmode^{p}\else\textsuperscript{\ensuremath{p}}\fi}
\newunicodechar{ᴺ}{\ifmmode^{N}\else\textsuperscript{\ensuremath{N}}\fi}
\newunicodechar{ᶜ}{\ifmmode^{c}\else\textsuperscript{\ensuremath{c}}\fi}
\newunicodechar{ʰ}{\ifmmode^{h}\else\textsuperscript{\ensuremath{h}}\fi}
\newunicodechar{ᵠ}{\ifmmode^{q}\else\textsuperscript{\ensuremath{q}}\fi}
\newunicodechar{ₙ}{\ifmmode_{n}\else\textsubscript{\ensuremath{n}}\fi}
\newunicodechar{ₐ}{\ifmmode_{a}\else\textsubscript{\ensuremath{a}}\fi}
\newunicodechar{ᵢ}{\ifmmode_{i}\else\textsubscript{\ensuremath{i}}\fi}
\newunicodechar{ᵣ}{\ifmmode_{r}\else\textsubscript{\ensuremath{r}}\fi}
\newunicodechar{ₖ}{\ifmmode_{k}\else\textsubscript{\ensuremath{k}}\fi}
\newunicodechar{ᵦ}{\ifmmode_{\beta}\else\textsubscript{\ensuremath{\beta}}\fi}
\newunicodechar{ₓ}{\ifmmode_{x}\else\textsubscript{\ensuremath{x}}\fi}
\newfontfamily\popdisplay{ArchivoBlack-Regular.ttf}[Path=../../../fonts/]
\newfontfamily\poplogo{SpaceGrotesk-Bold.ttf}[Path=../../../fonts/]
\newfontfamily\popsemi{PlusJakartaSans-SemiBold.ttf}[Path=../../../fonts/]
\newfontfamily\popxbold{PlusJakartaSans-ExtraBold.ttf}[Path=../../../fonts/]
\newfontfamily\popxboldls{PlusJakartaSans-ExtraBold.ttf}[Path=../../../fonts/,LetterSpace=3.0]

\setlist[enumerate]{leftmargin=1.7em,itemsep=5pt,topsep=4pt}
\setlist[itemize]{leftmargin=1.7em,itemsep=4pt,topsep=4pt,label={\color{poporange}\textbullet}}
\setlength{\parindent}{0pt}
\setlength{\parskip}{5pt}
\renewcommand{\arraystretch}{1.25}

% pgfplots : style Pop par défaut
\pgfplotsset{
  every axis/.append style={axis line style={popink,line width=1pt},tick style={popink},
    grid style={popline,line width=0.5pt},label style={font=\small},tick label style={font=\footnotesize}},
  popcurve/.style={poporange,line width=1.8pt,no markers,smooth},
}
% Même style disponible dans un \draw TikZ simple (hors pgfplots)
\tikzset{popcurve/.style={poporange,line width=1.8pt}}
\tikzset{popgrid/.style={popline,very thin}}

% ============================================================
% Pill / squiggle
% ============================================================
\newcommand{\poppill}[3]{%
  \tikz[baseline=-0.55ex]\node[draw=popink,line width=1.2pt,rounded corners=2.6pt,%
    fill=#2,inner xsep=6.5pt,inner ysep=3pt]{%
    \popxboldls\fontsize{7}{7}\selectfont\color{#3}\MakeUppercase{#1}};%
}
\newcommand{\popsquiggle}[1]{%
  \tikz[baseline=-0.4ex]{
    \draw[#1,line width=2.6pt,line cap=round]
      plot [smooth] coordinates {(0,0.09) (0.34,0.01) (0.68,0.21) (1.02,0.01) (1.36,0.10)};
  }%
}
% Mot-clé surligné (marqueur orange sous le texte)
\newcommand{\cle}[1]{{\popxbold\color{popdark}#1}}

% ============================================================
% En-tête / pied
% ============================================================
\pagestyle{fancy}
\fancyhf{}
\renewcommand{\headrulewidth}{0pt}
\renewcommand{\footrulewidth}{0pt}
\lhead{\raisebox{-0.15ex}{\poplogo\fontsize{13}{13}\selectfont\color{popink}OnBuch\color{poporange}+}}
\rhead{\poppill{\DOCMATIERE\ \textbullet\ \DOCNIVEAU\ \textbullet\ Leçon \DOCLECON}{white}{popink}}
\lfoot{\popxbold\fontsize{7.6}{7.6}\selectfont\color{popmuted}\MakeUppercase{Cours signé OnBuch+}\ \textbullet\ {\popsemi\DOCTITRE}}
\rfoot{\tikz[baseline=-0.6ex]\node[rounded corners=8.5pt,fill=popink,minimum height=17pt,minimum width=17pt,inner xsep=5pt,inner sep=0]{\popxbold\fontsize{8}{8}\selectfont\color{popcream}\thepage\,/\,\pageref*{LastPage}};}

% ============================================================
% Titres
% ============================================================
\newcommand{\chaptitre}[2]{%
  \begin{tcolorbox}[enhanced,colback=poporange,colframe=popink,boxrule=2.6pt,arc=11pt,
      drop shadow=popink,left=15pt,right=15pt,top=12pt,bottom=11pt,before skip=3pt,after skip=18pt]
    \poppill{#2}{popink}{popcream}\par\vspace{9pt}
    {\raggedright\hyphenpenalty=10000\exhyphenpenalty=10000\popdisplay\fontsize{22}{24}\selectfont\color{popcream}\MakeUppercase{#1}\par}
  \end{tcolorbox}
}
% Section numérotée : pastille orange avec le numéro + titre Archivo
\newcounter{popsec}
\newcommand{\coursec}[1]{%
  \stepcounter{popsec}\setcounter{popsub}{0}%
  \par\vspace{16pt}\noindent
  \begin{minipage}{\linewidth}
  \tikz[baseline=-0.7ex]\node[draw=popink,line width=1.6pt,fill=poporange,rounded corners=4pt,minimum size=22pt,inner sep=0]{\popdisplay\fontsize{12}{12}\selectfont\color{popcream}\thepopsec};%
  \hspace{8pt}{\popdisplay\fontsize{15}{17}\selectfont\color{popink}\MakeUppercase{#1}}\par
  \vspace{4pt}\noindent{\color{poporange}\rule{36pt}{3.6pt}}
  \end{minipage}\par\nopagebreak\vspace{8pt}
}
\newcounter{popsub}
\newcommand{\courssub}[1]{%
  \stepcounter{popsub}%
  \par\vspace{9pt}\noindent{\popxbold\fontsize{11.5}{13}\selectfont\color{popink}{\color{poporange}\thepopsec.\thepopsub}\hspace{6pt}#1}\par\nopagebreak\vspace{4pt}
}

% ============================================================
% Boîtes pédagogiques — même recette : fond blanc, bordure épaisse colorée,
% ombre dure encre, badge plein. Titre optionnel [#1].
% ============================================================
\tcbset{popbase/.style={breakable,enhanced,colback=white,boxrule=2.3pt,arc=9pt,
  shadow={3pt}{-3pt}{0pt}{popink},left=14pt,right=14pt,top=11pt,bottom=11pt,before skip=11pt,after skip=15pt,
  fontupper=\popsemi\fontsize{10.6}{14.5}\selectfont\color{popink},
  fontlower=\popsemi\fontsize{10.6}{14.5}\selectfont\color{popink}}}
% Coche dessinée (le glyphe ✓ n'existe pas dans les polices du document)
\newcommand{\popcheck}[1]{\tikz[baseline=-0.5ex]\draw[#1,line width=1.6pt,line cap=round,line join=round](-0.13,0.0)--(-0.04,-0.09)--(0.13,0.1);}
\providecommand{\coche}{\popcheck{popgreen}}
\newcommand{\popboxhead}[3]{\poppill{#1}{#2}{white}\ifx\relax#3\relax\else\hspace{6pt}{\popxbold\fontsize{10.6}{12}\selectfont\color{popink}#3}\fi\par\vspace{7pt}}

\newtcolorbox{aretenir}[1][]{popbase,colframe=popgreen,before upper={\popboxhead{À retenir}{popgreen}{#1}}}
\newtcolorbox{exemplebox}[1][]{popbase,colframe=poporange,before upper={\popboxhead{Exemple}{poporange}{#1}}}
\newtcolorbox{definition}[1][]{popbase,colframe=popblue,before upper={\popboxhead{Définition}{popblue}{#1}}}
\newtcolorbox{propriete}[1][]{popbase,colframe=poppurple,before upper={\popboxhead{Propriété}{poppurple}{#1}}}
\newtcolorbox{methode}[1][]{popbase,colframe=popgold,before upper={\popboxhead{Méthode}{popgold}{#1}}}
\newtcolorbox{attention}[1][]{popbase,colframe=poppink,before upper={\popboxhead{Attention}{poppink}{#1}}}
\newtcolorbox{experience}[1][]{popbase,colframe=popblue,colback=popblueL,before upper={\popboxhead{Expérience}{popblue}{#1}}}
% « Le savais-tu ? » : carte violette pleine (contexte, culture, Cameroun)
\newtcolorbox{savaistu}[1][]{popbase,colback=poppurple,colframe=popink,
  fontupper=\popsemi\fontsize{10.4}{14.2}\selectfont\color{white},
  before upper={\poppill{Le savais-tu\,?}{white}{poppurple}\ifx\relax#1\relax\else\hspace{6pt}{\popxbold\color{white}#1}\fi\par\vspace{7pt}}}
% Exercice résolu : énoncé en haut, \tcblower puis la solution sur fond crème
\newcounter{popexo}\newcounter{popexr}
\newtcolorbox{exoresolu}[1][]{popbase,colframe=poporange,colbacklower=popcream,
  segmentation style={popink,line width=1.2pt,dashed},
  before upper={\stepcounter{popexr}\popboxhead{Exercice résolu \thepopexr}{poporange}{#1}},
  before lower={\poppill{Solution}{popink}{popcream}\par\vspace{6pt}}}
% Exercice (non résolu en place) — la correction est dans la partie Corrigés
\newtcolorbox{exercice}[1][]{popbase,colframe=popink,
  before upper={\stepcounter{popexo}\popboxhead{Exercice \thepopexo}{popink}{#1}}}
% Corrigé
\newtcolorbox{corrige}[1][]{popbase,colframe=popgreen,colback=popgreenL,
  before upper={\popboxhead{Corrigé}{popgreen}{#1}}}
% Objectifs / prérequis / bilan
\newtcolorbox{objectifs}{popbase,colframe=popink,colback=poporangeL,
  before upper={\popboxhead{Objectifs de la leçon}{popink}{}}}
\newtcolorbox{prerequis}{popbase,colframe=popink,colback=popblueL,
  before upper={\popboxhead{Prérequis}{popblue}{}}}
\newtcolorbox{bilan}{popbase,colframe=popink,colback=white,boxrule=2.8pt,
  before upper={\popboxhead{Fiche bilan}{poporange}{L'essentiel à connaître pour l'examen}}}
% Figure encadrée avec légende
\newtcolorbox{popfigure}[1][]{enhanced,breakable=false,colback=white,colframe=popline,boxrule=1.4pt,arc=8pt,
  left=10pt,right=10pt,top=10pt,bottom=8pt,before skip=10pt,after skip=12pt,halign=center,
  lower separated=false,halign lower=center,
  fontlower=\popsemi\fontsize{9}{11.5}\selectfont\color{popmuted},
  #1}
\newcommand{\legende}[1]{\par\vspace{6pt}{\popsemi\fontsize{9}{11.5}\selectfont\color{popmuted}#1}}

% ============================================================
% Signature OnBuch+
% ============================================================
\newcommand{\poplogoinline}[1]{{\poplogo\fontsize{#1}{#1}\selectfont\color{popink}OnBuch\color{poporange}+}}
\newcommand{\signatureonbuch}{%
  \begin{tcolorbox}[enhanced,colback=popink,colframe=popink,boxrule=2.2pt,arc=10pt,
      drop shadow=poporange,left=14pt,right=14pt,top=10pt,bottom=10pt,before skip=0pt,after skip=0pt]
    \tikz[baseline=-0.7ex]\node[circle,fill=popgreen,draw=popcream,line width=1.3pt,minimum size=22pt,inner sep=0]{\popcheck{white}};%
    \hspace{9pt}\begin{minipage}[c]{0.62\linewidth}
      {\popxbold\fontsize{12}{13}\selectfont\color{popcream}Signé\ \textbullet\ {\poplogo OnBuch\color{poporange}+}}\par\vspace{2pt}
      {\raggedright\popsemi\fontsize{8}{10.5}\selectfont\color{popline}\MakeUppercase{Équipe pédagogique OnBuch+}\\\MakeUppercase{Conforme au programme officiel MINESEC}\par}
    \end{minipage}\hfill
    {\poplogo\fontsize{22}{22}\selectfont\color{popcream}OnBuch\color{poporange}+}
  \end{tcolorbox}%
}

% ============================================================
% Page de garde
% #1 titre · #2 matière · #3 niveau · #4 module · #5 n° leçon · #6 durée ·
% #7 description / objectifs (texte ou itemize)
% ============================================================
\newcommand{\pagegarde}[7]{%
  \thispagestyle{empty}
  \newgeometry{a4paper,margin=1.7cm,top=1.6cm,bottom=1.4cm}%
  \begin{tikzpicture}[remember picture,overlay]
    \fill[poporange!18] ([xshift=-3.2cm,yshift=-3.6cm]current page.north east) circle (4.6cm);
    \fill[poppurple!14] ([xshift=2.4cm,yshift=7.6cm]current page.south west) circle (3.8cm);
  \end{tikzpicture}%
  {\poplogo\fontsize{30}{30}\selectfont\color{popink}OnBuch\color{poporange}+}\hfill
  \raisebox{0.6ex}{\poppill{Cours signé \textbullet\ Premium}{popink}{popcream}}\par
  \vspace{1pt}\popsquiggle{poppurple}\par
  \vspace{8pt}
  \begin{tcolorbox}[enhanced,colback=poporange,colframe=popink,boxrule=2.8pt,arc=13pt,
      drop shadow=popink,left=18pt,right=18pt,top=12pt,bottom=13pt,after skip=10pt]
    \poppill{#2 \textbullet\ #3}{popink}{popcream}\hspace{5pt}\poppill{#4}{white}{popink}\par\vspace{8pt}
    {\popdisplay\fontsize{14}{14}\selectfont\color{popink}LEÇON #5}\par\vspace{4pt}
    {\raggedright\hyphenpenalty=10000\exhyphenpenalty=10000\popdisplay\fontsize{25}{27}\selectfont\color{popcream}\MakeUppercase{#1}\par}
    \vspace{8pt}
    {\popxbold\fontsize{9.5}{9.5}\selectfont\color{popink}\MakeUppercase{Durée conseillée : #6}}
  \end{tcolorbox}
  \begin{tcolorbox}[enhanced,colback=white,colframe=popink,boxrule=2.2pt,arc=10pt,
      drop shadow=popink,left=16pt,right=16pt,top=10pt,bottom=10pt,after skip=8pt]
    \poppill{Dans cette leçon}{poppurple}{white}\par\vspace{5pt}
    \popsemi\fontsize{9.3}{12.6}\selectfont\color{popink}#7
  \end{tcolorbox}
  \noindent\begin{minipage}[t]{0.487\linewidth}
    \begin{tcolorbox}[enhanced,colback=popgreen,colframe=popink,boxrule=2.2pt,arc=10pt,
        drop shadow=popink,left=11pt,right=11pt,top=10pt,bottom=10pt,height=3.1cm,valign=center]
      \tcbox[colback=white,colframe=popink,boxrule=1.3pt,arc=5pt,boxsep=0pt,nobeforeafter,
        left=3pt,right=3pt,top=3pt,bottom=3pt,valign=center]{\qrcode[height=1.9cm]{https://play.google.com/store/apps/details?id=com.luvvix.onbuch&hl=fr}}%
      \hspace{8pt}\begin{minipage}[c]{0.5\linewidth}
        {\popdisplay\fontsize{11}{12.5}\selectfont\color{white}\MakeUppercase{Télécharge OnBuch}}\par\vspace{4pt}
        {\raggedright\popsemi\fontsize{8.4}{11}\selectfont\color{white}Gratuit \textbullet\ Google Play\\Tuteur IA, annales, concours, résultats\par}
      \end{minipage}
    \end{tcolorbox}
  \end{minipage}\hfill
  \begin{minipage}[t]{0.487\linewidth}
    \begin{tcolorbox}[enhanced,colback=poppurple,colframe=popink,boxrule=2.2pt,arc=10pt,
        drop shadow=popink,left=11pt,right=11pt,top=10pt,bottom=10pt,height=3.1cm,valign=center]
      \tikz[baseline=-0.7ex]\node[circle,fill=white,draw=popink,line width=1.3pt,minimum size=1.6cm,inner sep=0]{\popdisplay\fontsize{24}{24}\selectfont\color{poppurple}+};%
      \hspace{8pt}\begin{minipage}[c]{0.6\linewidth}
        {\popdisplay\fontsize{11}{12.5}\selectfont\color{white}\MakeUppercase{Passe à OnBuch+}}\par\vspace{4pt}
        {\raggedright\popsemi\fontsize{8.4}{11}\selectfont\color{white}Tous les cours signés, Tuteur IA illimité, fiches PDF — onglet OnBuch+ de l'app\par}
      \end{minipage}
    \end{tcolorbox}
  \end{minipage}
  \vspace{8pt}
  \popcontenu
  \vfill
  \signatureonbuch
  \restoregeometry
  \newpage
}

% Rangée « Ce cours contient » (page de garde)
\newcommand{\poptile}[3]{% #1 couleur #2 gros texte #3 légende
  \begin{tcolorbox}[enhanced,colback=#1,colframe=popink,boxrule=2pt,arc=9pt,shadow={2.5pt}{-2.5pt}{0pt}{popink},
      left=6pt,right=6pt,top=9pt,bottom=9pt,halign=center,valign=center,height=1.75cm,width=\linewidth,nobeforeafter]
    {\popdisplay\fontsize{12.5}{13}\selectfont\color{white}\MakeUppercase{#2}}\par\vspace{3pt}
    {\centering\popsemi\fontsize{7.8}{9.5}\selectfont\color{white}#3\par}
  \end{tcolorbox}}
\newcommand{\popcontenu}{%
  \par\noindent{\popxboldls\fontsize{7.5}{7.5}\selectfont\color{popmuted}CE COURS SIGNÉ CONTIENT}\par\vspace{7pt}
  \noindent\begin{minipage}[t]{0.235\linewidth}\poptile{popblue}{Cours}{complet, illustré, pas à pas}\end{minipage}\hfill
  \begin{minipage}[t]{0.235\linewidth}\poptile{poporange}{Méthodes}{et exercices résolus}\end{minipage}\hfill
  \begin{minipage}[t]{0.235\linewidth}\poptile{poppink}{Exercices}{type examen + corrigés}\end{minipage}\hfill
  \begin{minipage}[t]{0.235\linewidth}\poptile{popgreen}{Fiche bilan}{l'essentiel à retenir}\end{minipage}\par}

% Sommaire
\newtcolorbox{sommaire}{enhanced,colback=white,colframe=popink,boxrule=2.2pt,arc=10pt,
  drop shadow=popink,left=18pt,right=18pt,top=13pt,bottom=13pt,before skip=6pt,after skip=20pt,
  before upper={\poppill{Sommaire}{poppurple}{white}\par\vspace{9pt}\popsemi\fontsize{10.6}{14.5}\selectfont\color{popink}}}

% Signature de fin de cours — poussée tout en bas de la dernière page
\newcommand{\findecours}{%
  \par\vfill\nopagebreak
  \begin{center}\popsquiggle{poporange}\end{center}\vspace{4pt}
  \signatureonbuch
}

%% FIGFIX-BEGIN v3 — correctifs globaux des figures (docs/onbuch-generation/tools/figfix)
\usepackage{adjustbox}\usepackage{etoolbox}
% toute figure TikZ/pgfplots plus large que la ligne est réduite à la largeur disponible
\BeforeBeginEnvironment{tikzpicture}{\begin{adjustbox}{max width=\linewidth}}
\AfterEndEnvironment{tikzpicture}{\end{adjustbox}}
% légendes pgfplots lisibles par-dessus les courbes
\pgfplotsset{every axis/.append style={legend style={fill=white,fill opacity=0.92,text opacity=1,draw=black!15,inner sep=3pt}}}
% étiquettes d'arêtes (syntaxe edge["..."]) avec fond blanc
\tikzset{every edge quotes/.append style={fill=white,inner sep=1.2pt}}
% bulles de cartes mentales : texte à la ligne dans la bulle, bulle un peu plus grande
\tikzset{every concept/.append style={text width=2.0cm,align=center,minimum size=2.3cm,inner sep=2pt}}
%% FIGFIX-END
\begin{document}

\pagegarde{Oscillateurs mécaniques}{Physique}{Tle C}{Module 2 : Mouvements et interactions — évolutions temporelles des systèmes mécaniques}{P09}{7 h}{Cette leçon étudie les oscillateurs mécaniques : pendule élastique, pendule simple et pendule de torsion. Elle établit les équations différentielles du mouvement, détermine les grandeurs caractéristiques (période, pulsation, amplitude, phase), analyse l'énergie mécanique et introduit les oscillations amorties, forcées et le phénomène de résonance.
\vspace{4pt}
\begin{multicols}{2}\small
\begin{itemize}
\item À la fin de cette leçon, je sais établir l'équation différentielle du mouvement d'un point matériel soumis à une force de rappel élastique (pendule horizontal et vertical).
\item À la fin de cette leçon, je sais reconnaître l'équation x'' + ω0²x = 0 et écrire sa solution générale x = Xm cos(ω0 t + φ).
\item À la fin de cette leçon, je sais déterminer l'amplitude et la phase à l'origine à partir des conditions initiales.
\item À la fin de cette leçon, je sais calculer la période propre d'un pendule élastique, d'un pendule simple et d'un pendule de torsion.
\item À la fin de cette leçon, je sais calculer l'énergie mécanique d'un oscillateur harmonique et montrer sa conservation.
\item À la fin de cette leçon, je sais exploiter un enregistrement d'oscillations (mesure de période, pseudo-période, amortissement).
\end{itemize}\end{multicols}}

\begin{sommaire}
\begin{multicols}{2}
\begin{enumerate}
\item Le pendule élastique horizontal
\item Le pendule élastique vertical
\item Pendule pesant et pendule simple
\item Pendule de torsion et étude énergétique
\item Oscillations amorties
\item Oscillations forcées et résonance
\item Activité d'intégration
\item Méthodes et exercices résolus
\item Exercices
\item Corrigés des exercices
\item Fiche bilan
\end{enumerate}
\end{multicols}
\end{sommaire}

\begin{objectifs}
\begin{itemize}
\item À la fin de cette leçon, je sais établir l'équation différentielle du mouvement d'un point matériel soumis à une force de rappel élastique (pendule horizontal et vertical).
\item À la fin de cette leçon, je sais reconnaître l'équation x'' + ω0²x = 0 et écrire sa solution générale x = Xm cos(ω0 t + φ).
\item À la fin de cette leçon, je sais déterminer l'amplitude et la phase à l'origine à partir des conditions initiales.
\item À la fin de cette leçon, je sais calculer la période propre d'un pendule élastique, d'un pendule simple et d'un pendule de torsion.
\item À la fin de cette leçon, je sais calculer l'énergie mécanique d'un oscillateur harmonique et montrer sa conservation.
\item À la fin de cette leçon, je sais exploiter un enregistrement d'oscillations (mesure de période, pseudo-période, amortissement).
\item À la fin de cette leçon, je sais distinguer qualitativement les régimes pseudo-périodique et apériodique.
\item À la fin de cette leçon, je sais décrire les oscillations forcées et expliquer le phénomène de résonance mécanique.
\end{itemize}
\end{objectifs}

\begin{prerequis}
\begin{itemize}
\item Deuxième loi de Newton et théorème du centre d'inertie (Première)
\item Tension d'un ressort : loi de Hooke T = k|x| (Première)
\item Énergie cinétique, énergie potentielle élastique et de pesanteur, énergie mécanique (Première)
\item Dérivation des fonctions sinusoïdales et notion d'équation différentielle (Mathématiques Tle)
\item Théorème du moment cinétique ou relation fondamentale de la dynamique de rotation (début d'année Tle)
\end{itemize}
\end{prerequis}

\newpage

\coursec{Le pendule élastique horizontal}

À Yaoundé, le mécanicien règle la suspension d'une moto-taxi pour qu'elle absorbe les nids-de-poule sans rebondir indéfiniment ; au marché, la balance à ressort oscille avant de se stabiliser ; le balancier d'une horloge traditionnelle bat le temps avec une régularité étonnante. Tous ces systèmes ont un point commun : ils \cle{oscillent} autour d'une position d'équilibre. Comment prévoir la période de ces oscillations ? Pourquoi certaines s'atténuent-elles et d'autres deviennent-elles dangereusement grandes, comme un pont qui entre en résonance ? Cette leçon donne les outils mathématiques et énergétiques pour répondre à ces questions.

Nous commençons par le plus simple de tous les oscillateurs : le \cle{pendule élastique horizontal}. C'est le système modèle à partir duquel toute la physique des oscillations se construit.

\courssub{Description du dispositif et repérage}

Considérons un solide (S) de masse $m$, assimilé à un point matériel $G$, accroché à l'extrémité libre d'un \cle{ressort} à spires non jointives, de masse négligeable et de \cle{constante de raideur} $k$ (en \SI{}{N.m^{-1}}). L'autre extrémité du ressort est fixée à un support immobile. Le solide glisse \emph{sans frottement} sur un support horizontal (une table à coussin d'air en TP, par exemple).

\begin{popfigure}
\begin{tikzpicture}[scale=1.0]
% Support mural
\fill[pattern=north east lines, pattern color=popmuted] (-0.4,-0.2) rectangle (0,2.2);
\draw[thick,popdark] (0,-0.2) -- (0,2.2);
% Ressort (zigzag)
\draw[thick,popblue] (0,1.0) -- (0.3,1.0);
\draw[thick,popblue] (0.3,1.0) -- (0.45,1.25) -- (0.75,0.75) -- (1.05,1.25) -- (1.35,0.75) -- (1.65,1.25) -- (1.95,0.75) -- (2.25,1.25) -- (2.4,1.0);
\draw[thick,popblue] (2.4,1.0) -- (2.8,1.0);
% Solide
\fill[poporange!80] (2.8,0.55) rectangle (4.2,1.45);
\draw[thick,popdark] (2.8,0.55) rectangle (4.2,1.45);
\node at (3.5,1.0) {\small (S), $m$};
% Table
\draw[thick,popdark] (0.2,0.35) -- (6.4,0.35);
\fill[pattern=horizontal lines, pattern color=popmuted] (0.2,0.05) rectangle (6.4,0.35);
% Centre G
\fill[popink] (3.5,1.0) circle (0.05);
\node[below,popink] at (3.5,0.62) {\small $G$};
% Axe Ox
\draw[->,thick,popink] (1.0,-0.5) -- (6.2,-0.5) node[right] {$x$};
\draw[thick,popink] (2.2,-0.62) -- (2.2,-0.38) node[above left] {\small $O$};
\draw[thick,popink] (3.5,-0.62) -- (3.5,-0.38);
\node[below] at (3.5,-0.62) {\small $G$};
\draw[<->,popgreen,thick] (2.2,-1.15) -- (3.5,-1.15) node[midway,below] {\small $x(t)$};
% Forces
\draw[->,very thick,popgreen] (3.5,1.45) -- (3.5,2.35) node[right] {\small $\vec{R}$};
\draw[->,very thick,popgreen] (3.5,0.55) -- (3.5,-0.35) node[right] {\small $\vec{P}$};
\draw[->,very thick,popred] (2.8,1.0) -- (1.6,1.0) node[above] {\small $\vec{T}$};
\end{tikzpicture}
\legende{Pendule élastique horizontal écarté de sa position d'équilibre : le ressort est étiré, la tension $\vec{T}$ rappelle le solide vers $O$.}
\end{popfigure}

\courssub{Position d'équilibre et élongation}

Lorsque le solide est au repos, le ressort n'est ni étiré ni comprimé (support horizontal : le poids ne joue aucun rôle le long de l'axe du mouvement). Le centre d'inertie $G$ occupe alors une position $O$ appelée \cle{position d'équilibre}.

On choisit un axe $(Ox)$ horizontal, dirigé dans le sens du mouvement possible, d'origine $O$. La position de $G$ à l'instant $t$ est repérée par son abscisse $x(t)$, appelée \cle{élongation} (ou écart à l'équilibre). C'est une grandeur \emph{algébrique} : $x > 0$ si le ressort est étiré, $x < 0$ s'il est comprimé.

\begin{attention}[Allongement ou élongation ?]
Ne confonds pas l'\emph{allongement} du ressort $\Delta \ell = \ell - \ell_0$ (mesuré par rapport à la longueur à vide $\ell_0$) et l'\emph{élongation} $x$ (mesurée depuis la position d'équilibre $O$). Pour le pendule \emph{horizontal}, les deux coincident car l'équilibre correspond au ressort non déformé. Mais pour le pendule \emph{vertical} (section suivante), le ressort est déjà allongé à l'équilibre, et c'est bien $x$ — pas l'allongement total — qui apparaît dans l'équation différentielle. Prends dès maintenant l'habitude de tout mesurer depuis $O$.
\end{attention}

\courssub{Mise en équation : la deuxième loi de Newton}

\begin{propriete}[Équation différentielle du pendule élastique horizontal — démonstration exigible]
Le solide (S) de masse $m$, accroché à un ressort de raideur $k$ et glissant sans frottement sur un support horizontal, effectue des oscillations dont l'élongation $x(t)$ vérifie :
\[
\ddot{x} + \frac{k}{m}\,x = 0
\]
Cette démonstration est exigible au Baccalauréat.
\end{propriete}

\begin{experience}[Démonstration guidée pas à pas]
\textbf{Système étudié :} le solide (S) de masse $m$, dans le référentiel terrestre supposé galiléen.

\textbf{Bilan des forces extérieures :}
\begin{itemize}
\item le poids $\vec{P} = m\vec{g}$, vertical, dirigé vers le bas ;
\item la réaction normale $\vec{R}$ du support, verticale, dirigée vers le haut (pas de frottement) ;
\item la tension du ressort $\vec{T} = -k\,x\,\vec{i}$, où $\vec{i}$ est le vecteur unitaire de l'axe $(Ox)$.
\end{itemize}

\textbf{Deuxième loi de Newton :}
\[
\vec{P} + \vec{R} + \vec{T} = m\,\vec{a}
\]

\textbf{Projection} sur l'axe horizontal $(Ox)$ : le poids et la réaction sont verticaux, leur projection est nulle. Il reste :
\[
-kx = m\,\ddot{x}
\]
soit, en divisant par $m$ et en regroupant :
\[
\boxed{\;\ddot{x} + \frac{k}{m}\,x = 0\;}
\]

\textbf{Vérification dimensionnelle :} $[k/m] = \dfrac{\mathrm{N.m^{-1}}}{\mathrm{kg}} = \dfrac{\mathrm{kg.m.s^{-2}.m^{-1}}}{\mathrm{kg}} = \mathrm{s^{-2}}$, et $[\ddot{x}] = \mathrm{m.s^{-2}}$, donc $[\ddot{x}/x] = \mathrm{s^{-2}}$. L'équation est homogène. $\square$
\end{experience}

Le point essentiel à retenir physiquement : la tension du ressort est une \cle{force de rappel}, toujours dirigée vers la position d'équilibre et proportionnelle à l'écart. C'est cette proportionnalité ($T = -kx$) qui engendre des oscillations sinusoïdales.

\begin{definition}[Oscillateur harmonique]
On appelle \cle{oscillateur harmonique} tout système physique dont l'évolution temporelle est régie par une équation différentielle de la forme
\[
\ddot{x} + \omega_0^2\, x = 0
\]
où $\omega_0$ est une constante positive appelée \cle{pulsation propre} (en \SI{}{rad.s^{-1}}). Pour le pendule élastique horizontal :
\[
\omega_0 = \sqrt{\frac{k}{m}}
\]
\end{definition}

\courssub{Solution de l'équation différentielle}

L'équation $\ddot{x} + \omega_0^2 x = 0$ demande une fonction dont la dérivée seconde est proportionnelle à l'opposé d'elle-même. Les fonctions sinus et cosinus possèdent exactement cette propriété.

\begin{propriete}[Solution générale]
La solution générale de l'équation différentielle $\ddot{x} + \omega_0^2 x = 0$ s'écrit :
\[
x(t) = X_m \cos(\omega_0 t + \varphi)
\]
où :
\begin{itemize}
\item $X_m > 0$ est l'\cle{amplitude} (en m) : élongation maximale ;
\item $\omega_0$ est la pulsation propre (en \SI{}{rad.s^{-1}}) ;
\item $\varphi$ est la \cle{phase à l'origine} (en rad) ; la quantité $\omega_0 t + \varphi$ est la \emph{phase} à l'instant $t$.
\end{itemize}
On admet que, pour des conditions initiales $x(0)$ et $v(0)$ données, cette solution est \emph{unique}.
\end{propriete}

\begin{experience}[Vérification par double dérivation]
Vérifions que $x(t) = X_m\cos(\omega_0 t + \varphi)$ satisfait bien l'équation différentielle.
\begin{align*}
x(t) &= X_m\cos(\omega_0 t + \varphi)\\
\dot{x}(t) &= -X_m\,\omega_0 \sin(\omega_0 t + \varphi)\\
\ddot{x}(t) &= -X_m\,\omega_0^2 \cos(\omega_0 t + \varphi) = -\omega_0^2\, x(t)
\end{align*}
Donc $\ddot{x} + \omega_0^2 x = 0$ : la fonction proposée est bien solution, quelles que soient les constantes $X_m$ et $\varphi$. $\square$
\end{experience}

Le mouvement est donc \cle{sinusoïdal} : on dit que le solide effectue des \emph{oscillations harmoniques non amorties}. La vitesse $v(t) = \dot{x}(t)$ est en quadrature de phase par rapport à l'élongation : la vitesse est maximale quand le solide passe par $O$ ($x=0$), et nulle aux extrémités ($x = \pm X_m$).

\begin{popfigure}
\begin{tikzpicture}
\begin{axis}[
width=\linewidth, height=6cm,
xlabel={$t$ (s)}, ylabel={$x(t)$ (m)},
xmin=0, xmax=1.3, ymin=-0.06, ymax=0.075,
xtick={0,0.1,0.5,0.9,1.3},
xticklabels={$0$, , $T_0$, , $2T_0$},
ytick={-0.05,0,0.05},
yticklabels={$-X_m$,$0$,$X_m$},
grid=both, grid style={popline!40},
axis lines=left, thick]
\addplot[popcurve,domain=0:1.3,samples=300]{0.05*cos(deg(2*pi*2.5*x + 0.6))};
% amplitude
\draw[<->,poporange,thick] (axis cs:1.15,0) -- (axis cs:1.15,0.05) node[midway,right]{\small $X_m$};
% phase marker
\fill[popred] (axis cs:0,0.0405) circle (1.5pt);
\node[popred,above right] at (axis cs:0.02,0.045) {\small $x(0)=X_m\cos\varphi$};
% period arrow
\draw[<->,popgreen,thick] (axis cs:0.038,-0.062) -- (axis cs:0.438,-0.062) node[midway,below]{\small $T_0$};
\end{axis}
\end{tikzpicture}
\legende{Élongation $x(t) = X_m\cos(\omega_0 t + \varphi)$ : oscillations harmoniques de période propre $T_0$, d'amplitude $X_m$ constante.}
\end{popfigure}

\courssub{Période propre et isochronisme}

\begin{propriete}[Période propre du pendule élastique]
La \cle{période propre} des oscillations du pendule élastique est :
\[
T_0 = \frac{2\pi}{\omega_0} = 2\pi\sqrt{\frac{m}{k}}
\qquad\text{et la fréquence propre :}\qquad
f_0 = \frac{1}{T_0} = \frac{1}{2\pi}\sqrt{\frac{k}{m}}
\]
$T_0$ s'exprime en secondes (s), $f_0$ en hertz (Hz).
\end{propriete}

\begin{aretenir}[Isochronisme des oscillations]
La période propre $T_0 = 2\pi\sqrt{m/k}$ ne dépend \emph{que} de la masse $m$ et de la raideur $k$. Elle est \textbf{indépendante de l'amplitude} $X_m$ : c'est l'\cle{isochronisme} des oscillations harmoniques. Qu'on écarte le solide de \SI{2}{cm} ou de \SI{8}{cm}, il battra au même rythme. C'est ce principe qui fait la régularité des horloges.
\end{aretenir}

Analyse dimensionnelle : $[m/k] = \dfrac{\mathrm{kg}}{\mathrm{kg.s^{-2}}} = \mathrm{s^2}$, donc $\sqrt{m/k}$ est bien homogène à un temps. La formule est cohérente : plus la masse est grande, plus le mouvement est lent ; plus le ressort est raide, plus les oscillations sont rapides.

\begin{exemplebox}[Calcul d'une période propre]
Un solide de masse $m = \SI{200}{g} = \SI{0,200}{kg}$ est accroché à un ressort de raideur $k = \SI{50}{N.m^{-1}}$.
\[
T_0 = 2\pi\sqrt{\frac{m}{k}} = 2\pi\sqrt{\frac{0,200}{50}} = 2\pi\sqrt{4,0\times 10^{-3}} \approx \SI{0,40}{s}
\]
La fréquence propre vaut $f_0 = \dfrac{1}{T_0} \approx \SI{2,5}{Hz}$ : le solide effectue environ 2,5 oscillations par seconde.
\end{exemplebox}

\courssub{Exploitation des conditions initiales}

Les constantes $X_m$ et $\varphi$ ne sont pas imposées par la physique du système : elles dépendent de la \emph{manière dont on lance} l'oscillateur, c'est-à-dire des \cle{conditions initiales} $x(0) = x_0$ et $v(0) = v_0$.

\begin{methode}[Déterminer $X_m$ et $\varphi$ à partir des conditions initiales]
\begin{enumerate}
\item Écrire la solution et sa dérivée à $t = 0$ :
\[
x(0) = X_m\cos\varphi = x_0
\qquad
v(0) = -X_m\,\omega_0\sin\varphi = v_0
\]
\item Éliminer $\varphi$ en sommant les carrés (avec $\cos^2\varphi + \sin^2\varphi = 1$) :
\[
X_m = \sqrt{x_0^2 + \frac{v_0^2}{\omega_0^2}}
\]
\item Éliminer $X_m$ par le quotient :
\[
\tan\varphi = -\frac{v_0}{\omega_0\, x_0}
\]
\item Choisir la bonne valeur de $\varphi$ grâce au \emph{signe de $\cos\varphi$} (signe de $x_0$) : la tangente seule donne deux solutions possibles, une seule est correcte.
\end{enumerate}
\end{methode}

\begin{attention}[Pièges fréquents sur la phase]
\begin{itemize}
\item Si le solide est lâché \emph{sans vitesse initiale} depuis $x_0 > 0$, alors $\varphi = 0$ et $X_m = x_0$ : $x(t) = x_0\cos(\omega_0 t)$.
\item $\tan\varphi = a$ admet deux solutions dans $]-\pi\,;\,\pi]$ : vérifie toujours le signe de $\cos\varphi$ avec $x_0$.
\item L'amplitude $X_m$ est \emph{toujours positive} ; une valeur négative indique une erreur de calcul.
\end{itemize}
\end{attention}

\begin{exoresolu}[Oscillations d'un pendule élastique]
\textbf{Énoncé.} Un solide de masse $m = \SI{200}{g}$ est fixé à un ressort horizontal de raideur $k = \SI{50}{N.m^{-1}}$. On écarte le solide de $x_0 = \SI{5,0}{cm}$ de sa position d'équilibre et on le lâche sans vitesse initiale à $t = 0$.
\begin{enumerate}
\item Établir l'équation différentielle du mouvement et calculer la période propre.
\item Déterminer $X_m$ et $\varphi$, puis écrire $x(t)$.
\item Calculer l'élongation et la vitesse à $t = \SI{0,10}{s}$.
\end{enumerate}
\tcblower
\textbf{Solution.}
\begin{enumerate}
\item Le bilan des forces projeté sur $(Ox)$ donne $m\ddot{x} = -kx$, soit $\ddot{x} + \dfrac{k}{m}x = 0$. La pulsation propre est
\[
\omega_0 = \sqrt{\frac{k}{m}} = \sqrt{\frac{50}{0,200}} = \sqrt{250} \approx \SI{15,8}{rad.s^{-1}},
\qquad
T_0 = \frac{2\pi}{\omega_0} \approx \SI{0,40}{s}.
\]
\item Conditions initiales : $x_0 = \SI{5,0e-2}{m}$ et $v_0 = 0$. Alors
\[
X_m = \sqrt{x_0^2 + \frac{v_0^2}{\omega_0^2}} = x_0 = \SI{5,0}{cm},
\qquad
\tan\varphi = 0 \ \text{et}\ \cos\varphi > 0 \Rightarrow \varphi = 0.
\]
D'où $x(t) = 5,0\times 10^{-2}\cos(15,8\,t)$ (en mètres).
\item À $t = \SI{0,10}{s}$ : $\omega_0 t = 15,8 \times 0,10 = \SI{1,58}{rad} \approx \dfrac{\pi}{2}$.
\[
x(0,10) \approx 5,0\times 10^{-2}\cos\!\left(\tfrac{\pi}{2}\right) \approx \SI{0}{m},
\qquad
v(0,10) = -\omega_0 X_m \sin\!\left(\tfrac{\pi}{2}\right) \approx -15,8 \times 0,050 \approx \SI{-0,79}{m.s^{-1}}.
\]
Après un quart de période, le solide repasse par la position d'équilibre avec sa vitesse maximale $v_{\max} = \omega_0 X_m \approx \SI{0,79}{m.s^{-1}}$, en sens négatif : il revient vers les élongations négatives.
\end{enumerate}
\end{exoresolu}

\begin{savaistu}[Pourquoi les ingénieurs adorent ce modèle]
Le pendule élastique horizontal est le modèle universel de toute petite oscillation : suspension de véhicule (les amortisseurs des moto-taxis de Douala), sismographes, microbalances, et même les vibrations des atomes dans un cristal. Dès qu'un système est rappelé vers sa position d'équilibre par un effet proportionnel à l'écart, il obéit à la même équation $\ddot{x} + \omega_0^2 x = 0$. Maîtriser ce cas simple, c'est déjà comprendre une grande partie de la physique des vibrations.
\end{savaistu}

\begin{aretenir}[L'essentiel de la section]
\begin{itemize}
\item Le pendule élastique horizontal sans frottement obéit à $\ddot{x} + \omega_0^2 x = 0$ avec $\omega_0 = \sqrt{k/m}$.
\item La solution est $x(t) = X_m\cos(\omega_0 t + \varphi)$ : oscillations harmoniques non amorties.
\item Période propre $T_0 = 2\pi\sqrt{m/k}$, indépendante de l'amplitude (isochronisme).
\item $X_m$ et $\varphi$ se déterminent à partir des conditions initiales : $X_m = \sqrt{x_0^2 + v_0^2/\omega_0^2}$ et $\tan\varphi = -v_0/(\omega_0 x_0)$.
\end{itemize}
\end{aretenir}

\coursec{Le pendule élastique vertical}

Dans la section précédente, nous avons étudié le pendule élastique \emph{horizontal} : le solide glissait sans frottement sur un support horizontal, et le poids était compensé à tout instant par la réaction du support. Une question naturelle se pose alors : que se passe-t-il si l'on suspend le ressort \emph{verticalement} ? Le poids n'est plus compensé par une réaction, il agit directement le long de l'axe du mouvement. Va-t-il modifier la nature des oscillations ? La période va-t-elle changer ? C'est à ces questions, d'une grande élégance physique, que cette section est consacrée. La réponse, tu vas le voir, est remarquable : la pesanteur \emph{décale} la position d'équilibre, mais ne modifie \emph{ni la pulsation ni la période} des oscillations.

\courssub{Dispositif expérimental et mise en évidence de l'équilibre}

\begin{definition}[Pendule élastique vertical]
Un \cle{pendule élastique vertical} est constitué d'un ressort à spires non jointives, de masse négligeable et de constante de raideur $k$, suspendu par son extrémité supérieure à un support fixe, et portant à son extrémité inférieure un solide $(S)$ de masse $m$, assimilé à un point matériel $G$. Le solide peut coulisser verticalement le long d'une tige de guidage (ou osciller librement) sans frottement.
\end{definition}

On note $\ell_0$ la longueur à vide du ressort. Lorsqu'on accroche le solide et qu'on attend que le système se stabilise, le ressort s'allonge : sa longueur devient $\ell_0 + \Delta\ell_0$, où $\Delta\ell_0$ est l'\cle{allongement à l'équilibre}. C'est une observation quotidienne : pense à la balance à ressort (le « peson ») utilisée sur les marchés de Yaoundé ou de Douala pour peser le poisson ou le manioc — plus la masse accrochée est grande, plus le ressort s'allonge.

\begin{experience}[Mise en évidence de l'allongement à l'équilibre]
\textbf{Matériel :} un ressort, une potence, une masse marquée $m = 100~\mathrm{g}$, une règle graduée.

\textbf{Protocole :}
\begin{enumerate}
\item Suspendre le ressort seul et repérer la position de son extrémité libre : longueur à vide $\ell_0$.
\item Accrocher la masse $m$ et attendre l'immobilisation : le ressort s'allonge de $\Delta\ell_0$.
\item Écarter légèrement la masse vers le bas et la lâcher : elle oscille verticalement autour de sa position d'équilibre.
\end{enumerate}

\textbf{Observations :} l'allongement $\Delta\ell_0$ est proportionnel à la masse accrochée. Les oscillations sont périodiques et s'effectuent \emph{autour de la position d'équilibre}, et non autour de la position de l'extrémité libre du ressort à vide.

\textbf{Interprétation :} à l'équilibre, la tension du ressort compense exactement le poids. Lors des oscillations, c'est l'écart par rapport à cette position d'équilibre qui gouverne le mouvement.
\end{experience}

\courssub{Condition d'équilibre}

\begin{propriete}[Condition d'équilibre du pendule élastique vertical]
À l'équilibre, l'allongement $\Delta\ell_0$ du ressort et la masse $m$ du solide sont liés par :
\[ k\,\Delta\ell_0 = m\,g \]
où $k$ est la constante de raideur du ressort (en $\mathrm{N\,m^{-1}}$) et $g$ l'intensité de la pesanteur (en $\mathrm{N\,kg^{-1}}$).
\end{propriete}

\textbf{Démonstration (exigible).} Étudions le système $\{$solide $(S)\}$ dans le référentiel terrestre supposé galiléen, à l'équilibre. Le solide est soumis à deux forces :
\begin{itemize}
\item son \cle{poids} $\vec{P} = m\vec{g}$, vertical, dirigé vers le bas ;
\item la \cle{tension du ressort} $\vec{T}_0$, verticale, dirigée vers le haut, de norme $T_0 = k\,\Delta\ell_0$.
\end{itemize}
Le solide étant au repos dans un référentiel galiléen, la première loi de Newton (principe d'inertie) s'applique :
\[ \vec{P} + \vec{T}_0 = \vec{0} \quad\Longrightarrow\quad T_0 = P \quad\Longrightarrow\quad k\,\Delta\ell_0 = m\,g. \]
Cette relation fondamentale sera la clé de toute la démonstration qui suit.

\begin{attention}[Vérification dimensionnelle]
La relation $k\Delta\ell_0 = mg$ est homogène : $[k\,\Delta\ell_0] = \mathrm{N\,m^{-1}} \times \mathrm{m} = \mathrm{N}$ et $[mg] = \mathrm{kg} \times \mathrm{m\,s^{-2}} = \mathrm{N}$. Prends l'habitude de vérifier systématiquement l'homogénéité de chaque relation avant de l'utiliser : c'est un réflexe qui sauve des points au Baccalauréat.
\end{attention}

\courssub{Choix du repère et bilan des forces en mouvement}

Pour étudier les oscillations, on choisit un axe vertical $(O;x)$ avec une convention qui va tout simplifier :

\begin{aretenir}[Choix du repère]
On choisit l'\cle{origine $O$ à la position d'équilibre} du centre d'inertie $G$ du solide, et l'axe $(O;x)$ vertical, \cle{orienté vers le bas}. L'abscisse $x$ de $G$ représente alors l'écart algébrique par rapport à la position d'équilibre : c'est l'\cle{élongation}.
\end{aretenir}

Lorsque le solide est à l'abscisse $x$ (comptée depuis l'équilibre, vers le bas), l'allongement total du ressort n'est plus $\Delta\ell_0$ mais :
\[ \Delta\ell = \Delta\ell_0 + x. \]
En effet, si le solide descend de $x$ ($x>0$), le ressort s'allonge davantage ; s'il monte ($x<0$), le ressort se détend.

Le solide en mouvement est soumis à :
\begin{itemize}
\item son poids $\vec{P} = m\vec{g}$, de projection $P_x = +mg$ (l'axe est orienté vers le bas) ;
\item la tension du ressort $\vec{T}$, dirigée vers le haut, de norme $T = k(\Delta\ell_0 + x)$, donc de projection $T_x = -k(\Delta\ell_0 + x)$.
\end{itemize}

\begin{popfigure}
\begin{center}
\begin{tikzpicture}[scale=0.9, every node/.style={font=\small}]
% ===== État 1 : ressort à vide =====
\begin{scope}[xshift=0cm]
\draw[thick, popdark] (-1,0) -- (1,0);
\foreach \x in {-0.8,-0.4,0,0.4,0.8} \draw[popdark] (\x,0) -- (\x+0.15,0.2);
\draw[decorate, decoration={coil, aspect=0.4, segment length=6pt, amplitude=5pt}, popblue, thick] (0,0) -- (0,-1.6);
\fill[popblue] (0,-1.6) circle (1.5pt);
\draw[<->, popmuted] (0.7,0) -- (0.7,-1.6) node[midway, right]{$\ell_0$};
\node[below, popdark] at (0,-2.1) {\textbf{Ressort à vide}};
\end{scope}
% ===== État 2 : équilibre =====
\begin{scope}[xshift=4.2cm]
\draw[thick, popdark] (-1,0) -- (1,0);
\foreach \x in {-0.8,-0.4,0,0.4,0.8} \draw[popdark] (\x,0) -- (\x+0.15,0.2);
\draw[decorate, decoration={coil, aspect=0.4, segment length=8pt, amplitude=5pt}, popblue, thick] (0,0) -- (0,-2.4);
\draw[fill=poporangeL, draw=poporange, thick] (-0.35,-2.4) rectangle (0.35,-3.1);
\node at (0,-2.75) {$m$};
\draw[<->, popmuted] (0.9,-1.6) -- (0.9,-2.75) node[midway, right]{$\Delta\ell_0$};
\draw[dashed, popmuted] (-0.9,-1.6) -- (0.9,-1.6);
\draw[->, popgreen, very thick] (0,-2.75) -- (0,-1.9) node[right]{$\vec{T}_0$};
\draw[->, poppink, very thick] (0,-2.75) -- (0,-3.6) node[right]{$\vec{P}$};
\fill[popink] (-0.9,-2.75) circle (1.5pt) node[left]{$O$};
\node[below, popdark] at (0,-3.9) {\textbf{Équilibre} : $k\Delta\ell_0 = mg$};
\end{scope}
% ===== État 3 : oscillation =====
\begin{scope}[xshift=8.6cm]
\draw[thick, popdark] (-1,0) -- (1,0);
\foreach \x in {-0.8,-0.4,0,0.4,0.8} \draw[popdark] (\x,0) -- (\x+0.15,0.2);
\draw[decorate, decoration={coil, aspect=0.4, segment length=9.5pt, amplitude=5pt}, popblue, thick] (0,0) -- (0,-3.15);
\draw[fill=poporangeL, draw=poporange, thick] (-0.35,-3.15) rectangle (0.35,-3.85);
\node at (0,-3.5) {$m$};
\draw[dashed, popmuted] (-0.9,-2.75) -- (1.6,-2.75);
\fill[popink] (-0.9,-2.75) circle (1.5pt) node[left]{$O$};
\draw[<->, poppurple] (1.4,-2.75) -- (1.4,-3.5) node[midway, right]{$x$};
\draw[->, popgreen, very thick] (0,-3.5) -- (0,-2.3) node[right]{$\vec{T}$};
\draw[->, poppink, very thick] (0,-3.5) -- (0,-4.35) node[right]{$\vec{P}$};
\draw[->, popdark, thick] (1.9,-2.6) -- (1.9,-3.4) node[midway, right]{$x$ (vers le bas)};
\node[below, popdark] at (0,-4.6) {\textbf{Oscillation} : allongement $\Delta\ell_0 + x$};
\end{scope}
\end{tikzpicture}
\end{center}
\legende{Les trois états du pendule élastique vertical : à vide, à l'équilibre (le poids est compensé par la tension $\vec{T}_0$), et en cours d'oscillation à l'élongation $x$ mesurée depuis la position d'équilibre $O$.}
\end{popfigure}

\courssub{Équation différentielle du mouvement}

\begin{propriete}[Équation différentielle — démonstration exigible au Bac]
Rapporté à sa position d'équilibre, le pendule élastique vertical obéit à l'équation différentielle :
\[ \ddot{x} + \frac{k}{m}\,x = 0 \qquad\text{soit}\qquad \ddot{x} + \omega_0^2\,x = 0 \quad \text{avec} \quad \omega_0 = \sqrt{\frac{k}{m}}. \]
C'est exactement la même équation que celle du pendule élastique horizontal.
\end{propriete}

\begin{experience}[Démonstration guidée pas à pas]
\textbf{Étape 1 — Système et référentiel.} Système : le solide $(S)$ de masse $m$. Référentiel terrestre supposé galiléen. Repère : axe $(O;x)$ vertical descendant, $O$ à la position d'équilibre.

\textbf{Étape 2 — Bilan des forces.} Le solide est soumis à son poids $\vec{P}$ (projection $+mg$) et à la tension du ressort $\vec{T}$ (projection $-k(\Delta\ell_0 + x)$). Les frottements sont négligés.

\textbf{Étape 3 — Deuxième loi de Newton.} Appliquée au centre d'inertie $G$, projetée sur l'axe $(O;x)$ :
\[ m\,\ddot{x} = mg - k(\Delta\ell_0 + x). \]

\textbf{Étape 4 — Simplification grâce à la condition d'équilibre.} Développons le second membre :
\[ m\,\ddot{x} = mg - k\,\Delta\ell_0 - k\,x. \]
Or la condition d'équilibre établie plus haut donne $k\,\Delta\ell_0 = mg$. Les deux premiers termes se compensent exactement :
\[ mg - k\,\Delta\ell_0 = 0 \quad\Longrightarrow\quad m\,\ddot{x} = -k\,x. \]

\textbf{Étape 5 — Forme canonique.} En divisant par $m$ :
\[ \ddot{x} + \frac{k}{m}\,x = 0. \]
C'est l'équation différentielle caractéristique d'un \cle{oscillateur harmonique} de pulsation propre $\omega_0 = \sqrt{k/m}$.
\end{experience}

\begin{attention}[Erreur fréquente : le mauvais choix d'origine]
Beaucoup d'élèves prennent l'origine à l'extrémité libre du ressort \emph{à vide}. L'abscisse $X$ mesurée depuis cette origine vérifie alors $m\ddot{X} = mg - kX$, soit $\ddot{X} + \dfrac{k}{m}X = g$ : le second membre $g$ parasite tout et l'équation n'est plus celle d'un oscillateur harmonique « pur ». Il faut alors effectuer le changement de variable $x = X - \Delta\ell_0$ pour retrouver la forme canonique. \textbf{Retiens la règle d'or : pour un oscillateur, on choisit toujours l'origine à la position d'équilibre.} L'équation différentielle est alors immédiatement homogène, et la solution s'écrit directement.
\end{attention}

\courssub{Solution de l'équation et période propre}

L'équation $\ddot{x} + \omega_0^2 x = 0$ admet la solution générale :
\[ x(t) = X_m\cos(\omega_0 t + \varphi), \]
où l'amplitude $X_m$ et la phase à l'origine $\varphi$ se déterminent à partir des conditions initiales $x(0)$ et $\dot{x}(0)$, exactement comme pour le pendule horizontal.

\begin{propriete}[Période propre du pendule élastique vertical]
La période propre des oscillations du pendule élastique vertical est :
\[ T_0 = \frac{2\pi}{\omega_0} = 2\pi\sqrt{\frac{m}{k}}. \]
Elle ne dépend que de la masse $m$ et de la raideur $k$ : \textbf{elle est indépendante de $g$ et de l'amplitude} (isochronisme des petites et grandes oscillations, tant que le ressort reste dans son domaine élastique).
\end{propriete}

\begin{aretenir}[Le rôle de la pesanteur]
La pesanteur ne modifie \emph{ni la pulsation propre $\omega_0$ ni la période propre $T_0$} du pendule élastique : elle se contente de \emph{décaler la position d'équilibre} de $\Delta\ell_0 = \dfrac{mg}{k}$ vers le bas. Autour de cette nouvelle position d'équilibre, tout se passe comme si la pesanteur n'existait pas. Un même ressort portant la même masse oscille avec la même période, qu'il soit horizontal sur une table ou vertical suspendu à une potence — et même sur la Lune, seule la position d'équilibre changerait !
\end{aretenir}

\begin{exemplebox}[Exemple chiffré : détermination de $k$ puis de $T_0$]
Un ressort suspendu verticalement s'allonge de $\Delta\ell_0 = 4{,}0~\mathrm{cm}$ lorsqu'on y accroche une masse $m = 100~\mathrm{g}$. On prend $g = 9{,}8~\mathrm{N\,kg^{-1}}$.

\textbf{1. Constante de raideur.} La condition d'équilibre $k\,\Delta\ell_0 = mg$ donne :
\[ k = \frac{mg}{\Delta\ell_0} = \frac{0{,}100 \times 9{,}8}{0{,}040} = 24{,}5~\mathrm{N\,m^{-1}}. \]

\textbf{2. Période propre.} On écarte légèrement la masse de sa position d'équilibre et on la lâche :
\[ T_0 = 2\pi\sqrt{\frac{m}{k}} = 2\pi\sqrt{\frac{0{,}100}{24{,}5}} = 2\pi\sqrt{4{,}08\times 10^{-3}} \approx 2\pi \times 0{,}0639 \approx 0{,}40~\mathrm{s}. \]
La masse effectue donc environ $2{,}5$ oscillations par seconde.
\end{exemplebox}

\courssub{Application : mesure de $g$ ou de $k$ par l'étude des oscillations}

La double relation qui gouverne le pendule élastique vertical offre deux méthodes de mesure très utilisées au laboratoire.

\begin{methode}[Mesurer $k$ et $g$ avec un pendule élastique vertical]
\begin{enumerate}
\item \textbf{Méthode statique (mesure de $k$) :} mesurer l'allongement à l'équilibre $\Delta\ell_0$ pour une masse connue $m$, puis appliquer $k = \dfrac{mg}{\Delta\ell_0}$. Pour plus de précision, tracer $mg = f(\Delta\ell_0)$ pour plusieurs masses : la pente de la droite obtenue est $k$.
\item \textbf{Méthode dynamique (mesure de $T_0$) :} chronométrer la durée $\Delta t$ de $N = 10$ ou $20$ oscillations complètes, puis calculer $T_0 = \dfrac{\Delta t}{N}$ pour réduire l'incertitude.
\item \textbf{Exploitation :} si $k$ est connue (méthode statique), la mesure de $T_0$ donne la masse : $m = \dfrac{k\,T_0^2}{4\pi^2}$. Réciproquement, en combinant les deux relations $k\Delta\ell_0 = mg$ et $T_0 = 2\pi\sqrt{m/k}$, on élimine $m$ et $k$ :
\[ T_0 = 2\pi\sqrt{\frac{\Delta\ell_0}{g}} \quad\Longrightarrow\quad g = \frac{4\pi^2\,\Delta\ell_0}{T_0^2}. \]
Cette dernière formule permet de mesurer $g$ avec un simple ressort, une masse et un chronomètre !
\end{enumerate}
\end{methode}

\begin{exoresolu}[Détermination expérimentale de $g$]
Au laboratoire du lycée, un élève suspend une masse $m = 200~\mathrm{g}$ à un ressort vertical. Il mesure un allongement à l'équilibre $\Delta\ell_0 = 8{,}0~\mathrm{cm}$. Il écarte ensuite la masse de $2{,}0~\mathrm{cm}$ vers le bas, la lâche sans vitesse initiale, et chronomètre $\Delta t = 11{,}4~\mathrm{s}$ pour $N = 20$ oscillations complètes.

\textbf{1.} Calculer la période propre $T_0$ des oscillations.

\textbf{2.} En déduire la valeur expérimentale de $g$.

\textbf{3.} Écrire l'équation horaire $x(t)$ du mouvement (axe vertical descendant, origine à l'équilibre).

\tcblower

\textbf{1. Période propre.} La durée de 20 oscillations est $\Delta t = 11{,}4~\mathrm{s}$, donc :
\[ T_0 = \frac{\Delta t}{N} = \frac{11{,}4}{20} = 0{,}57~\mathrm{s}. \]

\textbf{2. Valeur de $g$.} En combinant la condition d'équilibre $k\Delta\ell_0 = mg$ et l'expression de la période $T_0 = 2\pi\sqrt{m/k}$, on obtient $T_0 = 2\pi\sqrt{\Delta\ell_0/g}$, d'où :
\[ g = \frac{4\pi^2\,\Delta\ell_0}{T_0^2} = \frac{4\pi^2 \times 0{,}080}{(0{,}57)^2} = \frac{3{,}16}{0{,}325} \approx 9{,}7~\mathrm{m\,s^{-2}}. \]
Valeur tout à fait compatible avec la valeur tabulée $g = 9{,}8~\mathrm{m\,s^{-2}}$, compte tenu des incertitudes de chronométrage.

\textbf{3. Équation horaire.} La solution générale est $x(t) = X_m\cos(\omega_0 t + \varphi)$ avec :
\[ \omega_0 = \frac{2\pi}{T_0} = \frac{2\pi}{0{,}57} \approx 11{,}0~\mathrm{rad\,s^{-1}}. \]
Conditions initiales : $x(0) = +2{,}0~\mathrm{cm}$ (l'axe est orienté vers le bas et la masse est écartée vers le bas) et $\dot{x}(0) = 0$ (lâchée sans vitesse initiale). Donc $X_m\cos\varphi = 2{,}0~\mathrm{cm}$ et $-\omega_0 X_m\sin\varphi = 0$, ce qui impose $\varphi = 0$ et $X_m = 2{,}0~\mathrm{cm}$ :
\[ x(t) = 2{,}0\cos(11{,}0\,t) \quad \text{avec } x \text{ en cm et } t \text{ en s}. \]
\end{exoresolu}

\begin{attention}[Pièges de rédaction à éviter]
\begin{itemize}
\item Ne jamais écrire que la tension du ressort vaut $T = kx$ : elle vaut $T = k(\Delta\ell_0 + x)$, car le ressort est \emph{déjà allongé} à l'équilibre. C'est la simplification par $k\Delta\ell_0 = mg$ qui fait « disparaître » les termes constants.
\item Ne pas oublier de convertir les longueurs en mètres : $\Delta\ell_0 = 4{,}0~\mathrm{cm} = 0{,}040~\mathrm{m}$.
\item La relation $T_0 = 2\pi\sqrt{\Delta\ell_0/g}$ ressemble à celle du pendule simple, $T_0 = 2\pi\sqrt{\ell/g}$, mais il ne faut pas les confondre : ici $\Delta\ell_0$ est un \emph{allongement} de ressort, pas une longueur de fil.
\item La période ne dépend pas de l'amplitude : chronométrer avec des oscillations de $2~\mathrm{cm}$ ou de $4~\mathrm{cm}$ donne le même $T_0$ (dans la limite élastique du ressort).
\end{itemize}
\end{attention}

\begin{savaistu}[Les suspensions de nos véhicules]
Les suspensions des voitures et des mototaxis qui sillonnent les routes de Douala, Bafoussam ou Garoua sont des \emph{pendules élastiques verticaux amortis} : chaque roue est reliée à la caisse par un ressort (ou une lame) associé à un amortisseur. Lorsque le véhicule passe sur une bosse ou un nid-de-poule, le ressort s'enfonce puis oscille : sans amortisseur, la caisse rebondirait longtemps, rendant le trajet inconfortable et dangereux. L'amortisseur dissipe l'énergie et ramène rapidement le véhicule à sa position d'équilibre — nous étudierons ce régime « amorti » dans la section 5. Les ingénieurs choisissent la raideur $k$ des ressorts en fonction de la masse $m$ du véhicule pour obtenir une période propre $T_0 = 2\pi\sqrt{m/k}$ de l'ordre de la seconde, valeur jugée confortable pour le corps humain. Le même principe équipe les selles de vélo à ressort et les balances à ressort des marchés !
\end{savaistu}

\begin{exercice}[Oscillations verticales et pesanteur]
Un ressort de constante de raideur $k = 50~\mathrm{N\,m^{-1}}$ est suspendu verticalement. On y accroche un solide de masse $m = 500~\mathrm{g}$. On donne $g = 9{,}8~\mathrm{N\,kg^{-1}}$.

\textbf{1.} Calculer l'allongement du ressort à l'équilibre.

\textbf{2.} On écarte le solide de $3{,}0~\mathrm{cm}$ vers le haut à partir de sa position d'équilibre et on le lâche sans vitesse initiale. Établir l'équation différentielle du mouvement et donner l'équation horaire $x(t)$ (axe vertical descendant, origine à l'équilibre).

\textbf{3.} Calculer la période propre $T_0$. Cette période serait-elle différente si l'expérience était réalisée sur la Lune ($g_{\text{L}} = 1{,}6~\mathrm{N\,kg^{-1}}$) ? Justifier.
\end{exercice}

\begin{corrige}[Exercice : Oscillations verticales et pesanteur]
\textbf{1.} Condition d'équilibre : $k\,\Delta\ell_0 = mg$, donc
\[ \Delta\ell_0 = \frac{mg}{k} = \frac{0{,}500 \times 9{,}8}{50} = 0{,}098~\mathrm{m} = 9{,}8~\mathrm{cm}. \]

\textbf{2.} La deuxième loi de Newton projetée sur l'axe vertical descendant donne $m\ddot{x} = mg - k(\Delta\ell_0 + x)$, qui se simplifie par la condition d'équilibre en $\ddot{x} + \dfrac{k}{m}x = 0$, avec $\omega_0 = \sqrt{k/m} = \sqrt{50/0{,}500} = 10~\mathrm{rad\,s^{-1}}$. Conditions initiales : $x(0) = -3{,}0~\mathrm{cm}$ (écart vers le \emph{haut}, axe descendant) et $\dot{x}(0) = 0$, d'où $X_m = 3{,}0~\mathrm{cm}$ et $\varphi = \pi$ :
\[ x(t) = 3{,}0\cos(10\,t + \pi) = -3{,}0\cos(10\,t) \quad (x \text{ en cm}). \]

\textbf{3.} $T_0 = \dfrac{2\pi}{\omega_0} = \dfrac{2\pi}{10} \approx 0{,}63~\mathrm{s}$. Sur la Lune, la période serait \emph{identique}, car $T_0 = 2\pi\sqrt{m/k}$ ne dépend pas de $g$. Seule la position d'équilibre changerait : $\Delta\ell_0' = \dfrac{mg_{\text{L}}}{k} = \dfrac{0{,}500 \times 1{,}6}{50} = 1{,}6~\mathrm{cm}$. Le ressort serait moins allongé, mais oscillerait à la même cadence.
\end{corrige}

\coursec{Pendule pesant et pendule simple}

Dans les sections précédentes, nous avons étudié le pendule élastique, horizontal puis vertical : dans les deux cas, l'oscillateur est un solide en \emph{translation}, rappelé vers sa position d'équilibre par la force de tension d'un ressort, et nous avons établi l'équation différentielle $\ddot{x} + \omega_0^2 x = 0$ de période propre $T_0 = 2\pi\sqrt{m/k}$. Mais les oscillateurs mécaniques de la vie quotidienne ne sont pas tous des translations : le balancier d'une horloge, l'enfant sur une balançoire à Nkolbisson, la grue de chantier dont la charge oscille, sont des solides en \emph{rotation} autour d'un axe fixe, rappelés par leur propre \cle{poids}. C'est à cette nouvelle famille d'oscillateurs — les pendules pesants — que cette section est consacrée. Nous allons voir qu'avec une hypothèse simple (les petites oscillations), on retrouve exactement la même équation différentielle harmonique, avec une période propre d'une remarquable simplicité.

\courssub{Le pendule pesant : définition et grandeurs caractéristiques}

\begin{definition}[Pendule pesant]
Un \cle{pendule pesant} est un solide mobile autour d'un axe horizontal fixe $(\Delta)$, ne passant pas par son centre d'inertie $G$, et soumis à son poids. Écarté de sa position d'équilibre stable (où $G$ est à la verticale sous l'axe), il oscille sous l'action du moment du poids.
\end{definition}

Trois grandeurs décrivent complètement le pendule pesant :
\begin{itemize}
\item le \cle{moment d'inertie} $J_\Delta$ du solide par rapport à l'axe $(\Delta)$, en $\mathrm{kg\cdot m^2}$, qui mesure son « inertie de rotation » ;
\item la distance $d = OG$ entre l'axe et le centre d'inertie ;
\item la masse $m$ du solide, le champ de pesanteur $g$ étant supposé uniforme.
\end{itemize}

La position du pendule est repérée par l'angle $\theta$ que fait $OG$ avec la verticale descendante. Le poids $\vec{P} = m\vec{g}$ appliqué en $G$ a pour moment par rapport à $(\Delta)$ :
\[ \mathcal{M}_\Delta(\vec{P}) = -mgd\sin\theta \]
Le signe $-$ traduit le caractère de \emph{rappel} : quand $\theta > 0$, le moment est négatif et ramène le pendule vers $\theta = 0$. Le moment de la réaction de l'axe est nul (force dont la droite d'action rencontre l'axe).

\begin{propriete}[Équation différentielle du pendule pesant]
En négligeant les frottements, la relation fondamentale de la dynamique de rotation appliquée au pendule pesant s'écrit :
\[ J_\Delta\,\ddot{\theta} = -mgd\sin\theta \qquad\text{soit}\qquad \ddot{\theta} + \frac{mgd}{J_\Delta}\sin\theta = 0 \]
Cette équation n'est \emph{pas} linéaire à cause du $\sin\theta$ : en toute rigueur, le pendule pesant n'est pas un oscillateur harmonique. Elle le devient pour les petites oscillations.
\end{propriete}

\courssub{Le pendule simple : un modèle idéal}

\begin{definition}[Pendule simple]
Un \cle{pendule simple} est un point matériel $M$ de masse $m$, suspendu à un fil inextensible, sans masse, de longueur $\ell$, fixé en $O$, et pouvant osciller sans frottement dans un plan vertical. C'est le cas limite du pendule pesant pour lequel toute la masse est concentrée en $G \equiv M$ : $d = \ell$ et $J_\Delta = m\ell^2$.
\end{definition}

Une bille de pétanque suspendue à une ficelle, un plomb de maçon au bout de son cordeau : ce sont d'excellentes approximations du pendule simple, dès lors que le fil est léger devant la masse et celle-ci petite devant $\ell$.

\begin{popfigure}
\begin{center}
\begin{tikzpicture}[scale=2.2,>=Stealth]
% pivot et support
\fill[poppurple!25] (-0.9,0) rectangle (0.9,0.12);
\foreach \x in {-0.8,-0.6,...,0.8} \draw[popmuted] (\x,0.12) -- (\x+0.1,0);
\coordinate (O) at (0,0);
% verticale de reference
\draw[dashed,popmuted] (O) -- (0,-2.9) node[below,popdark]{\scriptsize verticale};
% pendule
\coordinate (M) at (35:2.6);
\draw[thick,popblue] (O) -- (M) node[midway,above left,popdark]{$\ell$};
\fill[poporange] (M) circle (0.09);
\node[popdark,right] at (M) {$M$};
% angle theta
\draw[popgreen,thick] (0,-1.1) arc (-90:-55:1.1) node[midway,below right,popgreen]{$\theta$};
% poids
\draw[->,very thick,popink] (M) -- ++(0,-1.5) node[below]{$\vec{P}$};
% composantes tangentielle et normale
\draw[->,very thick,poppink] (M) -- ++(-55:1.44) node[right]{$\vec{P}_t$};
\draw[->,very thick,popgold] (M) -- ++(215:0.87) node[left]{$\vec{P}_n$};
\draw[dashed,popmuted] ($(M)+(0,-1.5)$) -- ++(-55:0.87);
\draw[dashed,popmuted] ($(M)+(0,-1.5)$) -- ++(215:0.87);
% tension
\draw[->,very thick,popblue] (M) -- ($(O)!0.35!(M)$) node[midway,above right,popblue]{$\vec{T}$};
\fill (O) circle (0.03);
\node[left,popdark] at (O) {$O$};
\end{tikzpicture}
\end{center}
\legende{Pendule simple écarté d'un angle $\theta$. Le poids se décompose en composante normale $\vec{P}_n$ (compensée avec la tension $\vec{T}$ sur la direction du fil) et composante tangentielle $\vec{P}_t$, de norme $mg\sin\theta$, qui joue le rôle de force de rappel.}
\end{popfigure}

\courssub{Approximation des petites oscillations}

Retrouvons l'équation du mouvement par une deuxième méthode, très formatrice : le principe fondamental de la dynamique appliqué au point $M$, projeté sur la tangente à la trajectoire (le cercle de rayon $\ell$). La tension $\vec{T}$ et la composante normale $\vec{P}_n$ n'ont pas de projection tangentielle ; seule la composante tangentielle du poids intervient : $P_t = -mg\sin\theta$. L'accélération tangentielle vaut $a_t = \ell\ddot{\theta}$ (car l'abscisse curviligne est $s = \ell\theta$). D'où :
\[ m\ell\ddot{\theta} = -mg\sin\theta \quad\Longleftrightarrow\quad \ddot{\theta} + \frac{g}{\ell}\sin\theta = 0 \]
On retrouve l'équation du pendule pesant avec $J_\Delta = m\ell^2$ et $d = \ell$ : tout est cohérent.

Cette équation ne se résout pas simplement... sauf si l'angle reste petit. Pour $\theta$ \textbf{exprimé en radians}, le développement limité donne $\sin\theta = \theta - \dfrac{\theta^3}{6} + \dots$, donc :

\begin{aretenir}[Approximation des petits angles]
Pour $\theta$ en radians et $|\theta| \lesssim 10^{\circ}$ (soit environ $\num{0,17}$ rad) :
\[ \sin\theta \approx \theta \qquad\text{(erreur inférieure à 0{,}5\,\%)} \]
C'est l'\cle{approximation des petites oscillations}.
\end{aretenir}

\begin{popfigure}
\begin{center}
\begin{tikzpicture}
\begin{axis}[width=0.85\linewidth,height=6.2cm,
xlabel={$\theta$ (rad)},ylabel={},
xmin=-0.6,xmax=0.6,ymin=-0.6,ymax=0.6,
legend pos=north west,grid=both,grid style={popline!40},
samples=200]
\addplot[popcurve,domain=-0.6:0.6]{sin(deg(x))};
\addlegendentry{$\sin\theta$}
\addplot[popcurve,poporange,dashed,domain=-0.6:0.6]{x};
\addlegendentry{$\theta$}
\draw[popgreen,thick,dotted] (axis cs:0.1745,-0.6) -- (axis cs:0.1745,0.1745);
\draw[popgreen,thick,dotted] (axis cs:-0.1745,-0.6) -- (axis cs:-0.1745,-0.1745);
\node[popgreen,font=\scriptsize] at (axis cs:0.30,-0.45) {$10^{\circ} \approx 0{,}175$ rad};
\end{axis}
\end{tikzpicture}
\end{center}
\legende{Comparaison de $\sin\theta$ et de $\theta$ (en radians). Les deux courbes sont quasiment confondues entre $-10^{\circ}$ et $+10^{\circ}$ : l'approximation $\sin\theta \approx \theta$ y est excellente.}
\end{popfigure}

\begin{attention}[Radians obligatoires !]
L'approximation $\sin\theta \approx \theta$ n'est valable que si $\theta$ est exprimé en \textbf{radians}. Dire « $\sin 10 = 10$ » parce que l'angle vaut $10^{\circ}$ est une erreur grave : $\sin 10^{\circ} \approx \num{0,174}$, alors que $10^{\circ} = \num{0,1745}$ rad. Pense aussi à régler ta calculatrice en mode \textbf{rad} avant tout calcul de $\sin\theta$ ou de $\omega_0 t + \varphi$.
\end{attention}

\courssub{Équation linéarisée et période propre}

Avec $\sin\theta \approx \theta$, l'équation différentielle devient linéaire :

\begin{propriete}[Période propre du pendule simple — démonstration guidée]
Pour les petites oscillations, le pendule simple vérifie :
\[ \ddot{\theta} + \frac{g}{\ell}\,\theta = 0 \]
C'est l'équation d'un oscillateur harmonique de pulsation propre $\omega_0 = \sqrt{\dfrac{g}{\ell}}$, de solution $\theta(t) = \theta_m\cos(\omega_0 t + \varphi)$ et de période propre :
\[ \boxed{\ T_0 = 2\pi\sqrt{\frac{\ell}{g}}\ } \]
\end{propriete}

\begin{experience}[Démonstration guidée pas à pas]
\begin{enumerate}
\item \textbf{Bilan des forces} sur le point $M$ : poids $\vec{P} = m\vec{g}$ et tension $\vec{T}$ du fil.
\item \textbf{PFD} : $m\vec{a} = \vec{P} + \vec{T}$, projeté sur la tangente orientée dans le sens des $\theta$ croissants : $m\ell\ddot{\theta} = -mg\sin\theta$.
\item \textbf{Linéarisation} : pour $|\theta| \lesssim 10^{\circ}$, $\sin\theta \approx \theta$, d'où $\ddot{\theta} + \dfrac{g}{\ell}\theta = 0$.
\item \textbf{Identification} à la forme canonique $\ddot{\theta} + \omega_0^2\theta = 0$ : $\omega_0^2 = \dfrac{g}{\ell}$.
\item \textbf{Conclusion} : $T_0 = \dfrac{2\pi}{\omega_0} = 2\pi\sqrt{\dfrac{\ell}{g}}$.
\end{enumerate}
Vérification dimensionnelle : $\left[\sqrt{\ell/g}\right] = \sqrt{\dfrac{\mathrm{m}}{\mathrm{m\cdot s^{-2}}}} = \sqrt{\mathrm{s^2}} = \mathrm{s}$. La formule est homogène.
\end{experience}

\begin{aretenir}[Isochronisme des petites oscillations]
La période propre $T_0 = 2\pi\sqrt{\ell/g}$ ne dépend \textbf{ni de la masse} $m$, \textbf{ni de l'amplitude} $\theta_m$ (tant que celle-ci reste petite) : c'est l'\cle{isochronisme des petites oscillations}. Seuls la longueur $\ell$ et le champ de pesanteur $g$ fixent la période.
\end{aretenir}

\begin{exemplebox}[Le pendule qui « bat la seconde »]
Pour $\ell = \SI{1,0}{m}$ et $g = \SI{9,8}{N/kg}$ :
\[ T_0 = 2\pi\sqrt{\frac{1{,}0}{9{,}8}} \approx \SI{2,0}{s} \]
Un pendule d'environ un mètre effectue une demi-période (un « battement ») en une seconde : c'est le principe des horloges à balancier d'antan.
\end{exemplebox}

\begin{savaistu}[Huygens et la mesure du temps]
En 1673, Christiaan Huygens exploita l'isochronisme des petites oscillations pour construire les premières horloges à pendule précises, divisant par dix l'erreur de mesure du temps. Deux siècles plus tard, le pendule de Foucault (1851), long de 67 m, rendit visible la rotation de la Terre. Aujourd'hui encore, mesurer $T_0$ et $\ell$ reste l'une des méthodes les plus simples pour déterminer $g$ — y compris dans un laboratoire de lycée à Yaoundé ou à Douala.
\end{savaistu}

\courssub{Application : détermination expérimentale de g}

Puisque $T_0^2 = \dfrac{4\pi^2}{g}\,\ell$, le carré de la période est une fonction affine croissante de $\ell$, passant par l'origine. C'est la base d'une mesure précise de $g$.

\begin{methode}[Déterminer g expérimentalement]
\begin{enumerate}
\item Réaliser un pendule simple (petite bille dense, fil fin) et mesurer sa longueur $\ell$ au millimètre près, du point de suspension au centre de la bille.
\item Écarter le pendule d'un angle inférieur à $10^{\circ}$, le lâcher sans vitesse initiale, et mesurer la durée de $N = 10$ ou $20$ périodes : $T_0 = \dfrac{t}{N}$ (cela réduit l'erreur de chronométrage).
\item Recommencer pour au moins cinq longueurs différentes.
\item Tracer $T_0^2$ en fonction de $\ell$ : les points s'alignent sur une droite de pente $a = \dfrac{4\pi^2}{g}$.
\item En déduire $g = \dfrac{4\pi^2}{a}$ et comparer à la valeur de référence $\SI{9,81}{N/kg}$.
\end{enumerate}
\end{methode}

\begin{exoresolu}[Mesure de g au laboratoire]
Un groupe d'élèves de Terminale C mesure, pour un pendule de longueur $\ell = \SI{0,50}{m}$, la durée de 20 oscillations complètes : $t = \SI{28,4}{s}$. En traçant $T_0^2$ en fonction de $\ell$ pour plusieurs longueurs, ils obtiennent une droite de pente $a = \SI{4,04}{s^2/m}$. Déterminer la période pour $\ell = \SI{0,50}{m}$, puis la valeur expérimentale de $g$.
\tcblower
\textbf{Période propre.} La durée mesurée correspond à 20 périodes :
\[ T_0 = \frac{t}{N} = \frac{28{,}4}{20} = \SI{1,42}{s} \]
Vérification par la formule : $2\pi\sqrt{0{,}50/9{,}81} \approx \SI{1,42}{s}$. Cohérent.

\textbf{Valeur de g.} De $T_0^2 = \dfrac{4\pi^2}{g}\ell$, la pente du graphe $T_0^2 = f(\ell)$ vaut $a = \dfrac{4\pi^2}{g}$, donc :
\[ g = \frac{4\pi^2}{a} = \frac{4\times (3{,}1416)^2}{4{,}04} \approx \frac{39{,}48}{4{,}04} \approx \SI{9,77}{N/kg} \]
L'écart relatif avec la valeur de référence est $\dfrac{9{,}81 - 9{,}77}{9{,}81} \approx 0{,}4\ \%$ : excellente mesure pour un TP de lycée.
\end{exoresolu}

\begin{attention}[Pièges fréquents au Bac]
\begin{itemize}
\item Ne pas confondre la période mesurée et la durée totale : si l'énoncé donne « 20 oscillations en $t$ secondes », diviser par 20 avant tout calcul.
\item La formule $T_0 = 2\pi\sqrt{\ell/g}$ suppose les \textbf{petites oscillations} : pour $\theta_m = 60^{\circ}$, la période réelle est environ $7\ \%$ plus grande et l'isochronisme disparaît.
\item $\ell$ est la distance du point de suspension au \emph{centre} de la bille, pas la longueur du fil seul.
\item Dans $T_0^2 = \dfrac{4\pi^2}{g}\ell$, la pente est $\dfrac{4\pi^2}{g}$ et non $\dfrac{g}{4\pi^2}$ : vérifie l'axe des ordonnées ($T_0^2$) et celui des abscisses ($\ell$) avant d'exploiter le graphe.
\end{itemize}
\end{attention}

En résumé, le pendule simple aux petites oscillations est un oscillateur harmonique de période $T_0 = 2\pi\sqrt{\ell/g}$, indépendante de la masse et de l'amplitude : un résultat simple, puissant, et à la base de la mesure du temps pendant deux siècles. Il reste à examiner un troisième type de pendule — le pendule de torsion — où le rappel n'est plus dû au poids mais à l'élasticité de torsion d'un fil, puis à confronter tous ces oscillateurs idéaux à la réalité des frottements : ce sera l'objet des sections suivantes.

\coursec{Pendule de torsion et étude énergétique}

Dans la section précédente, nous avons étudié le pendule pesant et montré que, pour de petites oscillations, le pendule simple se comporte comme un oscillateur harmonique de période propre $T_0 = 2\pi\sqrt{\ell/g}$. Nous allons maintenant découvrir un troisième type d'oscillateur mécanique, le \cle{pendule de torsion}, qui joue un rôle historique considérable en physique, puis nous poserons une question plus profonde : que devient l'\cle{énergie} d'un oscillateur au cours du temps ? Cette étude énergétique, très prisée au Baccalauréat, donne un point de vue complémentaire et souvent plus rapide que l'étude dynamique.

\courssub{Le pendule de torsion}

\textbf{Dispositif et observation}\par
Suspendons un fil d'acier vertical à un support fixe et attachons à son extrémité inférieure une tige horizontale (ou un disque) homogène. Écartons la tige d'un angle $\theta_0$ autour de l'axe vertical $\Delta$ matérialisé par le fil, puis lâchons-la sans vitesse initiale. On observe que la tige oscille alternativement dans un sens puis dans l'autre : le fil se \emph{tord} et se \emph{détord}, exerçant sur la tige un couple qui la ramène toujours vers sa position d'équilibre. C'est l'analogue rotatif du ressort qui s'allonge et se comprime.

\begin{definition}[Pendule de torsion]
Un \cle{pendule de torsion} est constitué d'un solide de moment d'inertie $J_\Delta$ par rapport à un axe fixe $\Delta$, suspendu à un fil de torsion dont l'axe coïncide avec $\Delta$. Lorsque le solide tourne d'un angle $\theta$ par rapport à sa position d'équilibre, le fil exerce sur lui un \cle{couple de rappel} :
\[ \mathcal{M}_\Delta = -\,C\,\theta \]
où $C$ est la \cle{constante de torsion} du fil, en newton-mètre par radian (\si{N.m.rad^{-1}}). Le signe $-$ traduit le fait que le couple s'oppose toujours à la déformation, comme la tension $\vec{T} = -k\,x\,\vec{i}$ du ressort. Le radian étant sans dimension, $C$ s'exprime aussi en \si{N.m}.
\end{definition}

\begin{popfigure}
\begin{center}
\begin{tikzpicture}[scale=1.1]
% support vu de dessus : point de suspension
\fill[poppurple] (0,0) circle (2.2pt);
\draw[poppurple,thick] (0,0) circle (0.12);
% tige en position d'équilibre (pointillés)
\draw[popmuted,dashed,thick] (-2.2,0) -- (2.2,0);
\fill[popmuted] (-2.2,0) circle (2pt);
\fill[popmuted] (2.2,0) circle (2pt);
% tige écartée d'un angle theta
\draw[popblue,very thick] (150:2.2) -- (-30:2.2);
\fill[popblue] (150:2.2) circle (2.5pt);
\fill[popblue] (-30:2.2) circle (2.5pt);
% angle theta
\draw[poporange,thick,->] (1.3,0) arc (0:150:1.3);
\node[poporange] at (75:1.7) {$\theta$};
% sens du couple de rappel
\draw[popgreen,thick,->] (150:2.6) arc (150:60:2.6);
\node[popgreen] at (105:2.95) {$\mathcal{M} = -C\theta$};
% axe Delta
\node[poppurple] at (0.35,-0.35) {$\Delta$};
% légendes
\node[popmuted] at (2.2,-0.45) {\small position d'équilibre};
\node[popblue] at (-30:2.75) {\small tige};
\end{tikzpicture}
\end{center}
\legende{Pendule de torsion vu de dessus : la tige, écartée d'un angle $\theta$, subit un couple de rappel $\mathcal{M} = -C\theta$ qui la ramène vers sa position d'équilibre.}
\end{popfigure}

\textbf{Mise en équation}\par
Le solide est soumis à son poids et à la réaction de l'axe (de moments nuls par rapport à $\Delta$ car leurs droites d'action rencontrent l'axe), et au couple de torsion du fil. Appliquons le \cle{théorème du moment cinétique} (ou relation fondamentale de la dynamique de rotation) projeté sur l'axe $\Delta$ :
\[ J_\Delta\,\ddot{\theta} = \mathcal{M}_\Delta = -C\,\theta \]
d'où l'équation différentielle du mouvement :
\[ \boxed{\ \ddot{\theta} + \frac{C}{J_\Delta}\,\theta = 0\ } \]

C'est exactement la forme canonique $\ddot{\theta} + \omega_0^2\,\theta = 0$ rencontrée pour le pendule élastique, avec la pulsation propre :
\[ \omega_0 = \sqrt{\frac{C}{J_\Delta}} \]

\begin{propriete}[Période propre du pendule de torsion]
Le pendule de torsion effectue des oscillations harmoniques $\theta(t) = \theta_m\cos(\omega_0 t + \varphi)$ de période propre :
\[ T_0 = 2\pi\sqrt{\frac{J_\Delta}{C}} \]
La démonstration est \emph{strictement identique} à celle du pendule élastique : on vérifie que $\theta(t) = \theta_m\cos(\omega_0 t + \varphi)$ est solution de $\ddot{\theta} + \omega_0^2\theta = 0$ (car $\ddot{\theta} = -\omega_0^2\theta$), et la période s'obtient par $T_0 = 2\pi/\omega_0$. Analyse dimensionnelle : $[J_\Delta] = \si{kg.m^2}$, $[C] = \si{N.m.rad^{-1}} = \si{kg.m^2.s^{-2}}$, donc $[J_\Delta/C] = \si{s^2}$ : la racine carrée est bien homogène à un temps.
\end{propriete}

\begin{methode}[Tableau d'analogie translation $\leftrightarrow$ rotation]
Toutes les formules du pendule de torsion se retrouvent instantanément à partir de celles du pendule élastique grâce au dictionnaire suivant :
\begin{center}
\begin{tabularx}{\linewidth}{|l|X|X|}\hline
\rowcolor{popblueL} \textbf{Grandeur} & \textbf{Translation (ressort)} & \textbf{Rotation (torsion)} \\ \hline
Déplacement & $x$ (m) & $\theta$ (rad) \\ \hline
Vitesse & $v = \dot{x}$ & $\dot{\theta}$ (rad.s$^{-1}$) \\ \hline
Inertie & masse $m$ (kg) & $J_\Delta$ (\si{kg.m^2}) \\ \hline
Raideur & $k$ (\si{N.m^{-1}}) & $C$ (\si{N.m.rad^{-1}}) \\ \hline
Force / couple & $F = -k\,x$ & $\mathcal{M} = -C\,\theta$ \\ \hline
Loi de Newton & $m\ddot{x} = -k\,x$ & $J_\Delta\ddot{\theta} = -C\,\theta$ \\ \hline
Période propre & $T_0 = 2\pi\sqrt{m/k}$ & $T_0 = 2\pi\sqrt{J_\Delta/C}$ \\ \hline
Énergie cinétique & $\tfrac{1}{2}m v^2$ & $\tfrac{1}{2}J_\Delta\dot{\theta}^2$ \\ \hline
Énergie potentielle & $\tfrac{1}{2}k x^2$ & $\tfrac{1}{2}C\theta^2$ \\ \hline
\end{tabularx}
\end{center}
Retiens ce dictionnaire : il divise par deux le nombre de formules à mémoriser et te permet de vérifier rapidement tes résultats.
\end{methode}

\begin{savaistu}[La balance de torsion de Cavendish (1798)]
C'est grâce à un pendule de torsion que Henry Cavendish réalisa en 1798 l'une des plus belles expériences de l'histoire de la physique : la mesure de la constante de gravitation universelle $G$. Deux petites sphères de plomb, fixées aux extrémités d'une tige suspendue à un fil de torsion très fin, étaient attirées par deux grosses sphères de plomb voisines. La force d'attraction gravitationnelle, minuscule, faisait tourner la tige d'un angle mesurable jusqu'à ce que le couple de torsion $C\theta$ équilibre le couple gravitationnel. Connaissant $C$ (déterminé en mesurant la période $T_0 = 2\pi\sqrt{J_\Delta/C}$ du pendule !), Cavendish en déduisit $G \approx \SI{6,67e-11}{N.m^2.kg^{-2}}$ et, ce faisant, « pesa la Terre ». Le pendule de torsion reste aujourd'hui un instrument de précision pour mesurer des forces extrêmement faibles.
\end{savaistu}

\courssub{Énergie mécanique du pendule élastique}

Revenons au pendule élastique horizontal étudié dans les sections précédentes : un solide de masse $m$, accroché à un ressort de raideur $k$, glisse sans frottement sur un support horizontal. À un instant quelconque, son élongation est $x$ et sa vitesse $v = \dot{x}$.

\begin{definition}[Énergie mécanique d'un oscillateur élastique]
L'\cle{énergie mécanique} du système \{solide + ressort\} est la somme de son énergie cinétique et de son énergie potentielle élastique :
\[ E_m = E_c + E_{pe} = \frac{1}{2}m v^2 + \frac{1}{2}k x^2 \]
\end{definition}

\begin{propriete}[Conservation de l'énergie mécanique — démonstration exigible]
En l'absence de frottements, l'énergie mécanique d'un oscillateur harmonique est \cle{constante} et vaut :
\[ E_m = \frac{1}{2}k\,X_m^2 = \frac{1}{2}m\,\omega_0^2\,X_m^2 \]
\textbf{Démonstration par dérivation.} Calculons la dérivée de $E_m(t)$ :
\begin{align*}
\frac{dE_m}{dt} &= \frac{d}{dt}\left(\frac{1}{2}m\dot{x}^2 + \frac{1}{2}k x^2\right) \\
&= m\,\dot{x}\,\ddot{x} + k\,x\,\dot{x} \\
&= \dot{x}\,\left(m\ddot{x} + kx\right)
\end{align*}
Or l'équation différentielle du mouvement s'écrit précisément $m\ddot{x} + kx = 0$. Donc :
\[ \frac{dE_m}{dt} = 0 \quad\Longrightarrow\quad E_m = \text{constante} \]
Pour trouver sa valeur, évaluons-la à un instant particulier : au passage par un point extrême ($x = X_m$), la vitesse s'annule, donc $E_m = \tfrac{1}{2}kX_m^2$. Comme $\omega_0^2 = k/m$, on a aussi $E_m = \tfrac{1}{2}m\omega_0^2 X_m^2$. On pouvait aussi injecter directement $x(t) = X_m\cos(\omega_0 t + \varphi)$ et $v(t) = -X_m\omega_0\sin(\omega_0 t + \varphi)$ dans $E_m$ et utiliser $\cos^2 + \sin^2 = 1$.
\end{propriete}

\begin{attention}[Piège classique : $E_m = \tfrac{1}{2}kx^2$ ?]
L'erreur la plus fréquente est d'écrire $E_m = \tfrac{1}{2}k x^2$ avec l'élongation \emph{variable} $x$. C'est faux : $\tfrac{1}{2}kx^2$ n'est que l'énergie \emph{potentielle} à l'instant $t$ ! L'énergie mécanique totale, elle, vaut $E_m = \tfrac{1}{2}k\,\cle{X_m}^2$ avec l'\textbf{amplitude} $X_m$ (constante du mouvement). De même, $E_m = \tfrac{1}{2}m v_{\max}^2$ avec la vitesse \emph{maximale} $v_{\max} = \omega_0 X_m$, pas avec la vitesse instantanée. Retiens : l'énergie mécanique se calcule avec les valeurs \emph{extrémales} (amplitude ou vitesse maximale), jamais avec les valeurs instantanées seules.
\end{attention}

\begin{exemplebox}[Calcul d'une énergie mécanique]
Un solide de masse $m = \SI{200}{g}$ accroché à un ressort de raideur $k = \SI{50}{N.m^{-1}}$ oscille sans frottement avec une amplitude $X_m = \SI{5,0}{cm}$. Son énergie mécanique vaut :
\[ E_m = \frac{1}{2}kX_m^2 = \frac{1}{2}\times 50\times (0{,}050)^2 = \SI{6,3e-2}{J} \]
Vérification avec la vitesse maximale : $\omega_0 = \sqrt{k/m} = \sqrt{50/0{,}200} = \SI{16}{rad.s^{-1}}$, donc $v_{\max} = \omega_0 X_m = \SI{0,80}{m.s^{-1}}$ et $\tfrac{1}{2}mv_{\max}^2 = \tfrac{1}{2}\times 0{,}200\times (0{,}80)^2 = \SI{6,4e-2}{J}$. Les deux expressions donnent le même résultat à la précision des chiffres significatifs près (deux chiffres significatifs imposés par $k=50$ et $X_m=5,0$).
\end{exemplebox}

\courssub{Échanges énergétiques au cours d'une période}

Que se passe-t-il concrètement pendant une oscillation ? Prenons $x(t) = X_m\cos(\omega_0 t)$ (phase nulle pour simplifier). Alors :
\[ E_{pe}(t) = \frac{1}{2}kX_m^2\cos^2(\omega_0 t), \qquad E_c(t) = \frac{1}{2}kX_m^2\sin^2(\omega_0 t) \]

\begin{itemize}
\item À $t = 0$ : $x = X_m$, $v = 0$ : toute l'énergie est potentielle ($E_{pe} = E_m$, $E_c = 0$). Le ressort est au maximum de sa déformation.
\item À $t = T_0/4$ : $x = 0$, $v = v_{\max}$ : toute l'énergie est cinétique. Le solide traverse la position d'équilibre à vitesse maximale.
\item À $t = T_0/2$ : $x = -X_m$, $v = 0$ : l'énergie est de nouveau entièrement potentielle.
\item Et le cycle recommence.
\end{itemize}

Il y a donc une \cle{transformation continue} et périodique d'énergie potentielle élastique en énergie cinétique, et réciproquement, la somme restant constante. Remarque importante : comme $\cos^2$ et $\sin^2$ oscillent à la pulsation $2\omega_0$, les énergies $E_c$ et $E_{pe}$ ont une période égale à $T_0/2$.

\begin{popfigure}
\begin{center}
\begin{tikzpicture}
\begin{axis}[
width=0.95\linewidth, height=6.2cm,
xlabel={$t/T_0$}, ylabel={Énergie (J)},
xmin=0, xmax=1, ymin=0, ymax=0.07,
xtick={0,0.25,0.5,0.75,1},
xticklabels={$0$,$T_0/4$,$T_0/2$,$3T_0/4$,$T_0$},
legend style={at={(0.5,1.02)},anchor=south,legend columns=3,font=\small},
grid=both, grid style={popline!40},
]
\addplot[popcurve,domain=0:1,samples=200,color=popblue,very thick] {0.0625*cos(deg(2*pi*x))^2};
\addlegendentry{$E_{pe}(t)$}
\addplot[popcurve,domain=0:1,samples=200,color=poporange,very thick] {0.0625*sin(deg(2*pi*x))^2};
\addlegendentry{$E_c(t)$}
\addplot[popcurve,domain=0:1,samples=2,color=popgreen,very thick,dashed] {0.0625};
\addlegendentry{$E_m$ (constante)}
\end{axis}
\end{tikzpicture}
\end{center}
\legende{Échanges énergétiques sur une période (exemple de $k = \SI{50}{N.m^{-1}}$, $X_m = \SI{5,0}{cm}$) : $E_c$ et $E_{pe}$ s'échangent continuellement, leur somme $E_m = \SI{6,3e-2}{J}$ reste constante.}
\end{popfigure}

\begin{aretenir}[L'essentiel énergétique]
Pour tout oscillateur harmonique non amorti :
\[ E_m = E_c + E_p = \text{constante}, \qquad E_m = \frac{1}{2}kX_m^2 = \frac{1}{2}m v_{\max}^2 \]
L'énergie se « verse » alternativement du réservoir potentiel vers le réservoir cinétique, deux fois par période. C'est exactement l'analogue mécanique des échanges entre condensateur et bobine dans le circuit $LC$ que tu étudieras en électricité.
\end{aretenir}

\courssub{Énergie mécanique du pendule simple}

Pour le pendule simple (masse $m$ au bout d'un fil inextensible de longueur $\ell$), l'énergie cinétique est celle d'un mouvement circulaire : $E_c = \tfrac{1}{2}m v^2 = \tfrac{1}{2}m\ell^2\dot{\theta}^2$. L'énergie potentielle de pesanteur, en prenant l'origine à la position d'équilibre ($\theta = 0$), vaut $E_p = mg\ell(1 - \cos\theta)$, car la masse s'élève d'une hauteur $h = \ell(1-\cos\theta)$. D'où :

\begin{propriete}[Énergie mécanique du pendule simple]
En l'absence de frottements, l'énergie mécanique du pendule simple est conservée :
\[ E_m = \frac{1}{2}m\ell^2\dot{\theta}^2 + mg\ell(1 - \cos\theta) = \text{constante} \]
Si le pendule est lâché sans vitesse depuis l'angle $\theta_m$, alors $E_m = mg\ell(1-\cos\theta_m)$ et la \cle{vitesse maximale}, atteinte au passage à la position d'équilibre ($\theta = 0$), est :
\[ v_{\max} = \sqrt{2g\ell\,(1 - \cos\theta_m)} \]
\end{propriete}

\begin{exoresolu}[Le pendule du laboratoire de physique]
Un pendule simple est constitué d'une bille de masse $m = \SI{100}{g}$ suspendue à un fil de longueur $\ell = \SI{80}{cm}$. On l'écarte de $\theta_m = 30°$ de la verticale et on le lâche sans vitesse. On prend $g = \SI{9,8}{N.kg^{-1}}$.
\begin{enumerate}
\item Calculer l'énergie mécanique du pendule (origine des énergies potentielles à l'équilibre).
\item En déduire la vitesse maximale de la bille.
\item Calculer la vitesse de la bille lorsque $\theta = 15°$.
\end{enumerate}
\tcblower
\textbf{1. Énergie mécanique.} Au lâcher, la vitesse est nulle : toute l'énergie est potentielle.
\[ E_m = mg\ell(1-\cos\theta_m) = 0{,}100\times 9{,}8\times 0{,}80\times (1 - \cos 30°) \]
\[ \cos 30° = \frac{\sqrt{3}}{2} \approx 0{,}866 \quad\Rightarrow\quad 1-\cos 30° \approx 0{,}134 \]
\[ E_m = 0{,}784\times 0{,}134 \approx \SI{1,1e-1}{J} \quad\text{(2 chiffres significatifs)} \]
\textbf{2. Vitesse maximale.} À l'équilibre, $E_p = 0$, donc $E_m = \tfrac{1}{2}mv_{\max}^2$ :
\[ v_{\max} = \sqrt{\frac{2E_m}{m}} = \sqrt{\frac{2\times 0{,}11}{0{,}100}} = \sqrt{2{,}2} \approx \SI{1,5}{m.s^{-1}} \]
On retrouve le même résultat avec $v_{\max} = \sqrt{2g\ell(1-\cos\theta_m)}$.
\textbf{3. Vitesse à $\theta = 15°$.} Conservation de l'énergie : $\tfrac{1}{2}mv^2 = E_m - mg\ell(1-\cos 15°)$.
\[ \cos 15° \approx 0{,}966 \quad\Rightarrow\quad 1-\cos 15° \approx 0{,}034 \]
\[ mg\ell(1-\cos 15°) = 0{,}784\times 0{,}034 \approx \SI{2,7e-2}{J} \]
\[ v = \sqrt{\frac{2\times(0{,}11 - 0{,}027)}{0{,}100}} = \sqrt{1{,}66} \approx \SI{1,3}{m.s^{-1}} \]
La bille a déjà atteint une grande partie de sa vitesse maximale : l'énergie potentielle diminue vite dès que le pendule descend.
\end{exoresolu}

\begin{attention}[Conditions d'application de la conservation de l'énergie]
La conservation de $E_m$ suppose l'\textbf{absence totale de frottements}. Dès qu'il existe des frottements (air, axe, articulation), l'énergie mécanique \emph{diminue} au cours du temps : c'est le phénomène d'\cle{amortissement}, que nous étudierons dans la prochaine section. Autre piège : pour le pendule simple, n'oublie pas que $E_p = mg\ell(1-\cos\theta)$ n'est \emph{pas} proportionnelle à $\theta^2$ en général ; l'approximation harmonique $E_p \approx \tfrac{1}{2}mg\ell\,\theta^2$ (car $1-\cos\theta \approx \theta^2/2$) n'est valable que pour les petits angles (typiquement $\theta < 10°$).
\end{attention}

\begin{exercice}[Pendule de torsion d'un laboratoire yaoundéen]
Une tige homogène de moment d'inertie $J_\Delta = \SI{4,0e-3}{kg.m^2}$ est suspendue à un fil de torsion. On mesure une période d'oscillation $T_0 = \SI{2,5}{s}$.
\begin{enumerate}
\item Déterminer la constante de torsion $C$ du fil.
\item On écarte la tige de $\theta_m = 0{,}20$ rad et on la lâche sans vitesse. Calculer l'énergie mécanique du pendule et sa vitesse angulaire maximale.
\end{enumerate}
\end{exercice}

\begin{corrige}[Exercice 1]
\textbf{1.} De $T_0 = 2\pi\sqrt{J_\Delta/C}$ on tire $C = \dfrac{4\pi^2 J_\Delta}{T_0^2} = \dfrac{4\pi^2\times 4{,}0\times 10^{-3}}{(2{,}5)^2} \approx \SI{2,5e-2}{N.m.rad^{-1}}$ (2 chiffres significatifs).
\textbf{2.} Par analogie translation/rotation, $E_m = \tfrac{1}{2}C\theta_m^2 = \tfrac{1}{2}\times 2{,}5\times 10^{-2}\times (0{,}20)^2 = \SI{5,0e-4}{J}$. La vitesse angulaire maximale vérifie $\tfrac{1}{2}J_\Delta\dot{\theta}_{\max}^2 = E_m$, d'où $\dot{\theta}_{\max} = \sqrt{2E_m/J_\Delta} = \sqrt{\dfrac{2\times 5{,}0\times 10^{-4}}{4{,}0\times 10^{-3}}} = \SI{0,50}{rad.s^{-1}}$. On vérifie : $\dot{\theta}_{\max} = \omega_0\theta_m = \dfrac{2\pi}{2{,}5}\times 0{,}20 \approx \SI{0,50}{rad.s^{-1}}$.
\end{corrige}

\coursec{Oscillations amorties}

Dans la section précédente, nous avons étudié des oscillateurs \emph{idéaux} : pendule élastique, pendule simple, pendule de torsion. Pour chacun, l'énergie mécanique se conservait rigoureusement et les oscillations se poursuivaient indéfiniment avec une amplitude constante. Mais l'expérience quotidienne nous dit autre chose : une balançoire que l'on abandonne finit par s'immobiliser, le pendule d'une horloge de cuisine s'arrête si l'on ne remonte pas les poids, la lame d'une règle que l'on fait vibrer sur le bord d'une table s'apaise en quelques secondes. \textbf{Tout oscillateur réel voit l'amplitude de ses oscillations diminuer au cours du temps.} C'est cette réalité physique que nous allons maintenant étudier.

\courssub{Origine de l'amortissement : les forces de frottement}

\courssub{Observation et mise en évidence expérimentale}

\begin{experience}[Un pendule élastique dans l'air puis dans l'eau]
\textbf{Matériel :} un ressort de raideur $k$, une masse $m$ marquée, une cuve d'eau, une caméra ou un logiciel de pointage (type Aviméca ou Regressi, disponibles dans les salles multimédias des lycées).

\textbf{Protocole :}
\begin{enumerate}
\item Suspendre la masse au ressort, l'écarter de sa position d'équilibre et la lâcher sans vitesse initiale. Enregistrer le mouvement dans l'\textbf{air}.
\item Recommencer en immergeant la masse dans la \textbf{cuve d'eau}.
\item Pointer les positions et tracer $x(t)$ dans les deux cas.
\end{enumerate}

\textbf{Observations :} dans l'air, les oscillations persistent longtemps mais leur amplitude diminue lentement ; dans l'eau, l'amplitude décroît très vite, et si le liquide est très visqueux (huile, glycérine), la masse \emph{ne dépasse même plus} la position d'équilibre : elle y revient directement, sans osciller.

\textbf{Interprétation :} la masse subit, en plus de la force de rappel du ressort, une \cle{force de frottement} exercée par le fluide (air ou eau). Cette force, toujours opposée au mouvement, \emph{dissipe} l'énergie mécanique de l'oscillateur : celle-ci se transforme en énergie thermique (échauffement du fluide et du ressort). L'oscillateur n'est plus un système conservatif.
\end{experience}

\courssub{Les différentes sources d'amortissement}

L'amortissement d'un oscillateur mécanique réel provient de plusieurs causes, souvent simultanées :

\begin{itemize}
\item les \cle{frottements fluides} : résistance de l'air ou d'un liquide sur le mobile en mouvement ;
\item les \cle{frottements solides} : contacts au niveau des axes de rotation (pivot du pendule pesant), glissement de pièces les unes sur les autres ;
\item les \cle{déformations internes} des pièces élastiques : un ressort réel s'échauffe légèrement à chaque cycle, ce qui prélève de l'énergie mécanique.
\end{itemize}

Dans tous les cas, le point commun est le suivant : \textbf{ces forces sont non conservatives et leur travail est négatif}, de sorte que l'énergie mécanique $E_m = E_c + E_p$ de l'oscillateur diminue continûment au profit de l'énergie thermique du milieu extérieur.

\begin{definition}[Oscillations amorties]
On dit qu'un oscillateur mécanique effectue des \cle{oscillations amorties} lorsque l'amplitude de ses oscillations diminue au cours du temps, par suite de la dissipation de son énergie mécanique sous l'effet de forces de frottement.
\end{definition}

\courssub{Modélisation de la force de frottement fluide}

Pour un mobile de vitesse modérée se déplaçant dans un fluide (gaz ou liquide), l'expérience montre que la force de frottement est, en première approximation, \textbf{proportionnelle à la vitesse et de sens opposé} :

\[
\vec{f} = -\lambda\,\vec{v}
\]

où $\lambda$ est le \cle{coefficient de frottement fluide}, positif, qui dépend de la nature du fluide et de la géométrie du mobile. Son unité se déduit de la relation $f = \lambda v$ :

\[
[\lambda] = \frac{[f]}{[v]} = \frac{\mathrm{kg\cdot m\cdot s^{-2}}}{\mathrm{m\cdot s^{-1}}} = \mathrm{kg\cdot s^{-1}}.
\]

Le coefficient $\lambda$ s'exprime donc en $\mathrm{kg\cdot s^{-1}}$ (ou, ce qui est équivalent, en $\mathrm{N\cdot s\cdot m^{-1}}$).

\begin{attention}[Conditions de validité du modèle]
Le modèle $\vec{f} = -\lambda\,\vec{v}$ n'est valable que pour des vitesses \emph{faibles} et un écoulement \emph{laminaire} du fluide autour du mobile. À grande vitesse (chute d'un parachutiste, voiture lancée), la force de frottement devient proportionnelle à $v^2$ et le modèle linéaire ne s'applique plus. Au programme de Terminale, seul le modèle linéaire est utilisé, et sa résolution complète n'est \textbf{pas exigible}.
\end{attention}

\courssub{Équation différentielle du mouvement (présentation qualitative)}

Reprenons le pendule élastique horizontal de la section 1 : une masse $m$ glisse sur un support horizontal, accrochée à un ressort de raideur $k$, et subit de plus la force de frottement fluide $\vec{f} = -\lambda\,\vec{v}$. Le bilan des forces s'écrit :

\begin{itemize}
\item le poids $\vec{P}$ et la réaction $\vec{R}$ du support, qui se compensent (support horizontal sans frottement solide) ;
\item la force de rappel du ressort : $\vec{F} = -k\,x\,\vec{i}$ ;
\item la force de frottement fluide : $\vec{f} = -\lambda\,\dot{x}\,\vec{i}$.
\end{itemize}

La deuxième loi de Newton, projetée sur l'axe $(Ox)$, donne :

\[
m\,\ddot{x} = -k\,x - \lambda\,\dot{x}
\qquad\Longleftrightarrow\qquad
\boxed{\;\ddot{x} + \frac{\lambda}{m}\,\dot{x} + \frac{k}{m}\,x = 0\;}
\]

On reconnaît le terme $\dfrac{k}{m} = \omega_0^2$ de l'oscillateur non amorti, auquel s'ajoute le terme d'amortissement $\dfrac{\lambda}{m}\,\dot{x}$. C'est une équation différentielle linéaire du second ordre \emph{avec terme du premier ordre}. Sa résolution est hors programme, mais la nature de ses solutions dépend de l'intensité de l'amortissement : c'est ce qui définit les différents \textbf{régimes} que nous allons décrire.

\begin{popfigure}
\begin{center}
\begin{tikzpicture}[scale=0.95]
% Support hachuré
\fill[pattern=north east lines, pattern color=popmuted] (-0.4,1.6) rectangle (0,3.2);
\draw[popdark, thick] (0,1.6) -- (0,3.2);
% Ressort (zigzag)
\draw[popblue, thick] (0,2.4)
  -- (0.3,2.4)
  \foreach \i in {0,...,5} { -- ({0.45+0.25*\i},{2.4+0.18*(mod(\i,2)==0?1:-1)}) }
  -- (2.1,2.4);
% Masse
\fill[poporange!85] (2.1,1.9) rectangle (3.3,2.9);
\node[white, font=\bfseries] at (2.7,2.4) {$m$};
% Sol
\draw[popdark, thick] (-0.4,1.6) -- (5.6,1.6);
\fill[pattern=north west lines, pattern color=popmuted] (-0.4,1.3) rectangle (5.6,1.6);
% Axe x
\draw[-{Stealth}, popdark] (1.2,0.8) -- (4.6,0.8) node[below] {$x$};
\draw[popdark] (2.7,0.7) -- (2.7,0.9) node[above=6pt] {$O$ (équilibre)};
% Vecteurs forces
\draw[-{Stealth}, very thick, popgreen] (2.7,2.4) -- (1.7,2.4) node[above left, popgreen] {$\vec{F}$ (rappel)};
\draw[-{Stealth}, very thick, poppink] (2.7,2.15) -- (3.7,2.15) node[above right, poppink] {$\vec{f}=-\lambda\vec{v}$};
\draw[-{Stealth}, very thick, poppurple] (3.3,1.35) -- (4.3,1.35) node[below, poppurple] {$\vec{v}$};
\end{tikzpicture}
\end{center}
\legende{Pendule élastique horizontal amorti : la masse, animée d'une vitesse $\vec{v}$, subit la force de rappel $\vec{F}$ du ressort et la force de frottement fluide $\vec{f}$, opposée à $\vec{v}$.}
\end{popfigure}

\courssub{Les régimes d'oscillations amorties}

\courssub{Le régime pseudo-périodique (amortissement faible)}

Lorsque le coefficient de frottement $\lambda$ est \textbf{faible} (masse oscillant dans l'air, par exemple), la masse effectue encore des oscillations autour de sa position d'équilibre, mais l'amplitude de celles-ci \textbf{décroît au cours du temps}, approximativement selon une loi exponentielle. L'élongation peut se mettre sous la forme :

\[
x(t) = X_m\,e^{-\alpha t}\cos(\omega\,t + \varphi)
\]

où $\alpha$ est un coefficient positif lié à l'amortissement. Le facteur $e^{-\alpha t}$ joue le rôle d'une \emph{amplitude décroissante} : la courbe $x(t)$ reste comprise entre les deux courbes \textbf{enveloppes} $+X_m e^{-\alpha t}$ et $-X_m e^{-\alpha t}$.

On appelle \cle{pseudo-période} $T$ la durée qui sépare deux passages consécutifs du mobile par sa position d'équilibre \emph{dans le même sens} (ou, ce qui est équivalent, la durée entre deux maxima consécutifs de $x(t)$).

\begin{propriete}[Pseudo-période et période propre — admise]
Pour un amortissement \textbf{faible}, la pseudo-période $T$ des oscillations amorties est très voisine de la période propre $T_0$ de l'oscillateur non amorti :
\[
T \approx T_0 = 2\pi\sqrt{\frac{m}{k}}.
\]
Plus l'amortissement est faible, plus cette approximation est bonne. (Démonstration non exigible au Baccalauréat.)
\end{propriete}

\courssub{Le régime apériodique (amortissement fort)}

Si l'on augmente considérablement le frottement (masse plongée dans de la glycérine, par exemple), les oscillations \textbf{disparaissent} : écartée de sa position d'équilibre puis lâchée, la masse y revient \emph{lentement et sans la dépasser}. C'est le \cle{régime apériodique}. Il n'y a plus ni période, ni oscillation : la courbe $x(t)$ décroît de façon monotone vers zéro.

\courssub{Le régime critique}

Entre les deux régimes précédents existe un cas limite, appelé \cle{régime critique}, obtenu pour une valeur particulière de l'amortissement ($\lambda_c = 2\sqrt{km}$, valeur non exigible). Le mobile revient alors à sa position d'équilibre \textbf{sans osciller, et plus rapidement que dans tout régime apériodique}. Ce régime est recherché dans de nombreuses applications pratiques : c'est le réglage idéal des amortisseurs de véhicules.

\begin{popfigure}
\begin{center}
\begin{tikzpicture}
\begin{axis}[
  width=\linewidth, height=7.2cm,
  xlabel={$t$ (s)}, ylabel={$x$ (cm)},
  xmin=0, xmax=10, ymin=-1.3, ymax=1.3,
  grid=both, grid style={popline!40},
  axis line style={popdark},
  legend style={at={(0.98,0.35)}, anchor=east, font=\small, draw=popline},
  samples=300
]
% Enveloppes exponentielles
\addplot[popmuted, dashed, thick, domain=0:10] {exp(-0.25*x)};
\addplot[popmuted, dashed, thick, domain=0:10] {-exp(-0.25*x)};
% Régime pseudo-périodique
\addplot[popblue, very thick, domain=0:10] {exp(-0.25*x)*cos(deg(2*pi*x/2))};
% Régime apériodique
\addplot[poporange, very thick, domain=0:10] {0.5*(exp(-0.8*x)+exp(-2.2*x))};
\legend{enveloppe $\pm X_m e^{-\alpha t}$, , pseudo-périodique, apériodique}
\end{axis}
\end{tikzpicture}
\end{center}
\legende{Comparaison des régimes. En bleu : régime pseudo-périodique, oscillations d'amplitude décroissante comprises entre les enveloppes exponentielles (pointillés). En orange : régime apériodique, retour à l'équilibre sans oscillation.}
\end{popfigure}

\begin{attention}[Erreurs fréquentes à éviter]
\begin{itemize}
\item \textbf{Ne pas confondre pseudo-période et période propre.} La pseudo-période $T$ est mesurée sur l'oscillateur \emph{amorti} ; la période propre $T_0 = 2\pi\sqrt{m/k}$ est une caractéristique de l'oscillateur \emph{idéal}. Elles ne sont voisines que si l'amortissement est faible.
\item \textbf{Ce n'est pas la fréquence qui diminue, c'est l'amplitude !} Dans le régime pseudo-périodique, les oscillations se succèdent à intervalles de temps quasi constants ($T \approx T_0$) ; ce qui décroît, c'est l'écart maximal à l'équilibre. Dire que « l'oscillateur ralentit » est une erreur classique de rédaction.
\item La position d'équilibre \textbf{ne change pas} avec l'amortissement : pour le pendule élastique vertical, elle reste $x_{\text{éq}} = \dfrac{mg}{k}$.
\end{itemize}
\end{attention}

\courssub{Décroissance de l'énergie mécanique}

Pour l'oscillateur non amorti, nous avons établi que l'énergie mécanique se conserve : $E_m = \dfrac{1}{2}kX_m^2 = \text{constante}$. Qu'en est-il avec frottements ?

Multiplions l'équation différentielle par $\dot{x}$ :

\[
m\,\ddot{x}\,\dot{x} + k\,x\,\dot{x} = -\lambda\,\dot{x}^2
\]

Or $\ddot{x}\,\dot{x} = \dfrac{d}{dt}\!\left(\dfrac{\dot{x}^2}{2}\right)$ et $x\,\dot{x} = \dfrac{d}{dt}\!\left(\dfrac{x^2}{2}\right)$. On reconnaît donc :

\[
\frac{d}{dt}\underbrace{\left(\frac{1}{2}m\dot{x}^2 + \frac{1}{2}kx^2\right)}_{E_m} = -\lambda\,\dot{x}^2 \leq 0
\]

\begin{propriete}[Décroissance de l'énergie mécanique]
Pour un oscillateur soumis à une force de frottement fluide $\vec{f} = -\lambda\vec{v}$, l'énergie mécanique vérifie :
\[
\frac{dE_m}{dt} = -\lambda\,v^2 = \vec{f}\cdot\vec{v} \leq 0.
\]
L'énergie mécanique \textbf{diminue continûment} : la puissance des frottements est toujours négative. L'énergie perdue est convertie en énergie thermique (échauffement) dans le fluide et le ressort. L'énergie totale du système complet (oscillateur + milieu) reste, elle, conservée.
\end{propriete}

Comme $E_m \approx \dfrac{1}{2}kX_m(t)^2$ à chaque instant (en régime pseudo-périodique), la décroissance de l'amplitude entraîne une décroissance de l'énergie : si l'amplitude est divisée par $2$, l'énergie mécanique est divisée par $4$ (elle varie comme le \emph{carré} de l'amplitude).

\courssub{Exploitation d'un enregistrement d'oscillations amorties}

L'exploitation d'un enregistrement expérimental de $x(t)$ est une compétence classique au Baccalauréat. Voici la démarche à maîtriser.

\begin{methode}[Exploiter un enregistrement d'oscillations amorties]
\begin{enumerate}
\item \textbf{Identifier le régime} : oscillations d'amplitude décroissante $\Rightarrow$ régime pseudo-périodique ; retour à l'équilibre sans oscillation $\Rightarrow$ régime apériodique.
\item \textbf{Mesurer la pseudo-période avec précision} : ne pas mesurer une seule pseudo-période, mais repérer la durée $\Delta t$ de $n$ pseudo-périodes consécutives (par exemple $n = 5$ ou $10$), puis calculer $T = \dfrac{\Delta t}{n}$. L'incertitude de lecture est ainsi divisée par $n$.
\item \textbf{Relever les maxima successifs} $X_1, X_2, X_3, \dots$ (amplitudes aux passages par les sommets, du même côté de l'équilibre).
\item \textbf{Vérifier la décroissance exponentielle} : calculer les rapports $\dfrac{X_{n+1}}{X_n}$. S'ils sont (approximativement) constants, l'amplitude décroît bien de façon exponentielle : $X_n = X_1\,r^{\,n-1}$ avec $r < 1$.
\item \textbf{Conclure} : comparer $T$ à la période propre calculée $T_0 = 2\pi\sqrt{m/k}$ et vérifier que $T \approx T_0$ si l'amortissement est faible.
\end{enumerate}
\end{methode}

\begin{exoresolu}[Étude d'un pendule élastique amorti]
Un solide de masse $m = \SI{200}{g}$ est accroché à un ressort horizontal de raideur $k = \SI{20}{N.m^{-1}}$. On enregistre son mouvement amorti. Sur l'enregistrement, on lit :
\begin{itemize}
\item la durée de $5$ pseudo-périodes : $\Delta t = \SI{3,15}{s}$ ;
\item les amplitudes successives : $X_1 = \SI{8,0}{cm}$, $X_2 = \SI{6,4}{cm}$, $X_3 = \SI{5,1}{cm}$, $X_4 = \SI{4,1}{cm}$.
\end{itemize}
\textbf{1.} Quel est le régime d'oscillation ? \textbf{2.} Déterminer la pseudo-période $T$ et la comparer à la période propre $T_0$. \textbf{3.} Montrer que l'amplitude décroît de façon exponentielle. \textbf{4.} Calculer l'énergie mécanique initiale et l'énergie mécanique lors du 4\textsuperscript{e} maximum. Conclure.
\tcblower
\textbf{1.} Le solide oscille avec une amplitude qui diminue : c'est le \textbf{régime pseudo-périodique} (amortissement faible).

\textbf{2.} Pseudo-période :
\[
T = \frac{\Delta t}{n} = \frac{\num{3,15}}{5} = \SI{0,63}{s}.
\]
Période propre :
\[
T_0 = 2\pi\sqrt{\frac{m}{k}} = 2\pi\sqrt{\frac{\num{0,200}}{20}} = 2\pi\sqrt{\num{0,010}} \approx \SI{0,63}{s}.
\]
On constate que $T \approx T_0$ : l'amortissement est bien faible, la propriété est vérifiée.

\textbf{3.} Calculons les rapports d'amplitudes consécutives :
\[
\frac{X_2}{X_1} = \frac{\num{6,4}}{\num{8,0}} = \num{0,80} ; \qquad
\frac{X_3}{X_2} = \frac{\num{5,1}}{\num{6,4}} \approx \num{0,80} ; \qquad
\frac{X_4}{X_3} = \frac{\num{4,1}}{\num{5,1}} \approx \num{0,80}.
\]
Le rapport est constant ($r \approx \num{0,80}$) : l'amplitude est multipliée par $\num{0,80}$ à chaque pseudo-période, ce qui est la signature d'une \textbf{décroissance exponentielle} $X_n = X_1 \times \num{0,80}^{\,n-1}$.

\textbf{4.} Énergie mécanique initiale (au premier maximum, toute l'énergie est potentielle élastique) :
\[
E_{m,1} = \frac{1}{2}kX_1^2 = \frac{1}{2}\times 20 \times (\num{0,080})^2 = \SI{6,4e-2}{J} = \SI{64}{mJ}.
\]
Au quatrième maximum :
\[
E_{m,4} = \frac{1}{2}kX_4^2 = \frac{1}{2}\times 20 \times (\num{0,041})^2 \approx \SI{1,7e-2}{J} = \SI{17}{mJ}.
\]
\textbf{Conclusion :} l'énergie mécanique a diminué d'environ $\SI{47}{mJ}$ en trois pseudo-périodes ; cette énergie a été dissipée sous forme de chaleur par les frottements. On remarque que $E_{m,4}/E_{m,1} \approx \num{0,26} \approx \num{0,80}^{6}$ : l'énergie décroît comme le \emph{carré} de l'amplitude.
\end{exoresolu}

\courssub{Applications concrètes de l'amortissement}

Loin d'être toujours une gêne, l'amortissement est souvent \emph{recherché} et soigneusement réglé par les ingénieurs :

\begin{itemize}
\item \textbf{Amortisseurs de voitures et de motos} : après le passage sur une bosse (nombreuses sur nos routes, de Douala à Maroua !), la roue doit revenir à sa position \emph{le plus vite possible sans osciller} : le réglage idéal est proche du \textbf{régime critique}. Un amortisseur usé (amortissement trop faible) laisse le véhicule rebondir plusieurs fois : danger en virage et au freinage.
\item \textbf{Portes battantes et fermetures automatiques} : un vérin à frottement fluide empêche la porte de claquer ; elle revient en position fermée de façon apériodique, doucement.
\item \textbf{Suspension des balances de précision} : le fléau est amorti (par frottement dans l'air ou par amortisseur magnétique) pour que l'aiguille se stabilise rapidement sur la graduation d'équilibre, au lieu d'osciller longuement.
\item \textbf{Sièges, ponts, bâtiments} : toute structure susceptible de vibrer est équipée de dispositifs dissipant l'énergie des oscillations.
\end{itemize}

\begin{savaistu}[Les amortisseurs des motos-taxis]
Au Cameroun, les motos-taxis (les « benskin » de Douala ou de Yaoundé) transportent chaque jour des millions de passagers sur des chaussées souvent dégradées. Leurs amortisseurs arrière, fortement sollicités, sont conçus pour fonctionner au voisinage du \textbf{régime critique} : assez amortis pour éviter les rebonds successifs après un nid-de-poule (sécurité), mais pas trop raides pour ne pas transmettre brutalement les chocs au passager (confort). Un amortisseur « mort » se reconnaît facilement : la moto rebondit deux ou trois fois après une bosse — c'est le régime pseudo-périodique malvenu ! Les constructeurs testent ce réglage en écrasant brutalement la suspension : une bonne moto revient à l'équilibre en un seul mouvement, sans oscillation.
\end{savaistu}

\begin{aretenir}[L'essentiel sur les oscillations amorties]
\begin{itemize}
\item Les frottements (fluide : $\vec{f} = -\lambda\vec{v}$, ou solide) dissipent l'énergie mécanique : $\dfrac{dE_m}{dt} = -\lambda v^2 < 0$ ; l'amplitude décroît.
\item \textbf{Amortissement faible} : régime pseudo-périodique, amplitude en décroissance exponentielle, pseudo-période $T \approx T_0$.
\item \textbf{Amortissement fort} : régime apériodique, retour à l'équilibre sans oscillation.
\item \textbf{Régime critique} : retour à l'équilibre le plus rapide possible sans oscillation ; réglage idéal des suspensions.
\item Exploitation d'un enregistrement : mesurer $n$ pseudo-périodes ($T = \Delta t/n$), relever les maxima successifs et vérifier que leur rapport est constant.
\end{itemize}
\end{aretenir}

\begin{exercice}[Pour s'entraîner]
Un pendule simple de longueur $\ell = \SI{1,0}{m}$ est lâché avec un écart angulaire $\theta_1 = \SI{10}{\degree}$. Après $10$ oscillations, l'écart maximal n'est plus que $\theta_{11} = \SI{5,0}{\degree}$. On prendra $g = \SI{9,8}{m.s^{-2}}$.
\begin{enumerate}
\item Calculer la période propre $T_0$ du pendule. Peut-on l'identifier à la pseudo-période ?
\item En supposant la décroissance exponentielle, calculer le rapport $r$ par lequel l'amplitude est multipliée à chaque pseudo-période.
\item Au bout de combien d'oscillations l'amplitude sera-t-elle inférieure à $\SI{1}{\degree}$ ?
\item Que devient l'énergie mécanique perdue par le pendule ?
\end{enumerate}
\end{exercice}

\begin{corrige}[Exercice — Pour s'entraîner]
\textbf{1.} $T_0 = 2\pi\sqrt{\ell/g} = 2\pi\sqrt{\frac{\num{1,0}}{\num{9,8}}} \approx \SI{2,0}{s}$. L'amortissement est faible (l'amplitude n'est divisée que par $2$ en $10$ périodes), donc $T \approx T_0 \approx \SI{2,0}{s}$.

\textbf{2.} On a $\theta_{11} = \theta_1\,r^{10}$, d'où $r^{10} = \num{0,50}$ et $r = (\num{0,50})^{1/10} \approx \num{0,93}$. L'amplitude perd environ $7\,\%$ par oscillation.

\textbf{3.} On cherche $n$ tel que $10 \times \num{0,93}^{\,n} < 1$, soit $\num{0,93}^{\,n} < \num{0,1}$. En testant : $\num{0,93}^{30} \approx \num{0,11}$ et $\num{0,93}^{32} \approx \num{0,098}$. Il faut donc environ $32$ oscillations, soit une durée $t \approx 32 \times \num{2,0} = \SI{64}{s}$.

\textbf{4.} L'énergie mécanique perdue est convertie en énergie thermique : échauffement de l'air brassé par la masse et de l'axe de suspension (frottements au pivot).
\end{corrige}

\coursec{Oscillations forcées et résonance}

Dans la section précédente, nous avons vu que tout oscillateur réel, livré à lui-même, voit ses oscillations s'amortir : les frottements dissipent son énergie mécanique et l'amplitude décroît jusqu'à l'arrêt. Comment alors entretenir un mouvement oscillatoire ? Comment expliquer qu'une balançoire continue de monter si haut, qu'une carrosserie de voiture se mette à vibrer violemment à certaine vitesse, ou qu'un pont puisse osciller dangereusement sous le vent ? La réponse tient en une idée : il faut \emph{apporter de l'énergie de l'extérieur}, au bon rythme. C'est le monde des \cle{oscillations forcées} et du phénomène spectaculaire de \cle{résonance}.

\courssub{Excitateur et résonateur : qui commande, qui obéit ?}

\textbf{Une expérience pour commencer.}\par

Suspendons une masse $m$ à un ressort de raideur $k$, et fixons l'extrémité supérieure du ressort non plus à un support immobile, mais à la tige d'un moteur vibrant (vibreur électromagnétique) dont on peut régler la fréquence $f$. Le vibreur impose à l'extrémité du ressort un mouvement vertical sinusoïdal de fréquence $f$.

Observations :
\begin{itemize}
\item si le vibreur oscille \emph{lentement} ($f$ très petite), la masse suit tranquillement le mouvement, avec une faible amplitude ;
\item si le vibreur oscille \emph{rapidement} ($f$ grande), la masse bouge à peine : elle n'a « pas le temps » de suivre ;
\item mais pour une fréquence \emph{particulière}, proche de la fréquence propre $f_0$ du système masse--ressort, la masse se met à osciller avec une \emph{très grande amplitude}, parfois si grande que le ressort bute ou que la masse heurte le support !
\end{itemize}

\begin{definition}[Excitateur et résonateur]
Dans un dispositif d'oscillations forcées :
\begin{itemize}
\item l'\cle{excitateur} est le système qui \emph{impose} une action périodique (force, déplacement, tension) de fréquence $f$ réglable ;
\item le \cle{résonateur} est l'oscillateur qui \emph{subit} cette action et se met à osciller sous son effet.
\end{itemize}
Exemples : le moteur vibrant fixé sur une structure (excitateur) et la structure elle-même (résonateur) ; le vent (excitateur) et un pont suspendu (résonateur) ; les ondulations d'une piste (excitateur) et la suspension d'une voiture (résonateur).
\end{definition}

\begin{savaistu}[Le diapason et la caisse de résonance]
Quand on frappe un diapason, ses branches vibrent à leur fréquence propre (par exemple $f_0 = \SI{440}{Hz}$, le \emph{la} des musiciens). Posé sur sa caisse de résonance, l'air qu'elle contient est excité par les vibrations du diapason : si la fréquence propre de la colonne d'air est voisine de $\SI{440}{Hz}$, elle entre en résonance et le son devient nettement plus fort. C'est le même principe qui donne leur puissance sonore à la guitare, au balafon ou au tam-tam : la caisse ou la calebasse joue le rôle de résonateur.
\end{savaistu}

\courssub{Oscillations forcées : mise en équation}

Considérons le pendule élastique horizontal étudié dans les sections précédentes : masse $m$, ressort de raideur $k$, frottement fluide $\vec{f} = -\lambda\,\vec{v}$. L'excitateur exerce sur la masse une force supplémentaire \emph{sinusoïdale} de pulsation $\omega = 2\pi f$ :
\[ F_{\text{exc}}(t) = F_m\cos(\omega t). \]

Le bilan des forces appliquées à la masse, projeté sur l'axe $(Ox)$, donne :
\begin{itemize}
\item force de rappel du ressort : $-kx$ ;
\item force de frottement fluide : $-\lambda\,\dot{x}$ ;
\item force excitatrice : $F_m\cos(\omega t)$.
\end{itemize}

La deuxième loi de Newton ($\sum \vec{F} = m\vec{a}$) conduit à l'équation différentielle :
\[ m\ddot{x} = -kx - \lambda\dot{x} + F_m\cos(\omega t), \]
soit, en divisant par $m$ et en posant $\omega_0^2 = \dfrac{k}{m}$ :
\[ \ddot{x} + \frac{\lambda}{m}\,\dot{x} + \omega_0^2\,x = \frac{F_m}{m}\cos(\omega t). \]

\begin{attention}[Une équation à ne pas confondre]
Cette équation est \emph{différente} de l'équation de l'oscillateur libre amorti : le second membre n'est pas nul, il contient le terme excitateur $\dfrac{F_m}{m}\cos(\omega t)$. Au programme de Terminale C, on ne demande pas de résoudre cette équation : on en exploite qualitativement les conséquences, qui suivent.
\end{attention}

La solution générale est la somme de deux termes :
\begin{enumerate}
\item un \emph{régime transitoire}, de même nature que l'oscillation libre amortie (pseudo-périodique ou apériodique), qui \emph{disparaît} au bout de quelques périodes à cause de l'amortissement ;
\item un \emph{régime permanent}, solution sinusoïdale de pulsation $\omega$ :
\[ x(t) = X_m(\omega)\cos(\omega t + \varphi). \]
\end{enumerate}

\begin{propriete}[Régime permanent des oscillations forcées (admise)]
Après disparition du régime transitoire, l'oscillateur forcé vibre de façon sinusoïdale :
\begin{itemize}
\item à la \cle{fréquence $f$ de l'excitateur} (pulsation $\omega = 2\pi f$), et \emph{non} à sa fréquence propre $f_0$ ;
\item avec une amplitude $X_m$ qui \emph{dépend de la fréquence $f$} et de l'amortissement.
\end{itemize}
\end{propriete}

\begin{attention}[Erreur fréquente au Bac]
Beaucoup d'élèves écrivent qu'en régime forcé, l'oscillateur « vibre à sa fréquence propre ». C'est \textbf{faux} : en régime permanent, l'oscillateur vibre à la \cle{fréquence de l'excitateur} $f$, imposée de l'extérieur. La fréquence propre $f_0$ n'intervient que pour déterminer \emph{l'amplitude} de la réponse : c'est quand $f$ s'approche de $f_0$ que l'amplitude devient maximale. Retiens : l'excitateur impose le \emph{rythme}, la nature de l'oscillateur fixe l'\emph{ampleur} de la réponse.
\end{attention}

\courssub{La résonance mécanique}

\begin{definition}[Résonance mécanique]
On dit qu'un oscillateur entre en \cle{résonance mécanique} lorsque la fréquence $f$ de l'excitateur devient voisine de sa fréquence propre $f_0$ : l'amplitude des oscillations forcées passe alors par un \emph{maximum}, qui peut être très grand si l'amortissement est faible.
\end{definition}

\begin{propriete}[Propriétés de la résonance (admises)]
\begin{itemize}
\item À la résonance, la fréquence imposée par l'excitateur est voisine de la fréquence propre : $f \approx f_0$ (avec un amortissement faible, la fréquence de résonance est pratiquement égale à $f_0$).
\item L'amplitude maximale $X_{m,\max}$ dépend de l'amortissement : \emph{plus l'amortissement est faible, plus l'amplitude à la résonance est grande}.
\item On distingue :
\begin{itemize}
\item la \emph{résonance aiguë} (amortissement faible) : pic étroit et très élevé, potentiellement dangereux ;
\item la \emph{résonance floue} (amortissement fort) : pic large et peu marqué.
\end{itemize}
\end{itemize}
\end{propriete}

\textbf{Interprétation énergétique.}\par

L'excitateur fournit à chaque période un travail à l'oscillateur ; les frottements dissipent de l'énergie. En régime permanent, il y a équilibre : l'énergie reçue compense exactement l'énergie dissipée. Quand $f \approx f_0$, la force excitatrice agit « au bon moment », en phase avec le mouvement, comme quand on pousse une balançoire exactement quand elle repart vers l'avant : le transfert d'énergie est maximal, et l'amplitude croît jusqu'à ce que la dissipation (qui augmente avec l'amplitude) équilibre l'apport. Si l'amortissement est très faible, il faut une amplitude énorme pour dissiper autant que l'excitateur fournit : d'où le danger.

\begin{popfigure}
\begin{center}
\begin{tikzpicture}
\begin{axis}[
    width=0.95\linewidth, height=7cm,
    xlabel={Fréquence de l'excitateur $f$ (Hz)},
    ylabel={Amplitude $X_m$ (cm)},
    xmin=0, xmax=3, ymin=0, ymax=11,
    xtick={0,0.5,1,1.5,2,2.5,3},
    grid=major, grid style={popline!40},
    legend style={at={(0.97,0.97)}, anchor=north east, font=\small},
    domain=0.05:3, samples=300
]
% résonance aiguë : amortissement faible, f0 = 1.5 Hz
\addplot[popcurve, popblue, very thick] {1.5/sqrt(max(0,(1- (x/1.5)^2)^2 + 0.02*(x/1.5)^2))};
% résonance floue : amortissement fort
\addplot[popcurve, poporange, very thick] {1.5/sqrt(max(0,(1- (x/1.5)^2)^2 + 0.35*(x/1.5)^2))};
\addplot[popmuted, dashed] coordinates {(1.5,0) (1.5,10.5)};
\node[popmuted, font=\small, rotate=90, anchor=south] at (axis cs:1.5,5) {$f_0 = 1{,}5$ Hz};
\legend{Résonance aiguë (amortissement faible), Résonance floue (amortissement fort)}
\end{axis}
\end{tikzpicture}
\end{center}
\legende{Courbes de résonance d'un oscillateur de fréquence propre $f_0 = \SI{1,5}{Hz}$ : amplitude $X_m$ en fonction de la fréquence $f$ de l'excitateur. Le pic est d'autant plus haut et étroit que l'amortissement est faible.}
\end{popfigure}

\textbf{Lecture de la courbe de résonance.}\par

La courbe $X_m(f)$ présente trois zones :
\begin{itemize}
\item $f \ll f_0$ : l'amplitude est faible et presque constante ; la masse « suit » lentement l'excitateur ;
\item $f \approx f_0$ : l'amplitude passe par un maximum — c'est la \emph{résonance} ;
\item $f \gg f_0$ : l'amplitude tend vers zéro ; l'oscillateur, par son inertie, ne peut plus suivre les sollicitations trop rapides.
\end{itemize}

\courssub{Exemple chiffré : la suspension d'une voiture sur piste ondulée}

\begin{exemplebox}[Résonance d'une suspension sur la route Yaoundé--Bafoussam]
La suspension d'une voiture se comporte comme un oscillateur masse--ressort amorti de fréquence propre $f_0 = \SI{1,5}{Hz}$. La voiture roule sur une piste dont la surface présente des ondulations régulières (corrugations), distantes de $d = \SI{10}{m}$.

À la vitesse $v$, la voiture rencontre une bosse toutes les $\Delta t = \dfrac{d}{v}$ : la piste joue le rôle d'excitateur de fréquence
\[ f = \frac{1}{\Delta t} = \frac{v}{d}. \]

Il y a résonance lorsque $f = f_0$, c'est-à-dire pour la vitesse :
\[ v = f_0 \times d = 1{,}5 \times 10 = \SI{15}{m.s^{-1}} = \SI{54}{km.h^{-1}}. \]

À $\SI{54}{km.h^{-1}}$, la caisse oscille avec une amplitude maximale : le confort devient exécrable et la tenue de route dangereuse. Le conducteur prudent \emph{change de vitesse} (par exemple ralentir à $\SI{40}{km.h^{-1}}$ ou accélérer franchement) pour s'éloigner de la résonance. Les \emph{amortisseurs}, en augmentant la dissipation d'énergie, aplatissent le pic de résonance : c'est leur rôle précis.
\end{exemplebox}

\begin{methode}[Exploiter une situation de résonance mécanique]
\begin{enumerate}
\item Identifier l'\emph{excitateur} (qui impose la fréquence $f$) et le \emph{résonateur} (l'oscillateur de fréquence propre $f_0$).
\item Calculer la fréquence propre du résonateur : $f_0 = \dfrac{1}{2\pi}\sqrt{\dfrac{k}{m}}$ (ressort) ou $f_0 = \dfrac{1}{2\pi}\sqrt{\dfrac{g}{\ell}}$ (pendule simple).
\item Exprimer la fréquence de l'excitateur en fonction des données (vitesse, périodicité spatiale, régime moteur...).
\item Écrire la condition de résonance $f \approx f_0$ et en tirer la grandeur cherchée.
\item Conclure sur le rôle de l'amortissement : le réduire pour exploiter la résonance (instruments de musique), l'augmenter pour s'en protéger (amortisseurs).
\end{enumerate}
\end{methode}

\courssub{Résonances utiles et résonances nuisibles}

\textbf{Des résonances utiles.}\par

\begin{itemize}
\item \emph{Instruments de musique} : la caisse de résonance d'une guitare, d'un balafon ou d'un xylophone amplifie les sons émis par les cordes ou les lames, en vibrant à des fréquences voisines.
\item \emph{La balançoire} : l'enfant (ou celui qui pousse) fournit de l'énergie à la fréquence propre du pendule ; l'amplitude croît à chaque poussée.
\item \emph{Réception radio} : le circuit d'accord d'un poste radio est réglé pour entrer en résonance avec la fréquence de l'émetteur (RFI, CRTV, radios communautaires...) — il s'agit d'une résonance \emph{électrique}, mais le principe est identique.
\end{itemize}

\textbf{Des résonances nuisibles et leurs parades.}\par

\begin{itemize}
\item \emph{Ponts et passerelles} : une troupe qui marche au pas impose une force périodique ; si la cadence coïncide avec une fréquence propre du pont, les oscillations peuvent devenir dangereuses. Parade : les militaires \emph{rompent le pas} en traversant un pont suspendu ; les marcheurs d'une passerelle bondée sont invités à marcher sans cadence.
\item \emph{Vibrations de machines} : un moteur mal équilibré (mototaxi, groupe électrogène, machine à laver en essorage) excite son support ; à certains régimes, la résonance fait vibrer toute la structure. Parade : plots amortisseurs en caoutchouc, équilibrage des pièces tournantes, modification des masses ou des raideurs pour écarter $f_0$ des fréquences de fonctionnement.
\item \emph{Amortisseurs automobiles} : ils transforment la résonance aiguë de la suspension en résonance floue, limitant l'amplitude des oscillations de la caisse.
\end{itemize}

\begin{savaistu}[La catastrophe du pont de Tacoma]
Le 7 novembre 1940, le pont suspendu de Tacoma Narrows (États-Unis), ouvert à peine quatre mois plus tôt, se mit à osciller violemment sous l'effet d'un vent modéré d'environ $\SI{60}{km.h^{-1}}$. Le vent engendrait des tourbillons périodiques qui entretenaient les oscillations de torsion du tablier : l'amplitude crût jusqu'à plusieurs mètres, et le pont se rompit et s'effondra dans le détroit. Les images filmées de la catastrophe, mondialement célèbres, sont montrées à tous les étudiants ingénieurs. Depuis, tout grand pont est systématiquement étudié en \emph{soufflerie} sur maquette, et sa structure est raidie et amortie pour écarter ses fréquences propres des excitations possibles du vent.
\end{savaistu}

\courssub{Exercice résolu}

\begin{exoresolu}[Résonance d'un pendule élastique vertical]
Un solide de masse $m = \SI{200}{g}$ est suspendu à un ressort de raideur $k = \SI{20}{N.m^{-1}}$. L'extrémité supérieure du ressort est animée d'un mouvement vertical sinusoïdal de fréquence $f$ réglable. On prendra $g = \SI{10}{m.s^{-2}}$ et $\pi^2 \approx 10$.
\begin{enumerate}
\item Calculer la fréquence propre $f_0$ de cet oscillateur.
\item Pour quelle valeur de $f$ observe-t-on la résonance ? Que constate-t-on alors ?
\item On plonge le solide dans un liquide visqueux. Comment la courbe de résonance est-elle modifiée ?
\end{enumerate}
\tcblower
\textbf{Solution détaillée.}
\begin{enumerate}
\item La période propre du pendule élastique est $T_0 = 2\pi\sqrt{\dfrac{m}{k}}$, d'où :
\[ f_0 = \frac{1}{T_0} = \frac{1}{2\pi}\sqrt{\frac{k}{m}} = \frac{1}{2\pi}\sqrt{\frac{20}{0{,}200}} = \frac{1}{2\pi}\sqrt{100} = \frac{10}{2\pi} \approx \SI{1,6}{Hz}. \]
Vérification dimensionnelle : $\sqrt{\dfrac{[\mathrm{N.m^{-1}}]}{[\mathrm{kg}]}} = \sqrt{\dfrac{\mathrm{kg.s^{-2}}}{\mathrm{kg}}} = \mathrm{s^{-1}}$, homogène à une fréquence. La fréquence propre ne dépend pas de $g$ : l'allongement à l'équilibre n'intervient pas dans $T_0$.
\item La résonance se produit lorsque la fréquence de l'excitateur est voisine de la fréquence propre, soit $f \approx f_0 \approx \SI{1,6}{Hz}$. On constate alors que l'amplitude des oscillations du solide devient \emph{maximale} : le solide oscille vigoureusement, bien que le déplacement imposé en haut du ressort soit faible.
\item Le liquide augmente l'amortissement : le pic de résonance devient \emph{moins haut et plus large} (résonance floue). L'amplitude maximale diminue fortement, ce qui illustre le rôle protecteur des amortisseurs.
\end{enumerate}
\end{exoresolu}

\courssub{Vers la résonance électrique (complément utile au Bac)}

\begin{aretenir}[L'essentiel de la section]
\begin{itemize}
\item Dans une oscillation forcée, l'\cle{excitateur} impose sa fréquence $f$ au \cle{résonateur}, qui vibre en régime permanent à cette fréquence $f$ — et non à sa fréquence propre $f_0$.
\item Il y a \cle{résonance mécanique} lorsque $f \approx f_0$ : l'amplitude passe par un maximum.
\item Plus l'amortissement est faible, plus la résonance est \emph{aiguë} (pic étroit et élevé) et dangereuse ; un amortissement fort donne une résonance \emph{floue}.
\item La résonance est recherchée dans les instruments de musique et combattue dans les ponts, machines et suspensions (amortisseurs, rupture de cadence, études en soufflerie).
\end{itemize}
\end{aretenir}

\begin{savaistu}[De la mécanique à l'électricité]
Tout ce que tu viens d'apprendre se transposera presque mot pour mot au circuit électrique $RLC$ alimenté par un générateur de tension sinusoïdale : le générateur joue le rôle de l'excitateur, le circuit celui du résonateur, la résistance $R$ celui de l'amortissement, et l'intensité atteint son maximum lorsque la fréquence du générateur approche la fréquence propre $f_0 = \dfrac{1}{2\pi\sqrt{LC}}$ du circuit : c'est la \emph{résonance électrique}, au cœur de la réception radio. Les analogies entre oscillateurs mécaniques et électriques sont un grand classique du Baccalauréat C : garde-les en mémoire !
\end{savaistu}

\begin{exercice}[Pour t'entraîner]
Une passerelle piétonne a une fréquence propre d'oscillation verticale $f_0 = \SI{2,0}{Hz}$. Un groupe de marcheurs la traverse au pas cadencé, chaque pas durant $\SI{0,50}{s}$.
\begin{enumerate}
\item Calculer la fréquence des excitations dues à la marche.
\item Y a-t-il risque de résonance ? Justifier.
\item Proposer deux mesures permettant d'éviter le danger.
\end{enumerate}
\end{exercice}

\begin{corrige}[Exercice]
\begin{enumerate}
\item La fréquence des pas est $f = \dfrac{1}{T} = \dfrac{1}{0{,}50} = \SI{2,0}{Hz}$.
\item Oui : $f = f_0 = \SI{2,0}{Hz}$, la condition de résonance est réalisée ; l'amplitude des oscillations verticales de la passerelle peut devenir dangereusement grande.
\item On peut : rompre la cadence de la marche (modifier $f$ pour l'écarter de $f_0$) ; augmenter l'amortissement de la structure (amortisseurs) pour aplatir le pic de résonance ; ou raidir la passerelle pour déplacer sa fréquence propre.
\end{enumerate}
\end{corrige}

\newpage

\coursec{Activité d'intégration}

\courssub{Objectifs de l'activité}

À l'issue de cette activité d'intégration, tu seras capable de :
\begin{itemize}
\item dimensionner un oscillateur mécanique (pendule simple) à partir d'un cahier des charges portant sur sa \cle{période propre} ;
\item discuter la validité de l'\cle{approximation des petits angles} et évaluer l'erreur commise ;
\item déterminer l'\cle{amplitude} et la \cle{phase} d'un mouvement harmonique à partir des \cle{conditions initiales} ;
\item exploiter le \cle{théorème de l'énergie mécanique} pour relier amplitude, vitesse maximale et énergie ;
\item analyser un \cle{enregistrement} d'oscillations amorties : mesurer une pseudo-période, conclure sur l'amortissement, proposer une correction technique ;
\item interpréter le rôle d'un entretien à la fréquence propre à la lumière du phénomène de \cle{résonance}.
\end{itemize}

\courssub{Situation}

Le club scientifique du lycée de Nkol-Eton, à Yaoundé, prépare la semaine nationale de la science. Les élèves de Terminale C veulent construire un \textbf{chronomètre à pendule} capable de « battre la seconde » : à chaque passage du pendule par sa position d'équilibre, un contacteur émet un \emph{tic} sonore. Pour que deux tics successifs soient séparés d'exactement \SI{1,00}{s}, la demi-période du pendule doit valoir \SI{1,00}{s}, c'est-à-dire une période propre $T_0 = \SI{2,00}{s}$.

Le matériel disponible : un fil de pêche en nylon inextensible de masse négligeable, une bille d'acier de masse $m = \SI{100}{g}$ assimilable à un point matériel, un support métallique, un rapporteur, un chronomètre électronique et une caméra de téléphone portable filmant à 60 images par seconde.

\textbf{Données :} accélération de la pesanteur à Yaoundé : $g = \SI{9,78}{N/kg}$ ; le pendule sera lâché sans vitesse initiale depuis un angle $\theta_0 = 8,0°$ par rapport à la verticale ; on prendra $\pi^2 \approx 9,87$.

Après construction, les élèves filment le mouvement et relèvent les positions des maxima successifs du même côté :

\begin{center}
\begin{tabularx}{\linewidth}{|l|X|X|X|X|X|X|}
\hline
\rowcolor{popblueL} Rang $n$ du maximum & 0 & 1 & 2 & 3 & 4 & 5 \\
\hline Date $t_n$ (s) & 0,00 & 2,04 & 4,08 & 6,12 & 8,16 & 10,20 \\
\hline Amplitude $\theta_n$ (degrés) & 8,0 & 7,6 & 7,2 & 6,8 & 6,5 & 6,2 \\
\hline
\end{tabularx}
\end{center}

\courssub{Consignes}

\begin{enumerate}
\item \textbf{Dimensionnement.} Établir l'expression de la période propre $T_0$ du pendule simple dans l'approximation des petites oscillations, puis calculer la longueur $\ell$ du fil permettant d'obtenir $T_0 = \SI{2,00}{s}$. Vérifier l'homogénéité de la formule utilisée.
\item \textbf{Validité du modèle.} Justifier que l'approximation $\sin\theta \approx \theta$ est acceptable pour $\theta_0 = 8,0°$ en comparant $\sin\theta_0$ et $\theta_0$ (en radians). Conclure sur la légitimité du modèle du pendule simple harmonique.
\item \textbf{Conditions initiales.} L'élongation angulaire s'écrit $\theta(t) = \theta_m\cos(\omega_0 t + \varphi)$. Déterminer $\theta_m$, $\omega_0$ et $\varphi$ pour un lâcher sans vitesse initiale depuis $\theta_0 = 8,0°$.
\item \textbf{Étude énergétique.} Calculer l'énergie mécanique du pendule (on néglige les frottements et on prend l'origine de l'énergie potentielle au point le plus bas). En déduire la vitesse maximale $v_{\max}$ de la bille au passage à la verticale.
\item \textbf{Exploitation de l'enregistrement.} À partir du tableau, mesurer la pseudo-période $T$ du pendule réel. Comparer $T$ à $T_0$ et calculer l'écart relatif. Que dire de l'amortissement ? Si l'on assimile $T$ à la période propre, proposer une nouvelle longueur $\ell'$ du fil pour retrouver une période de \SI{2,00}{s}.
\item \textbf{Entretien des oscillations.} Expliquer, en mobilisant la notion de résonance, pourquoi de petites poussées périodiques appliquées à la fréquence propre du pendule permettraient d'entretenir les oscillations malgré l'amortissement.
\end{enumerate}

\courssub{Ressources à mobiliser}

\begin{aretenir}[Boîte à outils de l'activité]
\begin{itemize}
\item Pendule simple (petites oscillations) : $\ddot{\theta} + \dfrac{g}{\ell}\theta = 0$, de période propre $T_0 = 2\pi\sqrt{\dfrac{\ell}{g}}$ et pulsation $\omega_0 = \sqrt{\dfrac{g}{\ell}}$.
\item Solution générale : $\theta(t) = \theta_m\cos(\omega_0 t + \varphi)$ ; les conditions initiales $\theta(0)$ et $\dot{\theta}(0)$ fixent $\theta_m$ et $\varphi$.
\item Énergie mécanique conservée (sans frottement) : $E_m = \dfrac{1}{2}mv_{\max}^2 = mg\ell(1-\cos\theta_m)$.
\item Oscillations amorties : régime pseudo-périodique si l'amortissement est faible ; la pseudo-période $T$ est légèrement supérieure à $T_0$.
\item Résonance : un excitateur de fréquence voisine de la fréquence propre transfère un maximum d'énergie à l'oscillateur.
\end{itemize}
\end{aretenir}

\courssub{Démarche et corrigé commenté}

\begin{exoresolu}[Construction du chronomètre à pendule — corrigé détaillé]
Énoncé : dimensionner, analyser et corriger un pendule simple battant la seconde ($T_0 = \SI{2,00}{s}$, $g = \SI{9,78}{N/kg}$, $m = \SI{100}{g}$, $\theta_0 = 8,0°$), puis exploiter l'enregistrement des maxima.
\tcblower
\textbf{1. Dimensionnement de la longueur du fil.}

Dans l'approximation des petites oscillations, la période propre du pendule simple est :
\[ T_0 = 2\pi\sqrt{\frac{\ell}{g}} \quad\Longrightarrow\quad \ell = \frac{g\,T_0^2}{4\pi^2}. \]
Vérification dimensionnelle : $[\ell] = \dfrac{[\mathrm{L\,T^{-2}}]\times[\mathrm{T^2}]}{1} = [\mathrm{L}]$ : la formule est homogène à une longueur.
\[ \ell = \frac{\num{9,78}\times(\num{2,00})^2}{4\times \num{9,87}} = \frac{\num{39,12}}{\num{39,48}} \approx \SI{0,991}{m}. \]
Le fil devra mesurer environ \SI{99,1}{cm} (distance du point de suspension au centre de la bille).

\textbf{2. Validité de l'approximation des petits angles.}

Convertissons : $\theta_0 = 8,0° = 8,0\times\dfrac{\pi}{180} \approx \SI{0,1396}{rad}$. Or $\sin(\num{0,1396}) \approx \num{0,1391}$. L'écart relatif est :
\[ \frac{\num{0,1396}-\num{0,1391}}{\num{0,1391}} \approx \num{0,36}\,\%. \]
L'erreur commise en remplaçant $\sin\theta$ par $\theta$ est inférieure à \num{0,4}\,\% : l'approximation des petits angles est \textbf{excellente} et le modèle harmonique est légitime.

\textbf{3. Amplitude, pulsation et phase.}

La pulsation propre vaut :
\[ \omega_0 = \frac{2\pi}{T_0} = \frac{2\pi}{\num{2,00}} = \SI{3,14}{rad/s}. \]
Le pendule est lâché \emph{sans vitesse initiale} depuis $\theta_0$ : les conditions initiales sont $\theta(0) = \theta_0$ et $\dot{\theta}(0) = 0$.
\begin{align*}
\theta(0) &= \theta_m\cos\varphi = \theta_0,\\
\dot{\theta}(0) &= -\theta_m\,\omega_0\sin\varphi = 0 \;\Rightarrow\; \varphi = 0 \;(\text{mod } \pi).
\end{align*}
Comme $\theta_0 > 0$, on retient $\varphi = 0$ et $\theta_m = \theta_0 = 8,0° \approx \SI{0,140}{rad}$. D'où :
\[ \theta(t) = \num{0,140}\cos(\num{3,14}\,t) \quad (\theta \text{ en rad, } t \text{ en s}). \]

\textbf{4. Énergie mécanique et vitesse maximale.}

L'énergie mécanique, égale à l'énergie potentielle maximale (lâcher sans vitesse), s'écrit :
\[ E_m = mg\ell(1-\cos\theta_m). \]
Avec $\cos(8,0°) \approx \num{0,9903}$, donc $1-\cos\theta_m \approx \num{9,73e-3}$ :
\[ E_m = \num{0,100}\times \num{9,78}\times \num{0,991}\times \num{9,73e-3} \approx \SI{9,4e-3}{J}. \]
Au passage à la verticale, cette énergie est intégralement cinétique : $E_m = \dfrac{1}{2}mv_{\max}^2$, d'où :
\[ v_{\max} = \sqrt{\frac{2E_m}{m}} = \sqrt{\frac{2\times \num{9,4e-3}}{\num{0,100}}} \approx \SI{0,43}{m/s}. \]
Vérification : $v_{\max} = \ell\,\omega_0\,\theta_m = \num{0,991}\times \num{3,14}\times \num{0,140} \approx \SI{0,44}{m/s}$ : les deux méthodes concordent (écart dû aux arrondis).

\textbf{5. Exploitation de l'enregistrement.}

Deux maxima successifs du même côté sont séparés d'une pseudo-période :
\[ T = \frac{t_5 - t_0}{5} = \frac{\num{10,20}}{5} = \SI{2,04}{s}. \]
Écart relatif par rapport à la valeur visée :
\[ \frac{T - T_0}{T_0} = \frac{\num{2,04}-\num{2,00}}{\num{2,00}} = \num{2,0}\,\%. \]
Par ailleurs, l'amplitude décroît régulièrement (de $8,0°$ à $6,2°$ en cinq périodes) : le pendule réel est le siège d'un \textbf{amortissement faible} (frottements de l'air, frottements au point de suspension) ; le régime est \cle{pseudo-périodique}. La pseudo-période mesurée est légèrement supérieure à $T_0$, ce qui est cohérent avec la théorie des oscillations amorties.

Pour corriger, on assimile $T$ à la période propre du pendule corrigé : on cherche $\ell'$ telle que $T_0 = 2\pi\sqrt{\ell'/g} = \SI{2,00}{s}$ alors que le fil actuel donne $T = \SI{2,04}{s}$. Comme $T \propto \sqrt{\ell}$ :
\[ \frac{\ell'}{\ell} = \left(\frac{T_0}{T}\right)^2 = \left(\frac{\num{2,00}}{\num{2,04}}\right)^2 \approx \num{0,961}. \]
\[ \ell' \approx \num{0,991}\times \num{0,961} \approx \SI{0,953}{m}. \]
Il faut donc \textbf{raccourcir le fil d'environ \SI{3,8}{cm}} (soit \num{4}\,\% de sa longueur) pour compenser l'allongement de la période dû à l'amortissement et aux imperfections de réalisation.

\textbf{6. Entretien des oscillations et résonance.}

L'amortissement dissipe à chaque période une fraction de l'énergie mécanique : les oscillations finiraient par s'éteindre. Pour les entretenir, on applique au pendule de petites poussées périodiques (excitation) : le pendule devient un \cle{oscillateur forcé}. Lorsque la fréquence de l'excitateur est proche de la fréquence propre $f_0 = \dfrac{1}{T_0} = \SI{0,50}{Hz}$, le transfert d'énergie vers le pendule est maximal : c'est la \cle{résonance}. Une poussée bien synchronisée (à chaque passage par la verticale, par exemple) compense exactement l'énergie dissipée par les frottements : l'amplitude se stabilise et le chronomètre bat la seconde durablement. C'est le principe du balancier entretenu des horloges anciennes et du mécanisme d'échappement des pendules de précision.
\end{exoresolu}

\begin{attention}[Pièges fréquents dans ce type d'activité]
\begin{itemize}
\item Ne pas confondre la \textbf{période} $T_0 = \SI{2,00}{s}$ et la \textbf{demi-période} (\SI{1,00}{s}) : le « battement de seconde » correspond à deux passages par la verticale par période.
\item Toujours convertir l'angle en \textbf{radians} avant d'utiliser $\sin\theta \approx \theta$ ou les formules d'énergie.
\item Dans $E_m = mg\ell(1-\cos\theta_m)$, l'argument du cosinus peut rester en degrés si la calculatrice est réglée en degrés — mais la cohérence du réglage est exigée.
\item Une correction de longueur se déduit de $T \propto \sqrt{\ell}$ : pour réduire $T$ de \num{2}\,\%, il faut réduire $\ell$ d'environ \num{4}\,\%, pas de \num{2}\,\%.
\end{itemize}
\end{attention}

\courssub{Bilan de l'activité}

Cette activité a permis de mobiliser l'ensemble des savoir-faire de la leçon dans une situation authentique de conception :
\begin{itemize}
\item la \textbf{mise en équation} et l'exploitation de la période propre $T_0 = 2\pi\sqrt{\ell/g}$ ont conduit au dimensionnement $\ell \approx \SI{0,99}{m}$ ;
\item la discussion de l'approximation des petits angles a montré qu'un modèle physique n'est valable que sous des \textbf{conditions précises}, qu'il faut toujours vérifier quantitativement ;
\item les \textbf{conditions initiales} ont fixé sans ambiguïté amplitude et phase ($\theta_m = \SI{0,140}{rad}$, $\varphi = 0$) ;
\item le \textbf{bilan énergétique} a relié l'amplitude à la vitesse maximale ($v_{\max} \approx \SI{0,43}{m/s}$) sans résoudre l'équation différentielle ;
\item l'\textbf{exploitation d'un enregistrement} a révélé un régime pseudo-périodique faiblement amorti et a motivé une correction technologique du prototype ;
\item enfin, la notion de \textbf{résonance} a fourni la solution pour entretenir les oscillations, ouvrant la voie aux applications : horloges, sismographes, et même le balancier des montres mécaniques que l'on retrouve jusque dans les ateliers de réparation du marché central de Yaoundé.
\end{itemize}

\begin{savaistu}[Le pendule qui mesure $g$]
La relation $T_0 = 2\pi\sqrt{\ell/g}$ peut être « retournée » : en mesurant précisément $\ell$ et $T_0$, on détermine $g = \dfrac{4\pi^2\ell}{T_0^2}$. C'est l'une des méthodes historiques de mesure de la pesanteur. Au Cameroun, $g$ varie légèrement entre Yaoundé ($\SI{9,78}{N/kg}$) et le sommet du mont Cameroun ($\approx \SI{9,77}{N/kg}$) : un pendule suffisamment précis permettrait de détecter cette variation !
\end{savaistu}

\newpage

\coursec{Méthodes et exercices résolus}

\coursec{Oscillateurs mécaniques}

\courssub{Pendule élastique horizontal : modélisation et équation du mouvement}

\begin{definition}[Force de rappel élastique]
Un ressort idéal (de masse négligeable, sans frottement interne) exerce une force de rappel proportionnelle à son allongement (ou compression) $\Delta\ell$ par rapport à sa longueur au repos :
\[ \vec{F}_{\text{rappel}} = -k\,\vec{\Delta\ell} \]
où $k$ (en $\mathrm{N\,m^{-1}}$) est la \cle{raideur} du ressort. Le signe moins indique que la force s'oppose à l'allongement.
\end{definition}

\begin{propriete}[Équation différentielle et solution du pendule élastique horizontal]
Un solide de masse $m$ glissant sans frottement sur un plan horizontal, attaché à un ressort de raideur $k$ fixé en $O$, subit la force de rappel $\vec{F} = -k\,x\,\vec{u}_x$ (origine $O$ à la position d'équilibre, axe $Ox$ orienté dans le sens de l'allongement).
La projection de la relation fondamentale de la dynamique (RFD) sur $(Ox)$ donne :
\[ m\,\ddot{x} = -k\,x \quad\Longrightarrow\quad \boxed{\ddot{x} + \omega_0^2\,x = 0} \qquad\text{avec}\quad \omega_0 = \sqrt{\frac{k}{m}} \]
$\omega_0$ est la \cle{pulsation propre} (en $\mathrm{rad\,s^{-1}}$). La solution générale est :
\[ \boxed{x(t) = X_m\cos(\omega_0 t + \varphi)} \]
où $X_m \ge 0$ est l'\cle{amplitude} et $\varphi$ la \cle{phase à l'origine} (déterminée par les conditions initiales).
La \cle{période propre} et la \cle{fréquence propre} valent :
\[ T_0 = \frac{2\pi}{\omega_0} = 2\pi\sqrt{\frac{m}{k}}, \qquad f_0 = \frac{1}{T_0} = \frac{1}{2\pi}\sqrt{\frac{k}{m}}. \]
\end{propriete}

\begin{aretenir}[Isochronisme]
Pour un oscillateur harmonique non amorti, la période $T_0$ ne dépend \textbf{que} de $m$ et $k$ (ou $\ell$ et $g$ pour le pendule simple), \textbf{pas} de l'amplitude $X_m$. C'est la propriété d'\cle{isochronisme}.
\end{aretenir}

\begin{methode}[Mettre en équation le mouvement d'un point soumis à une force de rappel]
\begin{enumerate}
\item \textbf{Choisir un référentiel galiléen} (le laboratoire, assimilé à galiléen) et un \textbf{repère} : origine $O$ à la position d'équilibre, axe $(Ox)$ orienté dans le sens du mouvement possible.
\item \textbf{Faire le bilan des forces} : poids $\vec{P}$, réaction du support $\vec{R}$ (ou tension), force de rappel du ressort $\vec{F} = -k\,\vec{\Delta\ell}$. \textbf{Schéma obligatoire.}
\item \textbf{Appliquer la RFD} : $\sum \vec{F}_{\text{ext}} = m\,\vec{a}$, puis projeter sur $(Ox)$.
\item \textbf{Exprimer l'allongement} en fonction de $x$ (ressort horizontal : $\Delta\ell = x$).
\item \textbf{Identifier} l'équation canonique $\ddot{x} + \omega_0^2\,x = 0$ et en déduire $\omega_0 = \sqrt{k/m}$, $T_0 = 2\pi\sqrt{m/k}$.
\end{enumerate}
\textbf{Réflexes} : vérifier l'homogénéité ($[k/m] = \mathrm{N\,m^{-1}\,kg^{-1}} = \mathrm{s^{-2}}$) ; toujours partir du bilan de forces et de la RFD projetée.
\end{methode}

\begin{exoresolu}[Pendule élastique horizontal : lancement depuis l'équilibre]
Un solide de masse $m = \SI{200}{g}$, accroché à un ressort horizontal de raideur $k = \SI{20}{N.m^{-1}}$, glisse sans frottement sur un banc horizontal. À $t = 0$, on lance le solide depuis sa position d'équilibre avec une vitesse $v_0 = \SI{0,50}{m.s^{-1}}$ dans le sens positif.
\begin{enumerate}
\item Calculer la pulsation propre $\omega_0$ et la période $T_0$.
\item Déterminer l'équation horaire $x(t)$.
\item Calculer l'énergie mécanique de l'oscillateur.
\end{enumerate}
\tcblower
\textbf{1. Pulsation et période.}
\[ \omega_0 = \sqrt{\frac{k}{m}} = \sqrt{\frac{20}{0{,}200}} = \sqrt{100} = \SI{10}{rad.s^{-1}}, \qquad T_0 = \frac{2\pi}{\omega_0} = \frac{2\pi}{10} \approx \SI{0,628}{s}. \]

\textbf{2. Équation horaire.}
Solution générale : $x(t) = X_m\cos(\omega_0 t + \varphi)$, $v(t) = -X_m\omega_0\sin(\omega_0 t + \varphi)$.
Conditions initiales ($t=0$) : $x(0) = 0$, $v(0) = v_0 = +0,50\,\mathrm{m/s}$.
\[ x(0) = X_m\cos\varphi = 0 \;\Rightarrow\; \varphi = \pm\frac{\pi}{2}. \]
\[ v(0) = -X_m\omega_0\sin\varphi = v_0 > 0 \;\Rightarrow\; \sin\varphi < 0 \;\Rightarrow\; \varphi = -\frac{\pi}{2}. \]
Amplitude : $v_0 = X_m\omega_0 \;\Rightarrow\; X_m = \frac{v_0}{\omega_0} = \frac{0{,}50}{10} = \SI{5,0e-2}{m} = \SI{5,0}{cm}$.
\[ \boxed{x(t) = 5{,}0\times 10^{-2}\cos\left(10\,t - \frac{\pi}{2}\right) \quad (x \text{ en m},\, t \text{ en s})} \]
Vérification : $x(0)=0$, $\dot{x}(0)=+0,50\,\mathrm{m/s}$. Conforme.

\textbf{3. Énergie mécanique.}
\[ E_m = \frac{1}{2}k X_m^2 = \frac{1}{2}\times 20\times (5{,}0\times 10^{-2})^2 = \SI{2,5e-2}{J} = \SI{25}{mJ}. \]
\end{exoresolu}

\courssub{Pendule élastique vertical : rôle du poids}

\begin{propriete}[Équilibre et mouvement vertical]
Pour un ressort vertical de raideur $k$ portant une masse $m$, l'allongement à l'équilibre $\Delta\ell_0$ vérifie $k\,\Delta\ell_0 = mg$ (axe $Ox$ descendant).
Si on écarte la masse de $x$ par rapport à l'équilibre, l'allongement total est $\Delta\ell_0 + x$. La RFD projetée sur $(Ox)$ donne :
\[ m\ddot{x} = mg - k(\Delta\ell_0 + x) = \underbrace{mg - k\Delta\ell_0}_{=\,0} - kx. \]
On retrouve \textbf{exactement} la même équation $\ddot{x} + \frac{k}{m}x = 0$. Le poids ne fait que déplacer la position d'équilibre ; il ne modifie ni $\omega_0$, ni $T_0$.
\end{propriete}

\begin{exoresolu}[Pendule élastique vertical : la suspension de mototaxi]
Sur une mototaxi de Douala, on assimile la suspension à un ressort vertical de raideur $k = \SI{4,0e4}{N.m^{-1}}$. Une masse $m = \SI{100}{kg}$ (pilote + passager + partie arrière) repose sur le ressort. À l'équilibre, le ressort est comprimé de $\Delta\ell_0$. On écarte la masse de $a = \SI{3,0}{cm}$ vers le bas puis on la lâche sans vitesse initiale.
\begin{enumerate}
\item Établir l'équation différentielle du mouvement de la masse.
\item Calculer la période propre $T_0$ des oscillations.
\item Écrire l'équation horaire $x(t)$.
\end{enumerate}
\tcblower
\textbf{1. Équation différentielle.}
Référentiel terrestre supposé galiléen. Axe $(Ox)$ vertical descendant, origine $O$ = position d'équilibre.
Bilan des forces : $\vec{P} = m\vec{g}$ (vers le bas), $\vec{F}_{\text{rappel}} = -k(\Delta\ell_0 + x)\,\vec{u}_x$ (vers le haut).
Condition d'équilibre : $mg - k\Delta\ell_0 = 0$.
RFD : $m\ddot{x} = mg - k(\Delta\ell_0 + x) = -kx$.
\[ \boxed{\ddot{x} + \frac{k}{m}\,x = 0} \quad\text{avec}\quad \omega_0 = \sqrt{\frac{k}{m}}. \]

\textbf{2. Période propre.}
\[ T_0 = 2\pi\sqrt{\frac{m}{k}} = 2\pi\sqrt{\frac{100}{4{,}0\times 10^{4}}} = 2\pi\sqrt{2{,}5\times 10^{-3}} \approx \SI{0,314}{s} \approx \SI{0,31}{s}. \]

\textbf{3. Équation horaire.}
Lâché sans vitesse depuis $x_0 = +a = +\SI{3,0}{cm} = +\SI{3,0e-2}{m}$.
$\Rightarrow \varphi = 0$, $X_m = a = \SI{3,0e-2}{m}$.
\[ \boxed{x(t) = 3{,}0\times 10^{-2}\cos\left(\frac{2\pi}{0{,}314}\,t\right) \approx 3{,}0\times 10^{-2}\cos(20{,}0\,t) \quad (x \text{ en m},\, t \text{ en s})} \]
\end{exoresolu}

\begin{popfigure}
\begin{tikzpicture}[scale=1.2, >=Stealth]
% Sol
\draw[fill=popmuted!30] (-1.5,-0.2) rectangle (1.5,0);
% Ressort (comprimé à l'équilibre)
\def\coilheight{2.5}
\def\coilwidth{0.4}
\draw[popblue, thick] (0,0) -- (0,0.3); % fixation
\draw[popblue, thick, decorate, decoration={coil, aspect=0.4, segment length=3mm, amplitude=2.5mm}] (0,0.3) -- (0,2.5); % ressort
% Masse
\draw[fill=poporange!50, draw=popdark, thick] (-0.6,2.5) rectangle (0.6,3.1);
\node at (0,2.8) {$m$};
% Forces à l'équilibre (x=0)
\draw[->, thick, popgreen] (1.0, 2.8) -- (1.0, 3.5) node[right] {$\vec{P}=m\vec{g}$};
\draw[->, thick, popgreen] (1.0, 2.8) -- (1.0, 2.1) node[right] {$\vec{F}_0=-k\Delta\ell_0\,\vec{u}_x$};
% Forces en mouvement (x>0)
\begin{scope}[xshift=4cm]
\draw[fill=popmuted!30] (-1.5,-0.2) rectangle (1.5,0);
\draw[popblue, thick] (0,0) -- (0,0.3);
\draw[popblue, thick, decorate, decoration={coil, aspect=0.4, segment length=3mm, amplitude=2.5mm}] (0,0.3) -- (0,2.2);
\draw[fill=poporange!50, draw=popdark, thick] (-0.6,2.2) rectangle (0.6,2.8);
\node at (0,2.5) {$m$};
\draw[->, thick, popgreen] (1.0, 2.5) -- (1.0, 3.2) node[right] {$\vec{P}=m\vec{g}$};
\draw[->, thick, popred] (1.0, 2.5) -- (1.0, 1.8) node[right] {$\vec{F}=-k(\Delta\ell_0+x)\,\vec{u}_x$};
\draw[->, thick, popblue] (-1.2, 2.5) -- (-1.8, 2.5) node[left] {$\vec{x}$};
\node at (0, -0.5) {Écarté de $x>0$};
\end{scope}
\node at (0, -0.5) {Position d'équilibre ($x=0$)};
\end{tikzpicture}
\legende{Bilan des forces sur un pendule élastique vertical : à l'équilibre (gauche) et en mouvement (droite). Le poids est constant, la force de rappel varie avec $x$.}
\end{popfigure}

\courssub{Pendule pesant et modèle du pendule simple}

\begin{definition}[Pendule pesant et pendule simple]
Un \cle{pendule pesant} est un solide quelconque oscillant autour d'un axe horizontal fixe sous l'action de son poids. Le \cle{pendule simple} est un modèle idéalisé : une masse ponctuelle $m$ attachée à un fil inextensible, sans masse, de longueur $\ell$, oscillant sans frottement.
\end{definition}

\begin{propriete}[Équation du pendule simple (petites oscillations)]
Soit $\theta$ l'angle entre le fil et la verticale descendante. Le bilan des forces (poids $\vec{P}$, tension $\vec{T}$) et le théorème du moment cinétique (ou RFD en composantes tangentielles) donnent :
\[ m\ell\,\ddot{\theta} = -mg\sin\theta \quad\Longrightarrow\quad \ddot{\theta} + \frac{g}{\ell}\sin\theta = 0. \]
Pour de \textbf{petites oscillations} ($\theta \ll 1$ rad, soit $\theta_m \lesssim 10^\circ$), $\sin\theta \approx \theta$. L'équation devient harmonique :
\[ \boxed{\ddot{\theta} + \omega_0^2\,\theta = 0} \qquad\text{avec}\quad \omega_0 = \sqrt{\frac{g}{\ell}}, \quad T_0 = 2\pi\sqrt{\frac{\ell}{g}}. \]
La solution est $\theta(t) = \theta_m\cos(\omega_0 t + \varphi)$.
\end{propriete}

\begin{attention}[Condition des petites oscillations]
La formule $T_0 = 2\pi\sqrt{\ell/g}$ n'est valide que si $\theta_m \lesssim 10^\circ$ ($\approx 0,17$ rad). Au-delà, la période dépend de l'amplitude (anharmonicité) : $T > T_0$. Ne jamais l'appliquer sans vérifier cette condition.
\end{attention}

\begin{methode}[Exploiter un enregistrement d'oscillations (pendule simple)]
\begin{enumerate}
\item \textbf{Lire l'amplitude angulaire} $\theta_m$ (en degrés ou radians) : vérifier $\theta_m \lesssim 10^\circ$.
\item \textbf{Mesurer la période} $T$ : durée entre deux crêtes consécutives (ou mieux : mesurer $n$ périodes et diviser par $n$ pour réduire l'incertitude). $\omega_0 = 2\pi/T$.
\item \textbf{En déduire $g$} : $g = 4\pi^2\ell/T^2$.
\item \textbf{Calculer l'écart relatif} par rapport à la valeur tabulée ($g \approx \SI{9,81}{m.s^{-2}}$ au Cameroun) : $\frac{|g_{\text{mes}}-g_{\text{tab}}|}{g_{\text{tab}}}$.
\end{enumerate}
\end{methode}

\begin{exoresolu}[Exploitation d'un enregistrement au lycée de Ngaoundéré]
Un élève enregistre l'élongation angulaire $\theta(t)$ d'un pendule simple de longueur $\ell = \SI{1,00}{m}$. L'enregistrement montre des oscillations sinusoïdales d'amplitude $\theta_m = 8,0^\circ$ ; 10 oscillations complètes durent $\SI{20,1}{s}$.
\begin{enumerate}
\item Déterminer la période $T$ mesurée.
\item En déduire la valeur expérimentale de $g$.
\item Comparer à la valeur tabulée $g = \SI{9,81}{m.s^{-2}}$ (écart relatif).
\end{enumerate}
\tcblower
\textbf{1. Période mesurée.}
\[ T = \frac{20{,}1}{10} = \SI{2,01}{s}. \]

\textbf{2. Valeur expérimentale de $g$.}
L'amplitude $8,0^\circ < 10^\circ$ valide l'approximation $\sin\theta \approx \theta$.
\[ g = \frac{4\pi^2\ell}{T^2} = \frac{4\pi^2\times 1{,}00}{(2{,}01)^2} = \frac{39{,}48}{4{,}0401} \approx \SI{9,77}{m.s^{-2}}. \]

\textbf{3. Écart relatif.}
\[ \frac{|g_{\text{mes}} - g_{\text{tab}}|}{g_{\text{tab}}} = \frac{|9{,}77 - 9{,}81|}{9{,}81} \approx 0{,}41\,\%. \]
\textbf{Conclusion} : l'écart ($< 1\,\%$) est très faible, compatible avec les incertitudes de mesure (longueur, chronométrage). Le pendule simple est un excellent gravimètre de laboratoire.
\end{exoresolu}

\courssub{Pendule de torsion et étude énergétique}

\begin{definition}[Couple de rappel d'un fil de torsion]
Un fil de torsion (ou lame) exerce un couple de rappel proportionnel à l'angle de torsion $\theta$ :
\[ \mathcal{M} = -C\,\theta \]
où $C$ (en $\mathrm{N\,m\,rad^{-1}}$) est la \cle{constante de torsion} du fil.
\end{definition}

\begin{propriete}[Mouvement du pendule de torsion]
Un solide de moment d'inertie $J$ (par rapport à l'axe du fil) suspendu à un fil de constante $C$ obéit au théorème du moment cinétique : $J\,\ddot{\theta} = -C\,\theta$.
\[ \boxed{\ddot{\theta} + \frac{C}{J}\,\theta = 0} \quad\Rightarrow\quad \omega_0 = \sqrt{\frac{C}{J}}, \quad T_0 = 2\pi\sqrt{\frac{J}{C}}. \]
L'énergie mécanique est conservée (pas de frottement) :
\[ E_m = E_c + E_p = \frac{1}{2}J\dot{\theta}^2 + \frac{1}{2}C\theta^2 = \frac{1}{2}C\theta_m^2 = \text{constante}. \]
\end{propriete}

\begin{methode}[Étudier un pendule de torsion / Déterminer $J$ d'un solide]
\begin{enumerate}
\item Repérer le couple de rappel $\mathcal{M} = -C\theta$ et appliquer $J\ddot{\theta} = -C\theta$.
\item Identifier $\omega_0 = \sqrt{C/J}$ et $T_0 = 2\pi\sqrt{J/C}$.
\item \textbf{Énergie} : $E_m = \frac{1}{2}C\theta_m^2$.
\item \textbf{Méthode expérimentale pour $J$ inconnu} : mesurer $T_0$ (solide seul), puis $T_0'$ (solide + étalon de moment $J_0$ connu). Éliminer $C$ :
\[ \frac{T_0'^2}{T_0^2} = \frac{J+J_0}{J} \quad\Longrightarrow\quad J = \frac{J_0}{\frac{T_0'^2}{T_0^2} - 1}. \]
\end{enumerate}
\end{methode}

\begin{exoresolu}[Constante de torsion d'un fil et énergie]
Une tige homogène de moment d'inertie $J = \SI{2,0e-3}{kg.m^2}$ est suspendue horizontalement par un fil de torsion de constante $C$ inconnue. Écartée de sa position d'équilibre puis lâchée, elle effectue 20 oscillations en $\SI{25,1}{s}$.
\begin{enumerate}
\item Calculer la constante de torsion $C$ du fil.
\item Calculer l'énergie mécanique si l'amplitude angulaire est $\theta_m = 0{,}20$ rad.
\end{enumerate}
\tcblower
\textbf{1. Constante de torsion.}
Période mesurée : $T_0 = \dfrac{25{,}1}{20} = \SI{1,255}{s}$.
\[ T_0 = 2\pi\sqrt{\frac{J}{C}} \;\Rightarrow\; C = \frac{4\pi^2 J}{T_0^2} = \frac{4\pi^2\times 2{,}0\times 10^{-3}}{(1{,}255)^2} = \frac{7{,}90\times 10^{-2}}{1{,}575} \approx \SI{5,02e-2}{N.m.rad^{-1}}. \]
Vérification dimensionnelle : $[J/T^2] = \mathrm{kg\,m^2\,s^{-2}} = \mathrm{N\,m}$ (le radian est sans dimension). Cohérent.

\textbf{2. Énergie mécanique.}
\[ E_m = \frac{1}{2}C\theta_m^2 = \frac{1}{2}\times 5{,}02\times 10^{-2}\times (0{,}20)^2 = \SI{1,00e-3}{J} = \SI{1,00}{mJ}. \]
\end{exoresolu}

\begin{aretenir}[Énergie d'un oscillateur harmonique (rappel synthétique)]
\begin{center}
\begin{tabularx}{\linewidth}{|l|X|X|}\hline
\textbf{Oscillateur} & \textbf{Énergie potentielle $E_p$} & \textbf{Énergie mécanique $E_m$} \\ \hline
Ressort (horizontal/vertical) & $\frac{1}{2}k x^2$ & $\frac{1}{2}k X_m^2 = \frac{1}{2}m\omega_0^2 X_m^2$ \\ \hline
Pendule simple (petites osc.) & $mg\ell(1-\cos\theta) \approx \frac{1}{2}mg\ell\theta^2$ & $\frac{1}{2}mg\ell\theta_m^2 = \frac{1}{2}m\ell^2\omega_0^2\theta_m^2$ \\ \hline
Pendule de torsion & $\frac{1}{2}C\theta^2$ & $\frac{1}{2}C\theta_m^2 = \frac{1}{2}J\omega_0^2\theta_m^2$ \\ \hline
\end{tabularx}
\end{center}
Partout : $E_m = E_c + E_p = \text{cste}$ ; $v_{\max} = \omega_0 X_m$ (ou $\dot{\theta}_{\max} = \omega_0\theta_m$).
\end{aretenir}

\courssub{Oscillations amorties : régimes pseudo-périodique et apériodique}

\begin{definition}[Oscillateur amorti]
Lorsqu'un fluide (air, huile) ou des frottements solides agissent, une force de frottement $\vec{f} = -\lambda\vec{v}$ (visqueuse, $\lambda > 0$ en $\mathrm{kg\,s^{-1}}$) s'ajoute. L'équation devient :
\[ m\ddot{x} + \lambda\dot{x} + kx = 0 \quad\Longleftrightarrow\quad \ddot{x} + 2\beta\dot{x} + \omega_0^2 x = 0 \]
avec $\beta = \lambda/(2m)$ (coefficient d'amortissement, en $\mathrm{s^{-1}}$).
\end{definition}

\begin{propriete}[Deux régimes d'amortissement (qualitatif)]
Le discriminant $\Delta = \beta^2 - \omega_0^2$ détermine le régime :
\begin{itemize}
\item \textbf{Régime pseudo-périodique} ($\beta < \omega_0$, amortissement faible) : l'oscillateur oscille avec une \cle{pseudo-période} $T = \frac{2\pi}{\sqrt{\omega_0^2-\beta^2}} > T_0$ et une amplitude qui décroît exponentiellement : $X_m(t) = X_{m0}\,e^{-\beta t}$. L'énergie mécanique diminue (dissipée en chaleur).
\item \textbf{Régime apériodique} ($\beta > \omega_0$, amortissement fort) : pas d'oscillation, retour monotone à l'équilibre (somme de deux exponentielles décroissantes).
\item \textbf{Régime critique} ($\beta = \omega_0$) : retour le plus rapide à l'équilibre sans oscillation.
\end{itemize}
\end{propriete}

\begin{aretenir}[Effets de l'amortissement]
\begin{itemize}
\item La pseudo-période $T$ est \textbf{légèrement supérieure} à $T_0$ ($T \approx T_0$ si $\beta \ll \omega_0$).
\item L'amplitude et l'énergie mécanique décroissent exponentiellement.
\item Les amortisseurs des véhicules (mototaxis, voitures) sont conçus pour être en régime pseudo-périodique légèrement amorti (confort) ou critique (tenue de route).
\end{itemize}
\end{aretenir}

\courssub{Oscillations forcées et résonance mécanique}

\begin{definition}[Oscillateur forcé]
On soumet l'oscillateur à une force extérieure sinusoïdale $F(t) = F_0\cos(\omega t)$ de pulsation $\omega$ (imposée par l'excitateur). L'équation différentielle est :
\[ \ddot{x} + 2\beta\dot{x} + \omega_0^2 x = \frac{F_0}{m}\cos(\omega t). \]
En régime permanent (après extinction des transitoires), la réponse est sinusoïdale à la \textbf{même pulsation} $\omega$ que l'excitateur : $x(t) = A(\omega)\cos(\omega t + \varphi(\omega))$.
\end{definition}

\begin{propriete}[Courbe de résonance et résonance mécanique]
L'amplitude $A(\omega)$ dépend de la fréquence de l'excitateur :
\[ A(\omega) = \frac{F_0/m}{\sqrt{(\omega_0^2-\omega^2)^2 + 4\beta^2\omega^2}}. \]
\begin{itemize}
\item \textbf{Résonance} : $A(\omega)$ est maximale pour une pulsation $\omega_r \approx \omega_0$ (exactement $\omega_r = \sqrt{\omega_0^2 - 2\beta^2}$ pour l'amplitude, $\omega_0$ pour l'énergie).
\item \textbf{Pic de résonance} : plus l'amortissement $\beta$ est faible, plus le pic est \textbf{haut} et \textbf{fin} (résonance aiguë, dangereuse). Plus $\beta$ est grand, plus le pic est \textbf{large} et \textbf{bas} (résonance large, amortie).
\item En résonance, le déphasage entre l'excitateur et la réponse est $\varphi = -\pi/2$ (vitesse en phase avec la force : puissance absorbée maximale).
\end{itemize}
\end{propriete}

\begin{methode}[Oscillations forcées et résonance : réflexes qualitatifs]
\begin{enumerate}
\item \textbf{Identifier l'excitateur} et sa fréquence $f$ (ou pulsation $\omega$).
\item \textbf{Comparer $f$ à la fréquence propre $f_0$} de l'oscillateur (calculée sans l'excitateur).
\item \textbf{Résonance} : si $f \approx f_0$, l'amplitude devient très grande (dangereuse pour les structures, recherchée pour les capteurs, les antennes, la balançoire).
\item \textbf{Rôle de l'amortissement} : il limite l'amplitude en résonance et élargit la courbe. Un bon amortisseur "aplatit" la résonance.
\end{enumerate}
\end{methode}

\begin{exoresolu}[Résonance d'une suspension de mototaxi sur piste bosselée]
Une mototaxi et son pilote, de masse totale $m = \SI{150}{kg}$, reposent sur une suspension assimilée à un ressort de raideur $k = \SI{6,0e4}{N.m^{-1}}$. La moto roule à vitesse $v$ constante sur une piste dont les bosses sont espacées régulièrement de $d = \SI{5,0}{m}$.
\begin{enumerate}
\item Calculer la fréquence propre $f_0$ de la suspension.
\item À quelle vitesse $v$ la suspension entre-t-elle en résonance ? Commenter le danger.
\end{enumerate}
\tcblower
\textbf{1. Fréquence propre.}
\[ f_0 = \frac{1}{2\pi}\sqrt{\frac{k}{m}} = \frac{1}{2\pi}\sqrt{\frac{6{,}0\times 10^{4}}{150}} = \frac{1}{2\pi}\sqrt{400} = \frac{20}{2\pi} \approx \SI{3,18}{Hz}. \]

\textbf{2. Vitesse de résonance.}
Les bosses agissent comme un excitateur spatial de période $d$. La fréquence temporelle d'excitation est $f = \frac{v}{d}$ (une bosse franchie tous les $d/v$ secondes).
Résonance lorsque $f = f_0$ :
\[ v = f_0\,d = 3{,}18 \times 5{,}0 \approx \SI{15,9}{m.s^{-1}} \approx \SI{57}{km.h^{-1}}. \]
\textbf{Conclusion} : vers $\SI{57}{km.h^{-1}}$, la fréquence des secousses coïncide avec la fréquence propre de la suspension. L'amplitude des oscillations devient maximale : perte d'adhérence, inconfort, risque de casse mécanique. Le pilote \textbf{doit éviter cette vitesse} (ralentir ou accélérer franchement). Les amortisseurs (huile, gaz) dissipent l'énergie et réduisent le pic de résonance, mais ne l'éliminent pas.
\end{exoresolu}

\begin{savaistu}[Résonance et génie civil au Cameroun]
\begin{itemize}
\item \textbf{Pont sur le Wouri (Douala)} : le vent peut exciter des modes propres de la structure. Les ingénieurs calculent les fréquences propres et installent des \cle{amortisseurs de masse accordés} (TMD) pour éviter la résonance éolienne (type pont de Tacoma Narrows).
\item \textbf{Bâtiments antisismiques} (Yaoundé, zone sismique modérée) : on isole la base (paliers en caoutchouc/acier) pour abaisser la fréquence propre du bâtiment loin des fréquences dominantes du séisme ($\approx 1$ à $10\,\mathrm{Hz}$).
\item \textbf{Électrotechnique} : le circuit RLC série est l'analogue exact de l'oscillateur mécanique forcé ($L \leftrightarrow m$, $1/C \leftrightarrow k$, $R \leftrightarrow \lambda$). La résonance électrique ($f_0 = 1/(2\pi\sqrt{LC})$) sert à accorder les récepteurs radio (MTN, Orange, CRTV) sur une fréquence précise.
\end{itemize}
\end{savaistu}

\begin{attention}[Les 5 erreurs qui coûtent des points au Bac]
\begin{enumerate}
\item \textbf{Oublier le référentiel galiléen et le bilan des forces} avant la RFD. Sans ces deux lignes, l'équation différentielle n'est pas justifiée (0 point sur la mise en équation).
\item \textbf{Confondre allongement total et écart à l'équilibre} (pendule vertical) : la force de rappel vaut $-k(\Delta\ell_0 + x)$. C'est la condition d'équilibre $k\Delta\ell_0 = mg$ qui simplifie l'équation. L'oublier fait croire à tort que le poids change la période.
\item \textbf{Se tromper de quadrant pour la phase $\varphi$} : $\tan\varphi = -v_0/(\omega_0 x_0)$ ne suffit pas. Vérifier \textbf{les signes} de $x_0 = X_m\cos\varphi$ et $v_0 = -X_m\omega_0\sin\varphi$. Toujours conclure par une vérification numérique de $x(0)$ et $v(0)$.
\item \textbf{Erreurs de conversion d'unités} : masse en $\mathrm{kg}$ (pas g), amplitude en $\mathrm{m}$ (pas cm), $\omega_0$ en $\mathrm{rad/s}$ (pas Hz). Vérifier l'homogénéité : $[k/m] = \mathrm{s^{-2}}$, $[\sqrt{\ell/g}] = \mathrm{s}$.
\item \textbf{Appliquer $T_0 = 2\pi\sqrt{\ell/g}$ hors petites oscillations} : cette formule exige $\theta_m \lesssim 10^\circ$ ($\sin\theta \approx \theta$). Le préciser dans la copie montre la maîtrise des conditions d'application ; l'omettre peut coûter le point de justification.
\end{enumerate}
\end{attention}

\newpage

\coursec{Exercices}

\courssub{Je vérifie mes connaissances}

\begin{exercice}[QCM : une seule réponse exacte]
Pour chaque question, choisir la bonne réponse en justifiant brièvement.
\begin{enumerate}
\item La période propre d'un pendule élastique horizontal de masse $m$ et de raideur $k$ est :
\begin{enumerate}
\item $T_0 = 2\pi\sqrt{\dfrac{k}{m}}$ \quad
\item $T_0 = 2\pi\sqrt{\dfrac{m}{k}}$ \quad
\item $T_0 = \dfrac{1}{2\pi}\sqrt{\dfrac{m}{k}}$ \quad
\item $T_0 = 2\pi\sqrt{mk}$
\end{enumerate}
\item Si l'on quadruple la masse accrochée à un ressort, la période propre est :
\begin{enumerate}
\item divisée par 2 \quad \item multipliée par 2 \quad \item multipliée par 4 \quad \item inchangée
\end{enumerate}
\item Un pendule simple de longueur $\ell = \SI{1,0}{m}$ a pour période propre (prendre $g = \SI{9,8}{m.s^{-2}}$, $\pi^2 \simeq 9,87$) :
\begin{enumerate}
\item environ $\SI{0,5}{s}$ \quad \item environ $\SI{1,0}{s}$ \quad \item environ $\SI{2,0}{s}$ \quad \item environ $\SI{4,0}{s}$
\end{enumerate}
\item À la résonance mécanique d'un oscillateur faiblement amorti :
\begin{enumerate}
\item l'amplitude est nulle \quad \item la fréquence de l'excitateur est très différente de la fréquence propre \quad \item l'amplitude passe par un maximum \quad \item la période devient infinie
\end{enumerate}
\end{enumerate}
\end{exercice}

\begin{exercice}[Vrai ou faux ? Justifier]
Répondre par vrai ou faux en justifiant chaque réponse par une phrase ou un calcul.
\begin{enumerate}
\item L'équation différentielle $\ddot{x} + \omega_0^2\, x = 0$ admet pour solution $x(t) = X_m\cos(\omega_0 t + \varphi)$ quelles que soient les valeurs de $X_m$ et $\varphi$.
\item La période propre d'un pendule élastique dépend de l'amplitude des oscillations.
\item Un pendule simple transporté sur la Lune (où $g_{\text{L}} \simeq g/6$) oscille plus lentement que sur la Terre.
\item Dans un régime pseudo-périodique, l'amplitude des oscillations reste constante au cours du temps.
\item L'énergie mécanique d'un oscillateur harmonique non amorti est proportionnelle au carré de l'amplitude.
\end{enumerate}
\end{exercice}

\begin{exercice}[Définitions et vocabulaire]
\begin{enumerate}
\item Définir les grandeurs suivantes et donner leur unité SI : amplitude $X_m$, pulsation propre $\omega_0$, phase à l'origine $\varphi$, période propre $T_0$.
\item Qu'appelle-t-on \cle{régime pseudo-périodique} ? \cle{Régime apériodique} ? Citer un exemple de chaque dans la vie courante.
\item Qu'appelle-t-on \cle{résonance mécanique} ? Donner la condition sur les fréquences de l'excitateur et du résonateur.
\end{enumerate}
\end{exercice}

\begin{exercice}[Calculs directs]
\begin{enumerate}
\item Un solide de masse $m = \SI{200}{g}$ est accroché à un ressort de raideur $k = \SI{50}{N.m^{-1}}$. Calculer la période propre $T_0$ et la fréquence propre $f_0$ de l'oscillateur.
\item L'élongation d'un oscillateur est $x(t) = 4,0\cos\left(10\pi t + \dfrac{\pi}{3}\right)$ avec $x$ en cm et $t$ en s. Déterminer l'amplitude, la pulsation, la période et la phase à l'origine. Calculer l'élongation et la vitesse à $t = 0$.
\item Un pendule simple bat la seconde ($T_0 = \SI{2,0}{s}$). Calculer sa longueur (prendre $g = \SI{9,8}{m.s^{-2}}$).
\end{enumerate}
\end{exercice}

\courssub{Je m'entraîne}

\begin{exercice}[Pendule élastique vertical et pesanteur]
Un ressort de raideur $k = \SI{25}{N.m^{-1}}$, suspendu verticalement, porte une masse $m = \SI{250}{g}$. À l'équilibre, l'allongement du ressort est $\Delta\ell_0 = \SI{9,8}{cm}$.
\begin{enumerate}
\item Faire le bilan des forces à l'équilibre et en déduire la valeur de $g$ au lieu de l'expérience.
\item On écarte la masse de sa position d'équilibre et on la lâche. Établir l'équation différentielle du mouvement par rapport à la position d'équilibre.
\item Montrer que la période propre peut s'écrire $T_0 = 2\pi\sqrt{\dfrac{\Delta\ell_0}{g}}$ et conclure : la période dépend-elle de $g$ séparément de $\Delta\ell_0$ ? Calculer $T_0$.
\end{enumerate}
\end{exercice}

\begin{exercice}[Exploitation d'un enregistrement d'oscillations amorties]
Un mobile de masse $m = \SI{500}{g}$ accroché à un ressort horizontal oscille dans un liquide. L'enregistrement de $x(t)$ donne les résultats suivants : 10 pseudo-périodes s'écoulent en $\SI{6,3}{s}$ ; l'amplitude passe de $X_1 = \SI{6,0}{cm}$ à la première crête à $X_2 = \SI{4,2}{cm}$ à la crête suivante.
\begin{enumerate}
\item Déterminer la pseudo-période $T$ des oscillations.
\item En assimilant la pseudo-période à la période propre, calculer la raideur $k$ du ressort.
\item Calculer le rapport $\dfrac{X_2}{X_1}$. Ce rapport est-il compatible avec un amortissement faible ? Identifier le régime d'oscillations.
\item L'énergie mécanique est-elle conservée ? Justifier et préciser ce qu'elle devient.
\end{enumerate}
\end{exercice}

\begin{exercice}[Pendule de torsion]
Un disque homogène est suspendu par son centre à un fil de torsion de constante $C = \SI{2,5e-2}{N.m.rad^{-1}}$. Écarté de sa position d'équilibre puis abandonné, il effectue des oscillations de période propre $T_0 = \SI{1,2}{s}$.
\begin{enumerate}
\item Rappeler l'expression de la période propre d'un pendule de torsion et calculer le moment d'inertie $J$ du disque par rapport à l'axe du fil.
\item On fixe sur le disque deux surcharges identiques qui portent le moment d'inertie total à $J' = 1,44\,J$. Calculer la nouvelle période propre $T_0'$.
\item Le disque a pour rayon $R = \SI{10}{cm}$. En admettant que $J = \dfrac{1}{2}MR^2$, en déduire la masse $M$ du disque.
\end{enumerate}
\end{exercice}

\begin{exercice}[Suspension d'une voiture sur une piste camerounaise]
Une voiture de masse $m = \SI{800}{kg}$, à l'arrêt, s'enfonce de $d = \SI{5,0}{cm}$ sur ses suspensions lorsqu'elle est chargée. On modélise l'ensemble des suspensions par un ressort unique de raideur $k$ associé à un amortisseur. On prend $g = \SI{9,8}{m.s^{-2}}$.
\begin{enumerate}
\item Écrire la condition d'équilibre et calculer la raideur équivalente $k$.
\item En déduire la période propre $T_0$ puis la fréquence propre $f_0$ de la caisse sur ses suspensions.
\item La voiture roule sur une piste ondulée dont les bosses sont espacées de $L = \SI{12}{m}$. À quelle vitesse $v$ la fréquence des sollicitations est-elle égale à $f_0$ ? Quel phénomène se produit alors et pourquoi est-il dangereux ?
\item Quel est le rôle des amortisseurs ? Que risque-t-on si l'amortissement est trop fort ?
\end{enumerate}
\end{exercice}

\courssub{Je me prépare au Bac}

\begin{exercice}[Type Bac : oscillateur harmonique horizontal]
\textit{(D'après Baccalauréat C, session de juin 2018)}

Un solide $S$ de masse $m = \SI{100}{g}$, assimilable à un point matériel $G$, peut glisser sans frottement sur une tige horizontale. Il est accroché à un ressort à spires non jointives, de masse négligeable et de raideur $k = \SI{40}{N.m^{-1}}$, dont l'autre extrémité est fixe. On repère la position de $G$ par son abscisse $x$ sur l'axe horizontal $(Ox)$, l'origine $O$ correspondant à la position d'équilibre. On écarte $S$ de $X_m = \SI{3,0}{cm}$ dans le sens positif et on le lâche sans vitesse initiale à l'instant $t = 0$. On prend $\pi^2 = 10$.
\begin{enumerate}
\item Faire le bilan des forces appliquées à $S$ à un instant quelconque et les représenter sur un schéma.
\item En appliquant la relation fondamentale de la dynamique, établir l'équation différentielle du mouvement de $G$ et montrer qu'elle s'écrit $\ddot{x} + \omega_0^2\,x = 0$. Calculer $\omega_0$.
\item Vérifier que $x(t) = X_m\cos(\omega_0 t + \varphi)$ est solution de cette équation, puis déterminer $\varphi$ à l'aide des conditions initiales. Écrire l'expression complète de $x(t)$.
\item Calculer la période propre $T_0$ de l'oscillateur.
\item Donner l'expression de l'énergie mécanique $E_m$ du système \{solide + ressort\} en fonction de $k$ et $X_m$, puis calculer sa valeur.
\item En déduire la vitesse maximale $v_m$ du solide. À quel(s) point(s) de la trajectoire est-elle atteinte ?
\item Calculer l'énergie cinétique et l'énergie potentielle élastique lorsque $x = \SI{1,5}{cm}$.
\end{enumerate}
\end{exercice}

\begin{exercice}[Type Bac : mesure de $g$ à Yaoundé avec un pendule simple]
\textit{(D'après Baccalauréat C, session de juin 2016)}

Un groupe d'élèves de Terminale C du lycée de Yaoundé réalise un pendule simple constitué d'une petite bille métallique de masse $m = \SI{50}{g}$ suspendue à un fil inextensible de masse négligeable et de longueur réglable $\ell$. Pour différentes valeurs de $\ell$, ils mesurent la durée de 20 oscillations de faible amplitude et en déduisent la période $T_0$. Les résultats sont consignés dans le tableau ci-dessous.

\begin{center}
\begin{tabular}{|l|*{5}{c|}}
\hline
$\ell$ (m) & 0,20 & 0,40 & 0,60 & 0,80 & 1,00 \\
\hline
$T_0$ (s) & 0,90 & 1,27 & 1,55 & 1,80 & 2,01 \\
\hline
$T_0^2$ (s$^2$) & & & & & \\
\hline
\end{tabular}
\end{center}

\begin{enumerate}
\item Rappeler la condition d'application du modèle du pendule simple (petites oscillations) et l'expression de sa période propre.
\item Recopier et compléter la troisième ligne du tableau.
\item Montrer que $T_0^2$ est une fonction affine de $\ell$ et donner l'expression théorique de son coefficient directeur.
\item Tracer sur papier millimétré la courbe $T_0^2 = f(\ell)$ (échelles : $\SI{1}{cm}$ pour $\SI{0,10}{m}$ en abscisse ; $\SI{1}{cm}$ pour $\SI{0,40}{s^2}$ en ordonnée). Conclure quant à la nature de la courbe.
\item Déterminer graphiquement le coefficient directeur $a$ de la droite obtenue, puis en déduire la valeur expérimentale de l'intensité de pesanteur $g$ à Yaoundé. On prendra $\pi^2 = 9,87$.
\item Comparer à la valeur de référence $g_0 = \SI{9,81}{m.s^{-2}}$ en calculant l'écart relatif en pourcentage. Citer deux sources d'erreur possibles.
\item Le même pendule, réglé à $\ell = \SI{1,00}{m}$, est transporté à Limbé au bord de l'océan. Sa période change-t-elle sensiblement ? Justifier.
\end{enumerate}
\end{exercice}

\begin{exercice}[Type Bac : synthèse — oscillations libres, amorties et forcées]
\textit{(D'après Baccalauréat C, session de juin 2019)}

\textbf{Partie A : oscillations libres non amorties.} Un jouet d'enfant est constitué d'une figurine de masse $m = \SI{80}{g}$ fixée à l'extrémité d'un ressort vertical de raideur $k = \SI{20}{N.m^{-1}}$. On néglige d'abord les frottements. On prend $g = \SI{10}{m.s^{-2}}$ et $\pi^2 = 10$.
\begin{enumerate}
\item Calculer l'allongement $\Delta\ell_0$ du ressort à l'équilibre.
\item On tire la figurine vers le bas de $\SI{2,0}{cm}$ et on la lâche sans vitesse. Établir l'équation différentielle du mouvement autour de la position d'équilibre et calculer la période propre $T_0$.
\item Calculer l'énergie mécanique de l'oscillateur et la vitesse de la figurine au passage par la position d'équilibre.
\end{enumerate}

\textbf{Partie B : oscillations amorties.} En réalité, les frottements de l'air amortissent le mouvement. Un enregistrement montre que l'amplitude est divisée par 2 au bout de 5 pseudo-oscillations.
\begin{enumerate}
\setcounter{enumi}{3}
\item Nommer le régime observé. Représenter qualitativement l'allure de la courbe $x(t)$.
\item Expliquer, d'un point de vue énergétique, l'évolution de l'amplitude au cours du temps.
\end{enumerate}

\textbf{Partie C : oscillations forcées.} Le jouet est posé sur la table d'une machine à laver en fonctionnement, qui vibre à la fréquence $f_e$ réglable.
\begin{enumerate}
\setcounter{enumi}{5}
\item Identifier l'excitateur et le résonateur dans cette expérience.
\item Calculer la fréquence propre $f_0$ du jouet. Pour quelle valeur de $f_e$ observe-t-on la résonance ? Décrire alors le comportement de la figurine.
\item Expliquer pourquoi les ponts et les immeubles doivent être conçus en tenant compte du phénomène de résonance (vent, séismes, passages de camions sur les ponts du Wouri ou de la Sanaga).
\end{enumerate}
\end{exercice}

\newpage

\coursec{Corrigés des exercices}

\coursec{Exercices corrigés — Oscillateurs mécaniques}

\begin{corrige}[Exercice 1 — QCM : une seule réponse exacte]
\begin{enumerate}
\item \textbf{Réponse b.} L'analyse dimensionnelle impose $T_0$ homogène à un temps : $\left[\sqrt{m/k}\right] = \sqrt{\dfrac{\mathrm{kg}}{\mathrm{kg.s^{-2}}}} = \mathrm{s}$. La réponse a est l'inverse, la c est la fréquence propre $f_0$, la d n'est pas homogène à un temps.
\item \textbf{Réponse b.} $T_0 = 2\pi\sqrt{m/k}$ : si $m' = 4m$, alors $T_0' = 2\pi\sqrt{4m/k} = 2\,T_0$. La période varie comme la \emph{racine carrée} de la masse.
\item \textbf{Réponse c.} $T_0 = 2\pi\sqrt{\ell/g} = 2\pi\sqrt{\dfrac{1{,}0}{9{,}8}} \simeq 2\pi \times 0{,}319 \simeq \SI{2,0}{s}$. C'est le célèbre pendule qui « bat la seconde » (une demi-période vaut \SI{1}{s}).
\item \textbf{Réponse c.} À la résonance, la fréquence de l'excitateur est voisine de la fréquence propre du résonateur et l'amplitude des oscillations passe par un \cle{maximum}, d'autant plus aigu que l'amortissement est faible.
\end{enumerate}
\end{corrige}

\begin{corrige}[Exercice 2 — Vrai ou faux ? Justifier]
\begin{enumerate}
\item \textbf{Vrai.} Dérivons : $\dot{x} = -\omega_0 X_m\sin(\omega_0 t+\varphi)$ puis $\ddot{x} = -\omega_0^2 X_m\cos(\omega_0 t+\varphi) = -\omega_0^2 x$, donc $\ddot{x}+\omega_0^2 x = 0$ pour \emph{toutes} les valeurs des constantes $X_m$ et $\varphi$. Ce sont les conditions initiales qui fixent $X_m$ et $\varphi$.
\item \textbf{Faux.} $T_0 = 2\pi\sqrt{m/k}$ ne fait intervenir que $m$ et $k$, caractéristiques du système : c'est l'\cle{isochronisme} des oscillations (dans le cadre du modèle harmonique).
\item \textbf{Vrai.} $T_{0,\text{L}} = 2\pi\sqrt{\ell/g_{\text{L}}} = 2\pi\sqrt{6\ell/g} = \sqrt{6}\,T_{0,\text{T}} \simeq 2{,}45\,T_{0,\text{T}}$ : la période est plus grande, le pendule oscille donc plus lentement.
\item \textbf{Faux.} En régime pseudo-périodique, l'amplitude \emph{décroît} au cours du temps à cause de la dissipation d'énergie par les frottements ; seule la pseudo-période reste sensiblement constante.
\item \textbf{Vrai.} $E_m = \dfrac{1}{2}kX_m^2$ : l'énergie mécanique est bien proportionnelle à $X_m^2$ (et elle se conserve en l'absence d'amortissement).
\end{enumerate}
\end{corrige}

\begin{corrige}[Exercice 3 — Définitions et vocabulaire]
\begin{enumerate}
\item \begin{itemize}
\item \cle{Amplitude} $X_m$ : valeur maximale de l'élongation $x(t)$ ; unité SI : le mètre (\si{m}).
\item \cle{Pulsation propre} $\omega_0$ : vitesse angulaire caractéristique de l'oscillateur libre, $\omega_0 = \dfrac{2\pi}{T_0}$ ; unité SI : \si{rad.s^{-1}}.
\item \cle{Phase à l'origine} $\varphi$ : valeur à $t=0$ de la phase $\omega_0 t + \varphi$ ; elle dépend des conditions initiales ; unité : radian (\si{rad}).
\item \cle{Période propre} $T_0$ : durée d'une oscillation complète de l'oscillateur libre non amorti ; unité SI : la seconde (\si{s}).
\end{itemize}
\item \cle{Régime pseudo-périodique} : amortissement faible, le système oscille avec une amplitude décroissante et une pseudo-période $T$ proche de $T_0$. Exemple : masse suspendue à un ressort plongée dans l'air, balançoire abandonnée. \cle{Régime apériodique} : amortissement fort, le système revient vers sa position d'équilibre \emph{sans osciller}. Exemple : portière de voiture avec amortisseur, masse plongée dans un liquide très visqueux (huile, miel).
\item La \cle{résonance mécanique} est le phénomène d'augmentation brutale de l'amplitude d'un oscillateur (résonateur) soumis à une excitation périodique (excitateur) lorsque la fréquence $f_e$ de l'excitateur devient voisine de la fréquence propre $f_0$ du résonateur : condition $f_e \simeq f_0$.
\end{enumerate}
\end{corrige}

\begin{corrige}[Exercice 4 — Calculs directs]
\begin{enumerate}
\item Conversion : $m = \SI{200}{g} = \SI{0,200}{kg}$.
\[ T_0 = 2\pi\sqrt{\dfrac{m}{k}} = 2\pi\sqrt{\dfrac{0{,}200}{50}} = 2\pi\sqrt{4{,}0\times 10^{-3}} \simeq \SI{0,40}{s} \qquad ; \qquad f_0 = \dfrac{1}{T_0} \simeq \dfrac{1}{0{,}397} \simeq \SI{2,5}{Hz}. \]
\item Par identification avec $x(t) = X_m\cos(\omega_0 t + \varphi)$ :
\[ X_m = \SI{4,0}{cm} \quad ; \quad \omega_0 = 10\pi \simeq \SI{31,4}{rad.s^{-1}} \quad ; \quad T_0 = \dfrac{2\pi}{\omega_0} = \SI{0,20}{s} \quad ; \quad \varphi = \dfrac{\pi}{3}\ \si{rad}. \]
À $t=0$ : $x(0) = 4{,}0\cos\dfrac{\pi}{3} = 4{,}0\times 0{,}5 = \SI{2,0}{cm}$.
Vitesse : $v(t) = \dot{x} = -40\pi\sin\left(10\pi t + \dfrac{\pi}{3}\right)$ (en \si{cm.s^{-1}}), donc
\[ v(0) = -40\pi\sin\dfrac{\pi}{3} = -40\pi\times\dfrac{\sqrt{3}}{2} \simeq \SI{-1,1e2}{cm.s^{-1}} = \SI{-1,1}{m.s^{-1}}. \]
Le signe négatif signifie que le mobile se déplace d'abord dans le sens des $x$ décroissants.
\item De $T_0 = 2\pi\sqrt{\ell/g}$ on tire $\ell = \dfrac{gT_0^2}{4\pi^2} = \dfrac{9{,}8\times (2{,}0)^2}{4\times 9{,}87} = \dfrac{39{,}2}{39{,}48} \simeq \SI{0,99}{m}$, soit environ \SI{1,0}{m}.
\end{enumerate}
\end{corrige}

\begin{corrige}[Exercice 5 — Pendule élastique vertical et pesanteur]
\begin{enumerate}
\item À l'équilibre, la masse est soumise à son poids $\vec{P}$ (vertical, vers le bas, $P = mg$) et à la tension du ressort $\vec{T}$ (verticale, vers le haut, $T = k\,\Delta\ell_0$). L'équilibre impose $\vec{P}+\vec{T}=\vec{0}$, soit $mg = k\,\Delta\ell_0$, d'où
\[ g = \dfrac{k\,\Delta\ell_0}{m} = \dfrac{25\times 9{,}8\times 10^{-2}}{0{,}250} = \dfrac{2{,}45}{0{,}250} = \SI{9,8}{m.s^{-2}}. \]
\item Choisissons l'axe $(Ox)$ vertical descendant, d'origine la position d'équilibre $O$. À l'abscisse $x$, l'allongement total du ressort est $\Delta\ell_0 + x$. La RFD projetée sur $(Ox)$ s'écrit :
\[ m\ddot{x} = mg - k(\Delta\ell_0 + x) = \underbrace{mg - k\Delta\ell_0}_{=\,0\ \text{(équilibre)}} - kx \quad\Longrightarrow\quad \boxed{\ddot{x} + \dfrac{k}{m}\,x = 0}. \]
C'est l'équation d'un oscillateur harmonique de pulsation propre $\omega_0 = \sqrt{k/m}$ : la pesanteur ne modifie pas la nature du mouvement, elle décale seulement la position d'équilibre.
\item Comme $k\,\Delta\ell_0 = mg$, on a $\dfrac{m}{k} = \dfrac{\Delta\ell_0}{g}$, donc
\[ T_0 = 2\pi\sqrt{\dfrac{m}{k}} = 2\pi\sqrt{\dfrac{\Delta\ell_0}{g}}. \]
\textbf{Conclusion :} $T_0$ ne dépend pas de $g$ et de $\Delta\ell_0$ séparément, mais uniquement de leur rapport $\Delta\ell_0/g = m/k$ : deux lieux de pesanteur différente donnent des allongements différents mais la même période, à masse et ressort identiques. Numériquement :
\[ T_0 = 2\pi\sqrt{\dfrac{0{,}098}{9{,}8}} = 2\pi\sqrt{10^{-2}} = 2\pi\times 0{,}10 \simeq \SI{0,63}{s}. \]
\end{enumerate}
\end{corrige}

\begin{corrige}[Exercice 6 — Exploitation d'un enregistrement d'oscillations amorties]
\begin{enumerate}
\item La pseudo-période est la durée moyenne d'une oscillation :
\[ T = \dfrac{6{,}3}{10} = \SI{0,63}{s}. \]
\item En assimilant $T \simeq T_0 = 2\pi\sqrt{m/k}$, on tire $k = \dfrac{4\pi^2 m}{T^2}$ :
\[ k = \dfrac{4\times 9{,}87\times 0{,}500}{(0{,}63)^2} = \dfrac{19{,}74}{0{,}397} \simeq \SI{50}{N.m^{-1}}. \]
\item $\dfrac{X_2}{X_1} = \dfrac{4{,}2}{6{,}0} = 0{,}70$. L'amplitude ne décroît que de $30\,\%$ par pseudo-période et le système effectue encore de nombreuses oscillations : l'amortissement est \textbf{faible}, le régime est \cle{pseudo-périodique} (compatible avec l'approximation $T \simeq T_0$ faite à la question 2).
\item L'énergie mécanique $E_m = \dfrac{1}{2}kX_m^2$ \textbf{n'est pas conservée} : elle diminue à chaque oscillation (d'un facteur $0{,}7^2 \simeq 0{,}49$ par pseudo-période). Elle est dissipée par les frottements du liquide, c'est-à-dire convertie en \cle{énergie thermique} (échauffement du liquide et du ressort).
\end{enumerate}
\end{corrige}

\begin{corrige}[Exercice 7 — Pendule de torsion]
\begin{enumerate}
\item La période propre d'un pendule de torsion est $T_0 = 2\pi\sqrt{\dfrac{J}{C}}$, d'où
\[ J = \dfrac{C\,T_0^2}{4\pi^2} = \dfrac{2{,}5\times 10^{-2}\times (1{,}2)^2}{4\times 9{,}87} = \dfrac{3{,}6\times 10^{-2}}{39{,}48} \simeq \SI{9,1e-4}{kg.m^2}. \]
Vérification dimensionnelle : $[\mathrm{N.m.rad^{-1}.s^2}] = \mathrm{kg.m^2.s^{-2}.s^2} = \mathrm{kg.m^2}$. \checkmark
\item $T_0' = 2\pi\sqrt{\dfrac{J'}{C}} = 2\pi\sqrt{\dfrac{1{,}44\,J}{C}} = \sqrt{1{,}44}\;T_0 = 1{,}2\times 1{,}2 = \SI{1,44}{s}$.
\item De $J = \dfrac{1}{2}MR^2$ on tire
\[ M = \dfrac{2J}{R^2} = \dfrac{2\times 9{,}1\times 10^{-4}}{(0{,}10)^2} = \dfrac{1{,}82\times 10^{-3}}{10^{-2}} \simeq \SI{0,18}{kg} \quad (\text{environ } \SI{180}{g}). \]
\end{enumerate}
\end{corrige}

\begin{corrige}[Exercice 8 — Suspension d'une voiture sur une piste camerounaise]
\begin{enumerate}
\item À l'équilibre, le poids de la caisse est compensé par la force de rappel des suspensions : $mg = kd$, d'où
\[ k = \dfrac{mg}{d} = \dfrac{800\times 9{,}8}{5{,}0\times 10^{-2}} = \dfrac{7840}{0{,}050} \simeq \SI{1,6e5}{N.m^{-1}}. \]
\item $T_0 = 2\pi\sqrt{\dfrac{m}{k}} = 2\pi\sqrt{\dfrac{d}{g}} = 2\pi\sqrt{\dfrac{0{,}050}{9{,}8}} = 2\pi\times 0{,}0714 \simeq \SI{0,45}{s}$ ; \quad $f_0 = \dfrac{1}{T_0} \simeq \SI{2,2}{Hz}$.
\item Les bosses sollicitent la suspension à la fréquence $f_e = \dfrac{v}{L}$. La résonance survient pour $f_e = f_0$, soit
\[ v = f_0\,L = 2{,}23\times 12 \simeq \SI{27}{m.s^{-1}} \simeq \SI{96}{km.h^{-1}}. \]
C'est le phénomène de \cle{résonance mécanique} : l'amplitude de pompage de la caisse devient maximale, la tenue de route se dégrade (perte d'adhérence, inconfort, risque d'accident ou de rupture des suspensions).
\item Les \cle{amortisseurs} dissipent l'énergie mécanique en chaleur : ils réduisent l'amplitude à la résonance et ramènent rapidement la caisse à l'équilibre après une bosse (régime proche du critique). Si l'amortissement est \emph{trop fort} (régime apériodique raide), la suspension ne « travaille » plus : les chocs sont transmis directement à la caisse, d'où inconfort et perte de contact des roues avec la piste.
\end{enumerate}
\end{corrige}

\begin{corrige}[Exercice 9 — Type Bac : oscillateur harmonique horizontal]
\textbf{Barème indicatif (sur 10 pts) :} 1) 1 pt ; 2) 2 pts ; 3) 2 pts ; 4) 1 pt ; 5) 1 pt ; 6) 1,5 pt ; 7) 1,5 pt.

\begin{enumerate}
\item Le solide $S$ est soumis à trois forces :
\begin{itemize}
\item son \cle{poids} $\vec{P}$, vertical, vers le bas, $P = mg$ ;
\item la \cle{réaction} $\vec{R}$ de la tige, verticale, vers le haut (pas de frottement : $\vec{R}\perp$ tige) ;
\item la \cle{force de rappel} du ressort $\vec{F} = -kx\,\vec{\imath}$, horizontale, dirigée vers $O$.
\end{itemize}
\begin{popfigure}
\begin{tikzpicture}[scale=1.1]
\draw[thick] (-0.5,0) -- (5.5,0);
\draw[decorate,decoration={coil,segment length=4pt,amplitude=3pt}] (0,0.35) -- (2.2,0.35);
\draw[fill=popblueL,draw=popblue] (2.2,0.05) rectangle (3.0,0.65);
\fill (2.6,0.35) circle (1.5pt) node[above=2pt] {$G$};
\draw[-{Stealth},very thick,popdark] (2.6,0.35) -- (2.6,-0.55) node[below] {$\vec{P}$};
\draw[-{Stealth},very thick,popgreen] (2.6,0.65) -- (2.6,1.55) node[above] {$\vec{R}$};
\draw[-{Stealth},very thick,poporange] (2.2,0.35) -- (1.0,0.35) node[above left] {$\vec{F}$};
\draw[-{Stealth}] (3.4,0) -- (4.2,0) node[right] {$x$};
\fill (0,0) circle (1.5pt) node[below left] {$O$};
\end{tikzpicture}
\legende{Schéma des forces agissant sur le solide $S$ (pendule élastique horizontal).}
\end{popfigure}
\item RFD appliquée à $S$ : $\vec{P}+\vec{R}+\vec{F} = m\vec{a}$. En projection sur $(Ox)$, $\vec{P}$ et $\vec{R}$ n'ont pas de composante horizontale :
\[ m\ddot{x} = -kx \quad\Longrightarrow\quad \ddot{x} + \dfrac{k}{m}\,x = 0, \]
soit $\ddot{x}+\omega_0^2 x = 0$ avec $\omega_0 = \sqrt{\dfrac{k}{m}} = \sqrt{\dfrac{40}{0{,}100}} = \sqrt{400} = \SI{20}{rad.s^{-1}}$.
\item Soit $x(t) = X_m\cos(\omega_0 t+\varphi)$. Alors $\ddot{x} = -\omega_0^2 X_m\cos(\omega_0 t+\varphi) = -\omega_0^2 x$, donc $\ddot{x}+\omega_0^2 x = 0$ : c'est bien une solution quelles que soient $X_m$ et $\varphi$.
Conditions initiales : $x(0) = X_m\cos\varphi = X_m$ (lâché de $+X_m$) donne $\cos\varphi = 1$ ; $\dot{x}(0) = -\omega_0 X_m\sin\varphi = 0$ donne $\sin\varphi = 0$. D'où $\varphi = 0$ et
\[ \boxed{x(t) = 3{,}0\cos(20\,t)\ \text{(en cm, } t \text{ en s)}}. \]
\item $T_0 = \dfrac{2\pi}{\omega_0} = \dfrac{2\pi}{20} = \dfrac{\pi}{10} \simeq \SI{0,31}{s}$ (avec $\pi^2 = 10$, $T_0 = 2\pi\sqrt{m/k} = 2\pi\times 0{,}05$).
\item L'énergie mécanique se conserve (pas de frottement) et vaut l'énergie potentielle élastique maximale :
\[ E_m = \dfrac{1}{2}kX_m^2 = \dfrac{1}{2}\times 40\times (3{,}0\times 10^{-2})^2 = 20\times 9{,}0\times 10^{-4} = \SI{1,8e-2}{J}. \]
\item Conservation de l'énergie : $\dfrac{1}{2}mv_m^2 = E_m$, d'où
\[ v_m = \sqrt{\dfrac{2E_m}{m}} = \sqrt{\dfrac{3{,}6\times 10^{-2}}{0{,}100}} = \sqrt{0{,}36} = \SI{0,60}{m.s^{-1}} \quad (\text{on vérifie } v_m = \omega_0 X_m = 20\times 0{,}030). \]
Elle est atteinte au passage par la \textbf{position d'équilibre} $x = 0$, où l'énergie potentielle élastique est nulle.
\item Pour $x = \SI{1,5}{cm} = \SI{1,5e-2}{m}$ :
\[ E_{pe} = \dfrac{1}{2}kx^2 = \dfrac{1}{2}\times 40\times (1{,}5\times 10^{-2})^2 = 20\times 2{,}25\times 10^{-4} = \SI{4,5e-3}{J}, \]
\[ E_c = E_m - E_{pe} = 1{,}8\times 10^{-2} - 4{,}5\times 10^{-3} = \SI{1,35e-2}{J}. \]
On remarque que $x = X_m/2$ donne $E_{pe} = E_m/4$ et $E_c = 3E_m/4$.
\end{enumerate}
\textbf{Points de vigilance :} convertir la masse en kg et les élongations en mètres \emph{avant} tout calcul ; justifier que $\vec{P}$ et $\vec{R}$ se compensent (ou n'ont pas de projection sur $(Ox)$) ; pour déterminer $\varphi$, utiliser \emph{les deux} conditions initiales (position \textbf{et} vitesse) ; citer la conservation de $E_m$ avant de l'exploiter.
\end{corrige}

\begin{corrige}[Exercice 10 — Type Bac : mesure de $g$ à Yaoundé avec un pendule simple]
\textbf{Barème indicatif (sur 10 pts) :} 1) 1 pt ; 2) 1 pt ; 3) 1,5 pt ; 4) 2 pts ; 5) 2 pts ; 6) 1,5 pt ; 7) 1 pt.
\begin{enumerate}
\item Le modèle du pendule simple exige des oscillations de \textbf{faible amplitude} ($\alpha_m \lesssim 10^{\circ}$, de sorte que $\sin\alpha \simeq \alpha$ en radians), un fil inextensible de masse négligeable et une bille de dimensions petites devant $\ell$. Alors
\[ T_0 = 2\pi\sqrt{\dfrac{\ell}{g}}. \]
\item En élevant au carré les valeurs de $T_0$ :
\begin{center}
\begin{tabular}{|l|*{5}{c|}}
\hline
$\ell$ (m) & 0,20 & 0,40 & 0,60 & 0,80 & 1,00 \\
\hline
$T_0$ (s) & 0,90 & 1,27 & 1,55 & 1,80 & 2,01 \\
\hline
$T_0^2$ (\si{s^2}) & 0,81 & 1,61 & 2,40 & 3,24 & 4,04 \\
\hline
\end{tabular}
\end{center}
\item De $T_0 = 2\pi\sqrt{\ell/g}$ on tire $T_0^2 = \dfrac{4\pi^2}{g}\,\ell$ : $T_0^2$ est une fonction \textbf{affine} (en fait linéaire) de $\ell$, de coefficient directeur théorique
\[ a_{\text{th}} = \dfrac{4\pi^2}{g}. \]
\item Le tracé (non reproduit ici) place les points $(0{,}20\,;\,0{,}81)$, $(0{,}40\,;\,1{,}61)$, $(0{,}60\,;\,2{,}40)$, $(0{,}80\,;\,3{,}24)$, $(1{,}00\,;\,4{,}04)$. Les points sont \textbf{alignés sur une droite passant par l'origine} : cela confirme la proportionnalité $T_0^2 \propto \ell$, donc la validité du modèle du pendule simple.
\item En prenant deux points éloignés de la droite moyenne, par exemple $(0\,;\,0)$ et $(1{,}00\,;\,4{,}03)$ :
\[ a = \dfrac{\Delta T_0^2}{\Delta \ell} = \dfrac{4{,}03 - 0}{1{,}00 - 0} \simeq \SI{4,0}{s^2.m^{-1}}. \]
Comme $a = \dfrac{4\pi^2}{g}$, on en déduit
\[ g = \dfrac{4\pi^2}{a} = \dfrac{4\times 9{,}87}{4{,}03} = \dfrac{39{,}48}{4{,}03} \simeq \SI{9,8}{m.s^{-2}}. \]
\item Écart relatif :
\[ \dfrac{|g - g_0|}{g_0} = \dfrac{|9{,}80 - 9{,}81|}{9{,}81} \simeq 0{,}001 \simeq 0{,}1\,\%. \]
Accord excellent. Sources d'erreur possibles : imprécision du chronométrage (temps de réaction de l'opérateur), mesure de $\ell$ (position exacte du centre de la bille), amplitude angulaire trop grande, frottements de l'air.
\item La période ne dépend que de $\ell$ et de $g$. Or $g$ varie très peu entre Yaoundé (altitude $\sim\SI{750}{m}$) et Limbé (niveau de la mer) : l'écart relatif sur $g$ est de l'ordre de $10^{-4}$, donc celui sur $T_0$ de l'ordre de $5\times 10^{-5}$, totalement \textbf{insensible} avec le matériel utilisé. La période reste pratiquement égale à $\SI{2,01}{s}$.
\end{enumerate}
\textbf{Points de vigilance :} le coefficient directeur se calcule sur la \emph{droite moyenne}, pas entre deux points expérimentaux ; ne pas oublier que $T_0^2 = \dfrac{4\pi^2}{g}\ell$ impose une droite \emph{passant par l'origine} ; citer explicitement la condition des petites oscillations ($\sin\alpha \simeq \alpha$) dès la question 1.
\end{corrige}

\begin{corrige}[Exercice 11 — Type Bac : synthèse — oscillations libres, amorties et forcées]
\textbf{Barème indicatif (sur 10 pts) :} Partie A : 4 pts ; Partie B : 2,5 pts ; Partie C : 3,5 pts.

\textbf{Partie A :}
\begin{enumerate}
\item À l'équilibre : $k\,\Delta\ell_0 = mg$, d'où
\[ \Delta\ell_0 = \dfrac{mg}{k} = \dfrac{0{,}080\times 10}{20} = \SI{4,0e-2}{m} = \SI{4,0}{cm}. \]
\item En repérant la position par $x$ à partir de la position d'équilibre (axe vertical descendant), la RFD donne, comme pour tout pendule élastique vertical :
\[ m\ddot{x} = mg - k(\Delta\ell_0 + x) = -kx \quad\Longrightarrow\quad \ddot{x} + \dfrac{k}{m}\,x = 0. \]
\[ T_0 = 2\pi\sqrt{\dfrac{m}{k}} \quad ; \quad T_0^2 = \dfrac{4\pi^2 m}{k} = \dfrac{4\times 10\times 0{,}080}{20} = 0{,}16\ \si{s^2} \quad\Longrightarrow\quad T_0 = \SI{0,40}{s}. \]
\item L'amplitude est $X_m = \SI{2,0}{cm} = \SI{2,0e-2}{m}$ (lâché sans vitesse). D'où
\[ E_m = \dfrac{1}{2}kX_m^2 = \dfrac{1}{2}\times 20\times (2{,}0\times 10^{-2})^2 = 10\times 4{,}0\times 10^{-4} = \SI{4,0e-3}{J}. \]
Au passage par la position d'équilibre, toute l'énergie est cinétique : $\dfrac{1}{2}mv_m^2 = E_m$, soit
\[ v_m = \sqrt{\dfrac{2E_m}{m}} = \sqrt{\dfrac{8{,}0\times 10^{-3}}{0{,}080}} = \sqrt{0{,}10} \simeq \SI{0,32}{m.s^{-1}}. \]
\end{enumerate}
\textbf{Partie B :}
\begin{enumerate}
\setcounter{enumi}{3}
\item L'amplitude décroît progressivement (divisée par 2 en 5 oscillations) : c'est le régime \cle{pseudo-périodique} (amortissement faible).
\begin{popfigure}
\begin{tikzpicture}
\begin{axis}[width=0.75\linewidth,height=5cm,axis lines=middle,xlabel={$t$},ylabel={$x$},xtick=\empty,ytick=\empty,xmin=0,xmax=6.5,ymin=-1.4,ymax=1.4]
\addplot[popcurve,domain=0:6.3,samples=300]{exp(-0.28*x)*cos(deg(2*pi*x/0.4))};
\addplot[poporange,dashed,domain=0:6.3,samples=100]{exp(-0.28*x)};
\addplot[poporange,dashed,domain=0:6.3,samples=100]{-exp(-0.28*x)};
\end{axis}
\end{tikzpicture}
\legende{Oscillations pseudo-périodiques : amplitude décroissante, pseudo-période $T \simeq T_0 = 0,40\,\text{s}$. L'enveloppe exponentielle (pointillés) montre la décroissance de l'amplitude.}
\end{popfigure}
\item Les frottements de l'air effectuent un travail résistant : l'énergie mécanique $E_m = \dfrac{1}{2}kX_m^2$ diminue continûment, dissipée en \cle{chaleur}. Comme $E_m \propto X_m^2$, l'amplitude décroît au cours du temps jusqu'à l'arrêt de la figurine en position d'équilibre.
\end{enumerate}
\textbf{Partie C :}
\begin{enumerate}
\setcounter{enumi}{5}
\item L'\cle{excitateur} est la machine à laver (elle impose des vibrations de fréquence $f_e$) ; le \cle{résonateur} est le jouet \{figurine + ressort\}.
\item $f_0 = \dfrac{1}{T_0} = \dfrac{1}{0{,}40} = \SI{2,5}{Hz}$. La résonance est observée pour $f_e \simeq f_0 = \SI{2,5}{Hz}$ : l'amplitude des oscillations de la figurine devient alors \textbf{maximale} — la figurine « saute » vigoureusement, bien que l'excitation soit faible.
\item Toute structure (pont, immeuble) possède des fréquences propres de vibration. Si une excitation périodique (rafales de vent, ondes sismiques, passage cadencé de camions sur les ponts du Wouri ou de la Sanaga) a une fréquence voisine d'une fréquence propre, la \cle{résonance} produit des amplitudes énormes pouvant conduire à la rupture de l'ouvrage (exemple historique : pont de Tacoma, 1940). Les ingénieurs doivent donc écarter les fréquences propres des fréquences d'excitation probables et ajouter des dispositifs amortisseurs.
\end{enumerate}
\textbf{Points de vigilance :} distinguer clairement les trois régimes (libre non amorti, amorti, forcé) et ne pas confondre excitateur et résonateur ; pour la partie A, justifier que l'équation différentielle s'écrit \emph{par rapport à la position d'équilibre} (le poids est compensé par la tension d'équilibre) ; à la résonance, c'est l'\emph{amplitude} qui est maximale, pas la fréquence.
\end{corrige}

\newpage

\coursec{Fiche bilan}

\begin{popfigure}
\begin{tikzpicture}[mindmap, grow cyclic, every node/.style={concept, text=white, font=\small\bfseries},
  root concept/.append style={concept color=popdark, font=\large\bfseries},
  level 1/.append style={level distance=3.6cm, sibling angle=60},
  level 1 concept/.append style={font=\footnotesize\bfseries},
  scale=0.92, transform shape]
\node[root concept]{Oscillateurs mécaniques}
  child[concept color=popblue]{ node[align=center] {Pendule élastique horizontal\\ $T_0=2\pi\sqrt{m/k}$} }
  child[concept color=popgreen]{ node[align=center] {Pendule élastique vertical\\ m\^eme $T_0$ (r\'ef. d'\'equilibre)} }
  child[concept color=poporange]{ node[align=center] {Pendule simple\\ $T_0=2\pi\sqrt{\ell/g}$\\ (petits angles)} }
  child[concept color=poppurple]{ node[align=center] {Pendule de torsion\\ $T_0=2\pi\sqrt{J/C}$} }
  child[concept color=poppink]{ node[align=center] {\'Energie\\ $E=\tfrac12 kX_m^2$\\ $=\tfrac12 m\omega_0^2X_m^2$} }
  child[concept color=popgold]{ node[align=center] {Amortissement\\ pseudo-p\'eriodique\\ ap\'eriodique} }
  child[concept color=popmuted]{ node[align=center] {Oscillations forc\'ees\\ r\'esonance $\omega\simeq\omega_0$} };
\end{tikzpicture}
\legende{Carte mentale de la le\c{c}on : les sept id\'ees-forces des oscillateurs m\'ecaniques.}
\end{popfigure}

\begin{bilan}
\textbf{D\'efinitions cl\'es}
\begin{itemize}
  \item \cle{Oscillateur m\'ecanique} : syst\`eme qui, \'ecart\'e de sa position d'\'equilibre stable, effectue un mouvement p\'eriodique (ou pseudo-p\'eriodique) autour de celle-ci sous l'effet d'une \cle{force de rappel}.
  \item \cle{Amplitude} $X_m$ : valeur maximale de l'\'elongation (en m). \cle{Pulsation propre} $\omega_0$ (rad/s). \cle{Phase} $\omega_0 t+\varphi$ (rad) ; $\varphi$ : phase \`a l'origine.
  \item \cle{P\'eriode propre} $T_0=\dfrac{2\pi}{\omega_0}$ (s) ; \cle{fr\'equence} $f_0=\dfrac{1}{T_0}$ (Hz).
  \item \cle{R\'esonance} : pour un oscillateur forc\'e, l'amplitude devient maximale quand la pulsation de l'excitateur est voisine de $\omega_0$ (d'autant plus aigu\"e que l'amortissement est faible).
\end{itemize}

\textbf{Formules \`a conna\^itre par c\oe ur}
\begin{center}
\begin{tabularx}{\linewidth}{|l|X|X|}\hline
\rowcolor{popblueL} \textbf{Oscillateur} & \textbf{\'Equation / condition} & \textbf{P\'eriode propre} \\ \hline
Pendule \'elastique horizontal & $\ddot{x}+\dfrac{k}{m}x=0$ & $T_0=2\pi\sqrt{\dfrac{m}{k}}$ \\ \hline
Pendule \'elastique vertical & $\ddot{x}+\dfrac{k}{m}x=0$ (autour de l'\'equilibre) & $T_0=2\pi\sqrt{\dfrac{m}{k}}$ \\ \hline
Pendule simple ($\theta\lesssim 10^{\circ}$) & $\ddot{\theta}+\dfrac{g}{\ell}\theta=0$ & $T_0=2\pi\sqrt{\dfrac{\ell}{g}}$ \\ \hline
Pendule de torsion & $\ddot{\theta}+\dfrac{C}{J}\theta=0$ & $T_0=2\pi\sqrt{\dfrac{J}{C}}$ \\ \hline
\end{tabularx}
\end{center}

Solution g\'en\'erale : $x(t)=X_m\cos(\omega_0 t+\varphi)$, avec $X_m$ et $\varphi$ fix\'es par les \cle{conditions initiales} $x(0)$ et $\dot{x}(0)$.

\'Energie m\'ecanique (non amorti) : $E=E_c+E_p=\dfrac12 kX_m^2=\dfrac12 m\omega_0^2X_m^2$ = constante ; $E_c$ maximale \`a l'\'equilibre, $E_p$ maximale aux extr\'emit\'es.

\textbf{M\'ethodes express}
\begin{enumerate}
  \item \textbf{Mettre en \'equation} : rep\`ere d'origine \`a l'\'equilibre $\to$ bilan des forces $\to$ 2\up{e} loi de Newton projet\'ee $\to$ identifier $\omega_0^2$ $\to$ en d\'eduire $T_0$.
  \item \textbf{Conditions initiales} : $x(0)=X_m\cos\varphi$ et $\dot{x}(0)=-X_m\omega_0\sin\varphi$ ; si l\^ach\'e sans vitesse depuis $x_0$, alors $\varphi=0$ et $X_m=x_0$.
  \item \textbf{Exploiter un enregistrement} : mesurer $T_0$ sur plusieurs p\'eriodes, lire $X_m$, en d\'eduire $\omega_0$, puis $k$, $g$ ou $C$ selon l'oscillateur.
  \item \textbf{R\'egimes amortis} : frottement faible $\to$ pseudo-p\'eriodique ($T\simeq T_0$) ; frottement fort $\to$ ap\'eriodique (retour \`a l'\'equilibre sans oscillation).
\end{enumerate}
\end{bilan}

\begin{aretenir}[Checklist avant le Bac]
\begin{itemize}
  \item[$\square$] Je sais \'etablir $\ddot{x}+\omega_0^2x=0$ par la 2\up{e} loi de Newton, en choisissant l'origine \`a la position d'\'equilibre.
  \item[$\square$] Je sais que le pendule \'elastique vertical a la m\^eme p\'eriode que l'horizontal : le poids ne fait que d\'eplacer l'\'equilibre ($k\,\Delta\ell_0=mg$).
  \item[$\square$] Je connais par c\oe ur $T_0=2\pi\sqrt{m/k}$, $T_0=2\pi\sqrt{\ell/g}$, $T_0=2\pi\sqrt{J/C}$ et je v\'erifie l'homog\'en\'eit\'e.
  \item[$\square$] Je sais d\'eterminer $X_m$ et $\varphi$ \`a partir de $x(0)$ et $\dot{x}(0)$.
  \item[$\square$] Je sais que la condition $\theta\lesssim 10^{\circ}$ (petites oscillations, $\sin\theta\simeq\theta$ en rad) est indispensable pour le pendule simple.
  \item[$\square$] Je sais calculer $E=\tfrac12 kX_m^2$ et retrouver la vitesse maximale $v_{\max}=\omega_0 X_m$.
  \item[$\square$] Je sais mesurer une p\'eriode sur un enregistrement (sur 5 \`a 10 p\'eriodes pour la pr\'ecision) et lire une amplitude.
  \item[$\square$] Je distingue r\'egime pseudo-p\'eriodique et ap\'eriodique, et je sais que $T\simeq T_0$ si l'amortissement est faible.
  \item[$\square$] Je sais d\'efinir la r\'esonance et citer un exemple (pont, balan\c{c}oire, suspension de mototaxi).
  \item[$\square$] Je soigne les unit\'es : $m$ en kg, $k$ en N/m, $\ell$ en m, $J$ en kg\,m\up{2}, $C$ en N\,m/rad, angles en radians.
\end{itemize}
\end{aretenir}

\findecours

\end{document}
